%=============================================================================
% UIDT Ontology v3.9 — Vacuum Information Density as the Fundamental Geometric Scalar
% Author: Philipp Rietz (ORCID: 0009-0007-4307-1609)
% DOI (this manuscript, Ontology v3.9): 10.5281/zenodo.20319634
% DOI (UIDT Framework v3.9 parent): 10.5281/zenodo.17835200
% License: CC BY 4.0
%
% PHASE 3 SYNTHESIS NOTE (non-rendering):
% This manuscript retains the visual design language of v3.9.8.4 FINAL but
% adopts the corrected, audit-reconciled evidence discipline of the Phase 1/2
% MASTER_OUTLINE. It is an auditable research ontology, NOT a merge
% authorization, NOT a canonical freeze, NOT a proof of UIDT, and NOT an
% evidence-class upgrade. AI-authored manuscript text is advisory only
% (AUDIT_POLICY Sec. 1).
%=============================================================================
\PassOptionsToPackage{hyphens}{url}
\PassOptionsToPackage{table,svgnames,dvipsnames}{xcolor}

\documentclass[12pt,a4paper,twoside]{article}

%=============================================================================
% PACKAGES
%=============================================================================
\usepackage[utf8]{inputenc}
\usepackage[T1]{fontenc}
\usepackage{microtype}
\usepackage{fix-cm}
\usepackage{mathpazo}

\usepackage{geometry}
\usepackage{setspace}
\usepackage{titlesec}
\usepackage{fancyhdr}
\usepackage{booktabs}
\usepackage{longtable}
\usepackage{array}
\usepackage{multirow}
\usepackage{enumitem}
\usepackage{float}
\usepackage{caption}
\usepackage[numbers,sort&compress]{natbib}
\usepackage{etoolbox}
\usepackage{needspace}
\usepackage{placeins}

\usepackage{amsmath,amssymb,amsfonts,amsthm}
\usepackage{mathtools}
\usepackage{bm}
\usepackage{physics}

\usepackage{graphicx}
\usepackage{xcolor}

\usepackage{tikz}
\usetikzlibrary{shapes,arrows,positioning,shadows}
\usepackage[most]{tcolorbox}
\tcbuselibrary{skins,breakable}

\usepackage{hyperref}
\usepackage{cleveref}

%=============================================================================
% COLORS
%=============================================================================
% MONOCHROME PALETTE (master-class serious; no decorative colour)
% All evidence-class tags render in black; distinction is typographic, not chromatic.
% Greys are used only for box rules and structural hairlines.
\definecolor{catA}{RGB}{0,0,0}
\definecolor{catAminus}{RGB}{0,0,0}
\definecolor{catB}{RGB}{0,0,0}
\definecolor{catC}{RGB}{0,0,0}
\definecolor{catD}{RGB}{0,0,0}
\definecolor{catE}{RGB}{0,0,0}
\definecolor{darkblue}{RGB}{0,0,0}
\definecolor{dimgray}{RGB}{105,105,105}
\definecolor{pioverride}{RGB}{0,0,0}
\definecolor{uidtBlack}{RGB}{0,0,0}
\definecolor{uidtRule}{gray}{0.22}
\definecolor{uidtHairline}{gray}{0.62}
\definecolor{uidtSoft}{gray}{0.985}
\definecolor{uidtTitleShade}{gray}{0.955}
\definecolor{uidtTitleRule}{gray}{0.12}

% Ensure darkgray is defined regardless of xcolor option loading order
\providecolor{darkgray}{RGB}{64,64,64}

%=============================================================================
% PAGE SETUP
%=============================================================================
\geometry{a4paper, left=3cm, right=3cm, top=3cm, bottom=3cm}
\setlength{\headheight}{14.5pt}
\addtolength{\topmargin}{-2.5pt}
\clubpenalty=10000
\widowpenalty=10000
\raggedbottom
\setlength{\parskip}{0.65em plus 0.12em minus 0.04em}
\setlength{\parindent}{0pt}
\setstretch{1.32}

%=============================================================================
% SECTION HEADING STYLE
%=============================================================================
\titleformat{\section}
 {\normalfont\fontsize{16}{19}\selectfont\bfseries} {\thesection}{1em}{}
\titlespacing*{\section}{0pt}{3.0ex plus 1ex minus 0.2ex}{1.8ex plus 0.2ex}

\titleformat{\part}[display]
 {\normalfont\huge\bfseries\filcenter} {\partname\ \thepart}{0.6em}{\fontsize{20}{24}\selectfont}
\titlespacing*{\part}{0pt}{0pt}{2.5ex}

%=============================================================================
% HEADER/FOOTER
%=============================================================================
\fancypagestyle{uidtontology}{%
 \fancyhf{}
 \fancyhead[LE,RO]{\thepage}
 \fancyhead[RE]{\small\textit{UIDT v3.9: Ontological Foundations}}
 \fancyhead[LO]{\small\textit{P.\ Rietz}}
 \fancyfoot[C]{%
 \footnotesize\textcolor{darkgray}{%
 \copyright\ 2026 P.\ Rietz $\cdot$ CC BY 4.0 $\cdot$
 \href{https://doi.org/10.5281/zenodo.20319634}{DOI: 10.5281/zenodo.20319634} | \href{https://github.com/Mass-Gap/UIDT-Framework-v3.9-Canonical}{UIDTv3.9 Repository}}%
 }
 \renewcommand{\headrulewidth}{0.4pt}
 \renewcommand{\footrulewidth}{0.4pt}
}

%=============================================================================
% METADATA
%=============================================================================
\hypersetup{
 pdftitle={Vacuum Information Density as the Fundamental Geometric Scalar (UIDT v3.9)}, pdfauthor={Philipp Rietz}, pdfsubject={Structural-realist ontology of information-geometric physics}, pdfkeywords={Ontology, Information Density, Yang-Mills, Mass Gap, Vacuum Geometry, UIDT}, colorlinks=true, linkcolor=black, citecolor=black, urlcolor=black, pdflang={en}
}

%=============================================================================
% THEOREM ENVIRONMENTS
%=============================================================================
\BeforeBeginEnvironment{theorem}{\par\nopagebreak\needspace{4\baselineskip}}
\BeforeBeginEnvironment{definition}{\par\nopagebreak\needspace{4\baselineskip}}

\theoremstyle{plain}
\newtheorem{theorem}{Theorem}[section]
\newtheorem{proposition}[theorem]{Proposition}
\newtheorem{lemma}[theorem]{Lemma}
\newtheorem{corollary}[theorem]{Corollary}

\theoremstyle{definition}
\newtheorem{definition}[theorem]{Definition}
\newtheorem{axiom}{Axiom}
\newtheorem{prediction}{Testable Prediction}
\newtheorem{falsmark}{Falsification Criterion}

\theoremstyle{remark}
\newtheorem{remark}[theorem]{Remark}
\newtheorem{limitation}[theorem]{Limitation}
\newtheorem{openquestion}{Open Question}

% cleveref names for theorem-like custom environments (surgical v3.9.9-final fix)
\crefname{axiom}{axiom}{axioms}
\Crefname{axiom}{Axiom}{Axioms}
\crefname{prediction}{prediction}{predictions}
\Crefname{prediction}{Prediction}{Predictions}
\crefname{openquestion}{open question}{open questions}
\Crefname{openquestion}{Open Question}{Open Questions}
\crefname{limitation}{limitation}{limitations}
\Crefname{limitation}{Limitation}{Limitations}

% Custom Commands
\newcommand{\UIDT}{\textsc{uidt}}
\newcommand{\GeV}{\,\mathrm{GeV}}
\newcommand{\MeV}{\,\mathrm{MeV}}
\newcommand{\catmark}[1]{\textsuperscript{\textsf{\textbf{[#1]}}}}
\newcommand{\Sx}{S(x)}
\newcommand{\lambdaS}{5/12 \approx 0.4167} % [A] display alias for exact RG value 5\kappa^2/3.

%===================================================================
% EVIDENCE-DISCIPLINE MACROS (Phase 3, reconciled with MASTER_OUTLINE)
% These macros render evidence tags. They invent NO new evidence class.
% \Bdelta renders a clean [B] and (silently) hyperlinks back to the single
% authoritative PI-override box in Chapter 1. No visible asterisk is used,
% per the decision that the most rigorous form introduces no new glyph
% (cf. AI_AUDIT_POLICY Sec. 5; Epistemological Cleanup Handout Sec. IV).
%===================================================================
\newcommand{\Bdelta}{\hyperref[box:pi-override-delta]{\catmark{B}}}
\newcommand{\Bpending}{\catmark{B\,pending}}
%===================================================================
% UIDT PARAMETER LEDGER MACROS (v3.9, reconciled evidence classes)
% RACE CONDITION LOCK: do NOT move mp.dps = 80 to config.
% These macros are LaTeX display aliases only; they carry no numerical
% computation. The authoritative values live in LEDGER/CONSTANTS.md.
%===================================================================
\newcommand{\DeltaGap}{1.710 \pm 0.015\GeV} % [B] PI-override (was [A-]); see box:pi-override-delta
\newcommand{\kappaVal}{0.500} % [A] dimensionless EFT-normalised coefficient
\newcommand{\GammaGeom}{16.339} % [A-] calibrated, NEVER derived
\newcommand{\GammaMC}{16.374 \pm 1.005} % [A-] Monte Carlo mean; tension entry vs 16.339
\newcommand{\vVEV}{47.7\MeV} % [A]
\newcommand{\ETorsion}{2.44\MeV} % [E pending] observer-context only; NOT ego/life/death threshold
\newcommand{\LambdaGrid}{0.66\,\mathrm{nm}} % [D]
\newcommand{\KappaRG}{\tfrac{1}{2}\ \text{(exact)}} % [A] exact framework input; fixed-point reading untested by any defined flow; not an empirical pm
%===================================================================

%=============================================================================
% TCOLORBOX ENVIRONMENTS
%=============================================================================
\tcbset{%
 uidtboxbase/.style={%
  enhanced, breakable, sharp corners, colback=uidtSoft, colframe=uidtRule, coltitle=uidtBlack, boxrule=0pt, arc=0pt, left=11pt, right=11pt, top=9pt, bottom=9pt, boxsep=0pt, before skip=1.15em plus 0.25em minus 0.15em, after skip=1.15em plus 0.25em minus 0.15em, fontupper=\small\setstretch{1.18}, parbox=false, borderline north={0.45pt}{0pt}{uidtHairline}, borderline south={0.45pt}{0pt}{uidtHairline}, },
 uidttitle/.style={%
  colbacktitle=white,
  boxed title style={%
   sharp corners, colback=white, colframe=white, boxrule=0pt, left=0pt, right=0pt, top=0pt, bottom=0pt, }, attach boxed title to top left={xshift=11pt,yshift=-1.8mm}, },
}

\newtcolorbox{evidencebox}[2][]{%
 uidtboxbase, uidttitle, title={#1},
 fonttitle=\scriptsize\bfseries\sffamily\MakeUppercase\lsstyle,
 top=17pt, borderline west={1.1pt}{0pt}{uidtTitleRule},
 overlay unbroken and first={%
 \node[anchor=north east, font=\tiny\sffamily\bfseries\lsstyle, text=uidtHairline] at ([xshift=-7pt,yshift=14pt]frame.north east) {EVIDENCE}; },
}

\newtcolorbox{killswitch}{%
 uidtboxbase, colframe=uidtBlack, colback=white,
 fonttitle=\scriptsize\bfseries\sffamily\MakeUppercase\lsstyle,
 title={Falsification Criterion}, boxrule=0pt, borderline north={0.65pt}{0pt}{uidtBlack}, borderline south={0.65pt}{0pt}{uidtBlack}, borderline west={1.8pt}{0pt}{uidtBlack},
}

\newtcolorbox{axiombox}[1][]{%
 uidtboxbase, uidttitle, colframe=uidtRule,
 fonttitle=\scriptsize\bfseries\sffamily\MakeUppercase\lsstyle,
 title={#1}, borderline west={0.9pt}{0pt}{uidtRule},
}

\newtcolorbox{probbox}[2][]{%
 uidtboxbase, uidttitle, colframe=uidtBlack, colback=white,
 fonttitle=\scriptsize\bfseries\sffamily\MakeUppercase\lsstyle,
 title={#1 \hfill #2}, boxrule=0pt, borderline north={0.55pt}{0pt}{uidtBlack}, borderline south={0.55pt}{0pt}{uidtBlack}, borderline west={1.4pt}{0pt}{uidtBlack},
}

% DETERMINATION BOX: the single authoritative home for the Delta evidence-class
% reclassification. A heavier solid frame marks it as a governance artifact, not
% a physics result, distinguished by weight rather than colour.
\newtcolorbox{pioverridebox}[1][]{%
 uidtboxbase, uidttitle, colframe=uidtBlack, colback=white,
 fonttitle=\scriptsize\bfseries\sffamily\MakeUppercase\lsstyle,
 title={#1}, boxrule=0pt, borderline north={0.75pt}{0pt}{uidtBlack}, borderline south={0.75pt}{0pt}{uidtBlack}, borderline west={2.2pt}{0pt}{uidtBlack},
}

% GOVERNANCE / NO-MERGE banner box
\newtcolorbox{governancebox}[1][]{%
 uidtboxbase, uidttitle, colframe=uidtHairline,
 fonttitle=\scriptsize\bfseries\sffamily\MakeUppercase\lsstyle,
 title={#1}, boxrule=0pt, borderline west={1.1pt}{0pt}{uidtRule},
}

%=============================================================================
\begin{document}
\pagestyle{uidtontology}

%=============================================================================
% TITLE
%=============================================================================
\begin{titlepage}
\vspace*{1.6cm}
\begin{center}
{\LARGE\bfseries Vacuum Information Density\\[0.25em]
as the Fundamental Geometric Scalar}\\[0.9em]
{\small\bfseries A Structural-Realist Account of \(S(x)\) as an Ontological Primitive,\\[0.25em]
the Yang--Mills Spectral Gap as a Stress Test,\\[0.25em]
Emergent Spacetime, and the Open \(G_{\mathrm{SM}}\)-Origin Problem}\\[1.5em]
{\large\itshape Unified Information-Density Theory (UIDT) v3.9\\[0.25em]
Ontological Foundations, Audit-Reconciled Edition}\\[2.6em]
{\small Philipp Rietz}\\[0.5em]
{\small ORCID: \href{https://orcid.org/0009-0007-4307-1609}{0009-0007-4307-1609}}\\[0.2em]
{\small Independent Researcher}\\[3em]
{\normalsize UIDT Framework v3.9}\\[0.5em]
{\small This manuscript (Ontology v3.9):
\href{https://doi.org/10.5281/zenodo.20319634}{DOI 10.5281/zenodo.20319634}}\\[0.25em]
{\small UIDT Framework (parent record):
\href{https://doi.org/10.5281/zenodo.17835200}{DOI 10.5281/zenodo.17835200}}\\[0.5em]
{\small 31~May~2026 }\\[2.2em]
\rule{0.6\textwidth}{0.4pt}\\[1em]
{\small\itshape ``Reality is not a collection of things,\\
but a process of structural individuation\\
that we are licensed to describe only as far as we can falsify it.''}\\[1.4em]
\end{center}
\end{titlepage}

\newpage
\tableofcontents
\newpage

%=============================================================================
% ABSTRACT (audit-reconciled)
%=============================================================================
\begin{abstract}
\begingroup \normalfont\normalsize
\setlength{\parskip}{0.65em plus 0.12em minus 0.04em}%
\setlength{\parindent}{0pt}%
\noindent This document formulates the ontological and epistemological foundations of the Unified Information-Density Theory (\UIDT). Rather than postulating spacetime as a primordial arena, \UIDT\ investigates the hypothesis that a single real scalar field, the \emph{vacuum information density} \mbox{$\Sx$}, may serve as a fundamental geometric degree of freedom. The present edition is deliberately conservative: it is the product of a line-by-line honesty audit that downgraded several earlier claims and that treats the framework as a falsification-first research programme rather than a completed theory.

The framework rests upon four explicit axioms and addresses four areas: (I)~a quantum field-theoretic foundation in which the pure-$\mathrm{SU}(3)$ Yang--Mills spectral gap \mbox{$\Delta = \DeltaGap$\,\Bdelta} functions as a necessary \emph{stress test} of internal consistency rather than as a sufficient foundation, and in which the gap value is lattice-compatible under an explicitly recorded evidence-class reclassification (\cref{box:pi-override-delta}); (II)~cosmological calibration to Planck~2018, SH0ES, DESI~DR2, KiDS-Legacy, and DES~Year~6 data, capped at \catmark{C}\ and never presented as a resolution of the $H_0$ or $S_8$ tensions; (III)~laboratory-relevant predictions, including a candidate scalar resonance at \mbox{$m_S = 1.705 \pm 0.015\GeV$\,\catmark{D}} and a Casimir-scale anomaly at \mbox{$\lambda = \LambdaGrid$\,\catmark{D}}; and (IV)~a structural mapping between the vacuum scalar and effective refractive-index profiles, retained as an interpretive analogy.

The calibrated invariant \mbox{$\gamma = \GammaGeom$} is classified throughout as \catmark{A-}\ (calibrated, \emph{never} derived from renormalisation-group first principles); any future derivation must satisfy an explicit anti-target-leakage criterion. The naive scalar-gradient route to non-trivial gauge curvature, $A = \dd S \Rightarrow F = \dd A = \dd^2 S = 0$, is recorded as a clean \catmark{A}\ negative result. The central open problem, the \emph{$G_{\mathrm{SM}}$-Origin Gap}, is preserved as unsolved.
\newpage
Every quantitative claim carries an explicit evidence classification (A through E), every prediction is accompanied by a falsification criterion, every known limitation (L1--L12) is disclosed, and speculative observer-phenomenology material is retained strictly as bounded Category~\catmark{E}\ research-programme content within an interpretive stratum. The ontological position is a form of dynamic Ontic Structural Realism in which sufficiently constrained mathematical structure is taken to constitute, rather than merely describe, physical reality, subject always to the demarcation that interpretation never substitutes for empirical or formal closure.

\endgroup
\end{abstract}

\newpage

%=============================================================================
% PART I
%=============================================================================
\part{Evidentiary Ground, Source Integrity, and Empirical Boundary Conditions}
\label{part:evidentiary-ground}

\noindent The order of this manuscript is deliberate. Before a single ontological claim is advanced, before the scalar field $\Sx$ is endowed with any generative role, this first part establishes the discipline by which every later assertion will be judged. A theory of the foundations of physics earns the right to speculate only in proportion to the severity of the constraints it accepts in advance. Part~I therefore fixes three things: the corpus from which the manuscript is drawn and the integrity conditions on that corpus; the grammar of evidence classes and the separation of epistemic strata; and the governance boundary that prevents confident prose from ever functioning as a substitute for proof or for human review. Only then does Part~I turn to the empirical physics that the framework must respect, including pure $\mathrm{SU}(3)$ Yang--Mills theory, its topological observables, and the calibrated invariant $\gamma$, and it treats each of these as a constraint on \UIDT, not yet as a success of it.

%=============================================================================
% CHAPTER 1
%=============================================================================
\newpage
\section{Sources, Evidence, and Audit Discipline}
\label{sec:corpus-ledger-audit}

\subsection{The Source Base and Its Traceability}
\label{subsec:corpus-topology}

The material from which this manuscript is synthesised is not a single document but a layered, version-controlled body of source material. It is useful to make its structure explicit, because the credibility of an ontological claim depends in part on the traceability of the path from raw source to finished assertion. The sources fall into several kinds: a primary developmental record; additional corroborating material developed in parallel; independently sourced external documents; and the formal manuscripts and verification records produced along the way.\catmark{A}

The primary developmental record is treated as the main line of development, but it is explicitly \emph{not} the only source of record; the corroborating material and the independently sourced external documents carry independent weight, and a claim that rests on the developmental record alone is treated as weaker than one corroborated across several independent sources.\catmark{A}

A central methodological principle governs how new material is incorporated: nothing is overwritten, summarised away, or silently corrected. Superseded formulations are retained alongside the formulations that replaced them, rather than deleted.\catmark{A} This discipline has a direct epistemic payoff: because nothing is quietly removed, the history of superseded ideas remains visible and checkable, and the maturity of the framework can be measured by the visibility of its abandoned routes rather than concealed by their absence.

Two constraints on the use of external literature follow, and are stated here as standing requirements, because their violation has been a recurrent failure mode.
\newpage
\begin{governancebox}[Source-Handling Requirements]
\textbf{Native extraction precedes mapping.} Any external paper invoked in support of a \UIDT\ mapping must first be extracted natively, including its text, identifiers, and line structure recovered, before any interpretive use. A citation that cannot be resolved to an extracted source with a verifiable identifier is treated as contextual \catmark{E}\ material only, never as
evidence.\\[0.4em]
\textbf{Assistant acknowledgements are not evidence.} An affirmative or agreeable response generated by a language model in the course of the developmental dialogue does not constitute deep-research evidence. The presence of a confident reply in the corpus establishes only that the reply was made, not that its content is correct.
\end{governancebox}

\noindent These two requirements are unglamorous, but they are load-bearing. A significant portion of the honesty audit that produced this edition consisted of locating places where an assistant's fluent agreement had been mistaken for an established result, and of demoting such places to their proper status. The requirement that every final claim trace to a source family and an evidence tag is therefore not bureaucratic decoration; it is the mechanism by which the corpus defends itself against its own fluency.\catmark{A}

\subsection{Evidence-Class Grammar and Stratum Separation}
\label{subsec:evidence-grammar}

\UIDT\ assigns to every quantitative or structural claim a label drawn from a fixed grammar of evidence classes. The grammar is small by design, and its discipline lies in refusing to expand it. The classes are as follows.\catmark{A}

\begin{itemize}[leftmargin=2.2em]
 \item \textbf{\catmark{A}, Analytically/logically closed.} The claim is a formal, logical, or mathematical closure within a specified system. Class~A denotes \emph{internal} demonstrability, not rhetorical confidence and not external endorsement. For a numerical residual, the internal deterministic closure criterion is a residual below $10^{-14}$; this criterion certifies a stable algorithm, not agreement with nature. \item \textbf{\catmark{A-}, Calibrated.} The quantity is fixed by calibration to a chosen reference, not derived from first principles. A calibrated value may be precise and may be useful, but precision is not provenance. The invariant $\gamma$ lives here permanently.
 \newpage
 \item \textbf{\catmark{B}, Lattice/consensus compatible.} The claim is compatible with an externally established result, typically lattice QCD or Standard-Model consensus, but only after a full audit of the external source, its uncertainty budget, the matched observable, the units, and the covariance structure. Compatibility is not derivation. \item \textbf{\catmark{C}, Calibrated to cosmological data.} The claim is fitted to observational cosmology. This is the \emph{maximum} class any cosmological quantity may attain. No cosmological parameter is ever promoted above~C, and no fit to cosmological data is ever described as a resolution of a cosmological tension. \item \textbf{\catmark{D}, Predicted / research programme.} The claim is a prediction awaiting experimental test, or a structured programme of future work. It is held to be unverified by definition. \item \textbf{\catmark{E}, Speculative / interpretive / phenomenological.} The claim is interpretive, exploratory, or drawn from boundary-case phenomenology. Category~E material is retained for its heuristic and demarcational value and is never advanced as support for a physical claim.
\end{itemize}

The grammar admits no other symbols. Invented classes, including constructions such as $[\mathrm{A+}]$, $[\mathrm{B+}]$, $[\mathrm{B-}]$, $[\mathrm{C+}]$, or $[\mathrm{D+}]$ do not exist and are treated, where they appear in superseded material, as artifacts to be removed.\catmark{A} The reason for this severity is practical: a graduated tier invites the quiet inflation of a calibrated result into a near-derivation by the addition of a plus sign, and the history recorded in this corpus contains exactly that failure.

Orthogonal to the evidence classes, and equally binding, is the separation of the claim into three strata. The strata answer a different question from the classes: not \emph{how well supported} a statement is, but \emph{what kind of statement} it is.\catmark{A}

\begin{itemize}[leftmargin=2.2em]
 \item \textbf{Stratum~I: Empirical.} Direct external measurements with official error budgets. These are immutable inputs. Every Stratum~I value must carry an uncertainty; a measured quantity stated without $(\pm)$ is, by this rule, malformed.
 \\
 \item \textbf{Stratum~II: Consensus.} The current state of the field, the consensus of lattice QCD, of the Standard Model, of observational cosmology, including its own honest description of open tensions. Stratum~II requires proper external citation. \item \textbf{Stratum~III: UIDT interpretation.} The framework's own calibrated values, mappings, and interpretive structures. Stratum~III content may never be used as empirical ground truth; it may never serve as the reference against which a $z$-score or a $\sigma$-level agreement is computed, because to do so would be to test the theory against itself.
\end{itemize}

The most consequential consequence of this separation is the prohibition on circular validation: a UIDT-internal value (Stratum~III) cannot be promoted to play the role of an empirical anchor (Stratum~I). The honesty audit found this boundary crossed in several places, and the correction of those crossings is among the principal reasons this edition exists.

It is at this point, within the grammar of evidence, that the manuscript records its single most important class reclassification. The spectral gap value $\Delta = \DeltaGap$ is, within the \UIDT\ axiom system, established by an internal fixed-point construction whose residual meets the Class~A internal criterion. Yet the framework's own public documentation states explicitly that internal mathematical consistency is not external peer review, which is, word for word, the definition of Class~B rather than Class~A. To leave the gap labelled~A while the project's own note disavows the external content of that label would be to smooth over an inconsistency. The honest and inexpensive correction is to align the ledger to its own note.

\begin{pioverridebox}[Evidence-Class Determination: $\Delta$ Spectral Gap is \catmark{B}]
\label{box:pi-override-delta}
\textbf{Subject.} The pure-$\mathrm{SU}(3)$ spectral-gap value
$\Delta = 1.710 \pm 0.015\GeV$.\\[0.4em]
\textbf{Algorithmic class.} \catmark{A}\, the value satisfies the internal deterministic closure criterion of the \UIDT\ fixed-point construction (residual below the internal threshold). This certifies internal mathematical
consistency only.\\[0.4em]
\textbf{Determination.} The ledger class is set to \catmark{B}\ (lattice/consensus compatible). \emph{Rationale, verbatim from the framework's own Pillar~I note:} ``Category~A here denotes internal mathematical consistency. It does not imply external peer review.'' That statement is exactly the definition of Class~B. Aligning the ledger to its own note is the honest and
inexpensive correction; it neither inflates nor conceals.\\[0.4em]
\textbf{Mechanism.} The reclassification follows the same documented override procedure used elsewhere in this work: a value that is algorithmically~A is recorded as~B where its external content has not been peer-confirmed, with the rationale stated. The determination remains pending an external counter-signature in the sense of the framework's evidence discipline; until such a counter-signature exists, $\Delta$~is reported as~B and never as a derived or peer-confirmed
result.\\[0.4em]
\textbf{Rendering convention.} Throughout this manuscript, $\Delta$ carries a clean \catmark{B}\ tag that hyperlinks back to this box. No new evidence glyph is introduced; the override lives here, once, and is reflected as a formal entry in the Claims Table (\cref{sec:final-claims}) and the Tension Ledger (\cref{sec:tension-ledger}).
\end{pioverridebox}

\noindent This determination is the template for how \UIDT\ treats the gap between internal consistency and external confirmation everywhere: not by hiding the gap, and not by pretending it has been closed, but by labelling it precisely and recording where the boundary lies.

\subsection{Internal Consistency Is Not External Confirmation}
\label{subsec:governance-boundary}

A single methodological principle underwrites the entire evidence grammar, and it is worth stating plainly because its violation is the most common way a framework deceives itself. \emph{An internal consistency check is not an external confirmation.} That a value satisfies the deterministic closure criterion of the \UIDT\ construction establishes internal mathematical consistency only; it does not, by itself, license raising that value's evidence class to one that asserts external agreement. No internal calculation, and no automated analysis of the manuscript, can on its own promote a claim from one evidence class to a higher one. Such a promotion requires the independent, and where appropriate externally counter-signed, confirmation specified by the framework's stated evidence discipline.\catmark{A}

The motivation for stating this so emphatically is concrete. The developmental record contains an episode in which a numerologically constructed value was treated, on the strength of a confident internal check, as an established result, even though its own forward predictions failed by large factors and its prime-factor structure showed it had been built to hit a target rather than discovered. The lesson is not that internal checks are useless, but that they must never be load-bearing for a claim of external validity. A check that can raise an objection but cannot by itself confer confirmation cannot reproduce that failure.\catmark{A}

A direct corollary governs the language of finality. Terms such as ``definitive resolution,'' ``final proof,'' or ``settled'' are not permitted to describe the framework's status unless the exact formal conditions for a Class~A closure are demonstrably met. Where the developmental dialogue reached for such language, this edition replaces it with the accurate phrase: a \emph{current boundary pending external confirmation}. The distinction is not cosmetic: ``final'' asserts that no further work is owed, whereas ``current boundary'' asserts only that this is the present edge of what has been established, and that the edge is expected to move.\catmark{A}

The same discipline designates a class of \emph{museum} or \emph{quarantine} material: formulas, values, and derivational routes that have been tried and refuted. These are retained, visibly, as anti-patterns rather than hidden as embarrassments, and the rule governing them is absolute: a quarantined construction may be cited only as documentation of a failed approach, never as evidence for a live claim. The most important quarantined items are recorded explicitly later in this work. Taken together with the corpus discipline of \cref{subsec:corpus-topology} and the evidence grammar of this section, these principles are the standing conditions under which the rest of the manuscript operates. They are restrictive on purpose: a framework that proposes to identify a single scalar field as the root of physical structure must accept, in exchange for that ambition, an unusually severe accounting of what it has and has not established.\catmark{A}

\begin{governancebox}[Mathematical Scope]
This ontology does not provide a Wightman construction, an
Osterwalder--Schrader reconstruction, a nonperturbative construction of
four-dimensional Yang--Mills theory, a proof of the Clay Millennium problem,
a derivation of the Standard-Model gauge group $G_{\mathrm{SM}}$, or a
Born-rule theorem. It defines a research ontology together with a set of
falsification and audit constraints; nothing in it upgrades any evidence
class.
\end{governancebox}

%=============================================================================
% CHAPTER 2
%=============================================================================
\newpage
\section{The Empirical Physics Baseline Before UIDT Interpretation}
\label{sec:empirical-baseline}

\subsection{Pure \texorpdfstring{$\mathrm{SU}(3)$}{SU(3)} as a Stress-Test Sector}
\label{subsec:su3-stress-test}

The most disciplined way to test a candidate ontology of gauge structure is to confront it first with the cleanest non-trivial gauge theory available, in isolation from the complications it does not yet claim to address. For \UIDT, that arena is pure $\mathrm{SU}(3)$ Yang--Mills theory: the gluon sector without dynamical quarks. It is the natural stress-test sector precisely because it is both rich enough to possess a mass gap and simple enough to be computed to high precision on the lattice by an external community whose results owe nothing to \UIDT.\catmark{A}

The framing here is a deliberate inversion of an earlier temptation, and it is the single most important reframing in this chapter. The Yang--Mills spectral gap is to be understood as a \emph{necessary} condition that any serious candidate must pass, not as a \emph{sufficient} foundation from which the rest of physics follows. Reproducing, or being compatible with, the pure-gauge mass gap demonstrates that the framework is not immediately excluded; it does not demonstrate that the framework is correct, and it most emphatically does not derive the Standard Model.\catmark{A} A stress test is something a theory survives, not something it stands upon.

This requires three boundaries to be held firmly. First, pure gauge theory must be kept distinct from full quantum chromodynamics, which contains dynamical fermions, and both must be kept distinct from the full Standard-Model gauge group $G_{\mathrm{SM}} = \mathrm{SU}(3)\times\mathrm{SU}(2)\times\mathrm{U}(1)$. Success in the first arena says nothing automatic about the second or third.\catmark{A} Second, the standard pure-gauge benchmark families, including the glueball spectrum~\cite{morningstar1999,michael1989,athenodorou2020,athenodorou2021}, the flow scales $t_0$ and $w_0$~\cite{luscher2010}, and the topological susceptibility $\chi_{\mathrm{top}}$~\cite{durr2007}, are to be treated as external benchmarks to be matched \emph{after} a source audit, and never as quantities the framework may quietly assume. Until such an audit is complete for a given observable, that observable's status is \Bpending.\catmark{B} Third, and as a standing limitation, the \UIDT\ spectral gap $\Delta$\,\Bdelta\ must not be naively identified with the mass of the lightest scalar glueball, the $0^{++}$ state. This is the limitation designated~L12, and it is developed in \cref{subsec:delta-glueball-l12}.

\begin{evidencebox}[Pure $\mathrm{SU}(3)$ as Stress Test]{catB}
\textbf{Claim.} Pure $\mathrm{SU}(3)$ Yang--Mills theory is the appropriate isolated sector against which the \UIDT\ gauge mapping is first
tested.\catmark{A}\\[0.3em]
\textbf{Status of the gap value.} $\Delta = \DeltaGap$ is reported as \catmark{B}\ under the recorded determination (\cref{box:pi-override-delta}): internally consistent, lattice-compatible, not externally peer-confirmed and not
derived.\\[0.3em]
\textbf{What is \emph{not} claimed.} Compatibility with the pure-gauge mass gap does not derive $G_{\mathrm{SM}}$, does not constitute a solution of the Yang--Mills existence-and-mass-gap problem, and does not by itself confirm any ontological thesis of \UIDT.\catmark{A}
\end{evidencebox}

\subsection{Topology and the Topological Susceptibility}
\label{subsec:topology-chi}

Among the pure-gauge benchmarks, the topological susceptibility $\chi_{\mathrm{top}}$ occupies a special position, because topology is where a structural-realist theory of gauge fields is most exposed. The quantity $\chi_{\mathrm{top}}^{1/4}$ is treated here as a pure-gauge topological benchmark \emph{seed}: a target for post-hoc comparison, carried at status \Bpending\ until its provenance is fully audited.\catmark{B}

The audit required before any such comparison may be taken seriously is demanding, and it is demanding for good reason. A bare number for $\chi_{\mathrm{top}}^{1/4}$ is meaningless without the apparatus that produced it. The framework therefore requires, for every external topological value it wishes to compare against, an explicit operator definition; the smoothing or gradient-flow prescription used to define the topological charge density; the set of lattice spacings entering the result; the continuum-extrapolation procedure; the units and scale-setting convention; the quoted uncertainty; and the covariance status relative to any other quantity used in the same comparison.\catmark{A} Only ratios in which the scale-setting ambiguity has been harmonised across the compared quantities are admissible, and even then only at predictive status~\catmark{D}.

Two prohibitions accompany these requirements. The first concerns a specific historical route. An early attempt derived topology through a leading-order sum-rule (SVZ-type) estimate; this route is recorded as a museum artifact and is admissible only as documentation of a failed approach, never as a live calculation.\catmark{A} The second is structural: agreement of $\chi_{\mathrm{top}}$ with a lattice value, however close, may never be presented as a closure of the $G_{\mathrm{SM}}$-origin problem. A topological benchmark is a constraint on the operator content and the scale bridge of the framework; it is not, and cannot be, a derivation of the Standard-Model gauge group.\catmark{A} A mismatch, correspondingly, is to be read first as a constraint on operator choice, truncation, or scale-setting, and only after those are exhausted, as a more serious tension, not as immediate, total refutation of the entire ontology.\catmark{A}

\begin{governancebox}[Topology: Admissibility Conditions]
A topological benchmark comparison is admissible only if it carries: operator definition; flow/smoothing prescription; lattice spacings; continuum extrapolation; units and scale-setting; uncertainty; covariance status. The SVZ leading-order route is failure documentation only. No $\chi_{\mathrm{top}}$ agreement closes the $G_{\mathrm{SM}}$-origin problem.
\end{governancebox}

\subsection{Calibration and the Discipline of \texorpdfstring{$\gamma$}{gamma}}
\label{subsec:gamma-discipline}

The invariant $\gamma = \GammaGeom$ is the most load-bearing single number in the framework, and for that reason it is also the number around which the strictest discipline is enforced. Its classification is fixed and not subject to upgrade: $\gamma$ is \catmark{A-}, calibrated. It is \emph{not} \catmark{A}, and in particular it is not derived from renormalisation-group first principles.\catmark{A-} The framework's own limitations register this as a standing open problem: the value emerges from algebraic self-consistency within the \UIDT\ construction, in a manner structurally analogous to a leading-order truncation scheme, but the discrepancy between a renormalisation-group-derived $\gamma$ and the calibrated kinetic value remains an unresolved question, openly recorded rather than concealed. To be precise about what ``calibrated'' means here: $\gamma$ is fixed by matching to the kinetic vacuum-expectation structure of $\Sx$ in the effective Lagrangian of \cref{eq:lagrangian}, a calibration to an internal kinetic reference, not a derivation from a renormalisation-group flow. The public-facing description of $\gamma$ must therefore read ``calibrated to the kinetic VEV,'' never ``derived from the kinetic VEV''; the latter phrasing violates the \catmark{A-} discipline and is corrected wherever it appears.\catmark{A-}

From the calibrated status of $\gamma$ follows the framework's central methodological safeguard against self-deception, the anti-target-leakage rule.
\newpage
\begin{governancebox}[Anti-Target-Leakage Rule]
A solver that has access to the value it is meant to predict cannot be said to predict it. Therefore the calibrated value $\gamma = 16.339$, together with all quarantined museum constants and all external benchmark values, must remain inaccessible to any code path that claims to compute $\gamma$, $\Delta$, or related quantities from within \UIDT. Blind computation and post-hoc comparison must be separated programmatically; benchmark values are stored with solver access disabled. A derivation that imports its own target is a tautology wearing the costume of a result.
\end{governancebox}

\noindent The rule is stated abstractly here and is given its formal treatment as a target-leakage theorem in Part~III. Its force in the present context is to constrain what a future derivation of $\gamma$ would have to look like in order to count. Any such derivation must specify, in advance and in public, the operator whose flow it integrates; the flow equation itself; the regulator; the boundary conditions; an explicit proof that no target value has leaked into the computation; and an executable reproduction command. Absent any one of these, a claimed derivation of $\gamma$ remains at predictive status~\catmark{D} at best, and more honestly is not a derivation at all.\catmark{A}

A final point of discipline concerns conflicting values. The corpus contains more than one number for $\gamma$: alongside the calibrated kinetic value $16.339$, a Monte-Carlo determination yields a mean of $\gamma_{\mathrm{MC}} = \GammaMC$. These are not to be reconciled by quiet selection of the more convenient figure, nor are they to be treated as interchangeable. They are recorded as a \emph{tension entry}: two determinations that must be displayed together, with their methods and uncertainties, so that the disagreement is visible and available for resolution.\catmark{A} The full tension ledger appears in \cref{sec:tension-ledger}; the point here is that a discrepancy between two internal determinations is information to be preserved, not noise to be filtered.

\noindent With the corpus discipline, the evidence grammar, the governance boundary, and the empirical baseline now in place, the framework has stated the terms of its own accountability. Part~II takes up the constructive task these terms have been protecting: the definition of the primitive hierarchy, beginning with the precise sense in which the vacuum information density $\Sx$ is, and is not, a foundation.

\begin{evidencebox}[The Status of $\gamma$]{catAminus}
\textbf{Classification.} $\gamma = \GammaGeom$ is \catmark{A-}\ (calibrated), permanently. It is never \catmark{A}\ and never described as
``derived.''\\[0.3em]
\textbf{Open problem.} A first-principles renormalisation-group derivation of $\gamma$ does not exist; the kinetic-versus-RG discrepancy is an openly recorded
limitation, not a closed result.\\[0.3em]
\textbf{Tension.} The Monte-Carlo value $\gamma_{\mathrm{MC}} = \GammaMC$ is held alongside $16.339$ as a recorded tension, not resolved by
selection.\\[0.3em]
\textbf{Leakage prohibition.} $\gamma = 16.339$ must never enter a blind solver path; any future derivation must satisfy the anti-target-leakage rule in full.
\end{evidencebox}
%=============================================================================
% PART II
%=============================================================================
\newpage
\part{The Formal UIDT Boundary and the Primitive Hierarchy}
\label{part:primitive-hierarchy}

\noindent Having fixed the terms of accountability, the manuscript now undertakes its constructive task: to state precisely what \UIDT\ takes as primitive, and to keep that statement free of the two failures that most often afflict foundational proposals, including hidden circularity, in which the primitive is secretly defined in terms of what it is meant to generate, and category confusion, in which a tool of mathematical description is mistaken for a stage in the genesis of being. Part~II separates the \emph{ontological} axis, the vertical line of states from which physical reality is held to crystallise, from the \emph{formal} layer of morphisms, functors, and categories that describe that line. The separation is the central methodological result of this part, and it is what allows the framework to use the full apparatus of modern mathematics without confusing the map for the territory.

%=============================================================================
% CHAPTER 3 (NEW) — AXIOMATICS AND THE EFFECTIVE LAGRANGIAN
%=============================================================================
\newpage
\section{Axiomatics and the Effective Lagrangian}
\label{sec:axiomatics}

\noindent Before the ontological hierarchy is developed in prose, the framework must make its formal commitments explicit. A book-length ontology that reasons about the spectral gap $\Delta$, the calibrated invariant $\gamma$, the coupling $\kappa$, and the self-coupling $\lambda_S$ owes the reader the Lagrangian in which those quantities actually live. This section states the four axioms of \UIDT, the explicit effective Lagrangian, and the dimensional bookkeeping that fixes the theory's status as an effective field theory. The axioms are not arbitrary postulates; each is motivated by a requirement of mathematical consistency or physical coherence, and each is tagged with its honest evidence class.

\subsection{The Four Axioms}
\label{subsec:four-axioms}

\begin{axiom}[Primacy of Information Density]
\label{ax:primacy}
There exists a real scalar field $\Sx$, the vacuum information density, of canonical mass dimension $[\Sx]=1$. Within the manuscript grammar of \cref{subsec:variant-b} (Variant~B), this field is taken as the sole fundamental degree of freedom; all spatiotemporal, geometric, and material structures are hypothesised, within that grammar, to be emergent consequences of its dynamics.\catmark{A/E}
\end{axiom}

\begin{axiom}[Gauge-Invariant Self-Interaction]
\label{ax:gauge}
The dynamics of $\Sx$ are governed by a Lagrangian density $\mathcal{L}_{\UIDT}$ invariant under $\mathrm{SU}(3)$ gauge transformations and respecting vacuum stability ($V''(v)>0$). It is an effective field theory valid below a cutoff $\Lambda_{\UIDT}$. The non-minimal coupling $(\bar\kappa/\Lambda_{\UIDT})\,S\,F_{\mu\nu}^a F^{a\mu\nu}$ has mass dimension~5 and is therefore not renormalisable by strict power counting; its validity is confined to the EFT domain (Limitation~L7). A full non-perturbative UV completion is an open research question.\catmark{A}
\end{axiom}

\noindent The minimal realisation satisfying \cref{ax:gauge} is the effective Lagrangian
\begin{equation}
 \mathcal{L}_{\UIDT}
 = -\tfrac{1}{4}F_{\mu\nu}^a F^{a\mu\nu}
 + \tfrac{1}{2}(\partial_\mu S)(\partial^\mu S)
 - \tfrac{1}{2}\mu^2 S^2
 - \frac{\lambda_S}{4!}S^4
 + \frac{\bar\kappa}{\Lambda_{\UIDT}}\,S\,F_{\mu\nu}^a F^{a\mu\nu},
 \label{eq:lagrangian}
\end{equation}
where $\bar\kappa$ is the dimensionless EFT-normalised non-minimal gauge--scalar coupling and $\lambda_S$ the scalar self-coupling; the dimension-5 operator therefore carries the dimensionful prefactor $\bar\kappa/\Lambda_{\UIDT}$, in agreement with \cref{tab:eft-dimensions}. Because $S$ is a real gauge singlet, its kinetic term uses the ordinary derivative $\partial_\mu$. In the framework's ledger and renormalisation-group statements this dimensionless coefficient is written simply $\kappa$ (so $\kappa\equiv\bar\kappa$, with ledger value $\kappa=\kappaVal$). The exact internal coupling constraint of the framework is $5\kappa^2 = 3\lambda_S$ (a candidate fixed-point condition, untested by any defined functional renormalisation-group flow), which for $\kappa = \tfrac{1}{2}$ gives the exact value $\lambda_S = \lambdaS$ and the vacuum expectation value $v = \vVEV$; like $\lambda_S$, $v$ is fixed by this exact constraint rather than independently measured, which is why it carries no uncertainty. This constraint closes to a residual below $10^{-14}$ at eighty-digit working precision (\cref{prop:target-leakage} discipline applies: the closure is an internal-consistency certificate, not an external agreement).

\begin{table}[H]
\scriptsize
\caption{Dimensional and EFT status of the terms entering \cref{ax:gauge}.
The dimension-5 operator marks \UIDT\ as an effective field theory.}
\label{tab:eft-dimensions}
\renewcommand{\arraystretch}{1.25}
\begin{tabular*}{\textwidth}{@{}p{0.31\textwidth}@{\extracolsep{\fill}}ccp{0.34\textwidth}@{}}
\toprule
\textbf{Term} & \textbf{Op.\ dim.} & \textbf{Coupling dim.} & \textbf{Renormalisability} \\
\midrule
$F_{\mu\nu}^a F^{a\mu\nu}$ & $4$ & dimensionless & Strictly renormalisable Yang--Mills kinetic term \catmark{A} \\
$(D_\mu S)(D^\mu S)$ & $4$ & dimensionless & Canonical scalar kinetic term \catmark{A} \\
$\mu^2 S^2$ & $2$ & $+2$ & Relevant scalar mass operator \catmark{A} \\
$\lambda_S S^4$ & $4$ & dimensionless & Power-counting renormalisable self-interaction \catmark{A} \\
$(\bar\kappa/\Lambda_{\UIDT})\,S F_{\mu\nu}^a F^{a\mu\nu}$ & $5$ & $-1$ & Non-renormalisable; EFT-only below $\Lambda_{\UIDT}$ \catmark{A} \\
\bottomrule
\end{tabular*}
\end{table}

\begin{axiom}[Geometric Emergence]
\label{ax:emergence}
The effective spacetime metric $g_{\mu\nu}^{\mathrm{eff}}$ is not fundamental but a derived quantity determined by the vacuum configuration of $\Sx$. A candidate parametrisation is
\begin{equation}
 g_{\mu\nu}^{\mathrm{eff}} = \eta_{\mu\nu}
 + \alpha\,\partial_\mu\Sx\,\partial_\nu\Sx
 + \beta\,\Sx^2\,\eta_{\mu\nu},
 \label{eq:metric-emergence}
\end{equation}
with $\eta_{\mu\nu}$ the Minkowski metric and $\alpha,\beta$ dimensionful constants fixed by the requirement that the emergent geometry reproduce general relativity in the low-energy limit. So that each additive term in $g_{\mu\nu}^{\mathrm{eff}}$ is dimensionless given $[\Sx]=1$, the constants carry mass dimension $[\alpha]=-4$ (since $[\partial_\mu\Sx\,\partial_\nu\Sx]=4$) and $[\beta]=-2$ (since $[\Sx^2]=2$). The specific form of \cref{eq:metric-emergence} and the values of $\alpha,\beta$ are model-dependent and carried at predictive status.\catmark{D}
\end{axiom}

\begin{axiom}[Falsifiability]
\label{ax:falsifiability}
The theory defines a finite set of experimentally testable predictions (\cref{sec:falsification-tension}), each with an explicit falsification threshold. If any critical prediction is excluded at the specified confidence level, the corresponding \UIDT\ mapping, calibration, or prediction must be revised or abandoned.\catmark{A}
\end{axiom}

\begin{remark}[Falsifiability and critical rationalism]
\label{rem:popper}
Including falsifiability as a formal axiom is a methodological commitment to the critical-rationalist tradition: a theory that cannot in principle be wrong has no claim to being right. \cref{ax:falsifiability} operationalises this by assigning every central prediction a quantitative falsification threshold. The axiom is not a physical law but a meta-theoretical commitment to epistemic honesty, continuous with the principle of \cref{subsec:governance-boundary} that internal consistency is not external confirmation.
\end{remark}

\subsection{The Ontological Hierarchy}
\label{subsec:ontological-hierarchy}

The architecture of \UIDT\ is strictly hierarchical: each level is held to emerge from the preceding one through mathematical consistency rather than by fiat. The ordering is (1) a logical substratum (formal logic, set theory, real analysis); (2) the information density $\Sx$ of \cref{ax:primacy,ax:gauge}; (3) a vacuum geometry generated by the self-interaction of $\Sx$, including the spectral gap $\Delta$ and the calibrated invariant $\gamma$; (4) an emergent effective metric (\cref{ax:emergence}); (5) material content as excitations above the structured vacuum; and (6) cosmological dynamics governed by the vacuum-energy hierarchy. The ordering is strict, and no level may be invoked to explain a level that precedes it: spacetime does not explain $\Sx$; within the framework's grammar, $\Sx$ explains spacetime.\catmark{A/E} The one decisive caveat, established in \cref{sec:d2-obstruction}, is that the passage from level~(2) to the gauge content of level~(5) is \emph{not} achieved by the scalar gradient alone; the $G_{\mathrm{SM}}$-origin gap of \cref{oq:gsm-gap} sits precisely at that transition.

\begin{evidencebox}[Axiomatic Core: Standing Status]{catA}
\textbf{Axioms.} Primacy \catmark{A/E}; Gauge-Invariant Self-Interaction \catmark{A} (EFT, dimension-5 operator); Geometric Emergence \catmark{D}
(model-dependent metric); Falsifiability \catmark{A} (meta-theoretical).\\[0.3em]
\textbf{Lagrangian.} \cref{eq:lagrangian} with $\kappa=\kappaVal$, $\lambda_S=\lambdaS$, $v=\vVEV$; internal coupling constraint $5\kappa^2=3\lambda_S$ closes
internally to $<10^{-14}$ at eighty digits.\\[0.3em]
\textbf{Status.} The Lagrangian is an effective field theory below $\Lambda_{\UIDT}$; renormalisability is explicit per \cref{tab:eft-dimensions}; the gauge sector is coupled, not generated (\cref{subsec:coarse-geometry-eft}).
\end{evidencebox}

%=============================================================================
% CHAPTER 4 (was 3) — UNMARKED STATE
%=============================================================================
\section{Unmarked State, First Distinction, and \texorpdfstring{$S(x)$}{S(x)}}
\label{sec:unmarked-distinction-sx}

\subsection{Why \texorpdfstring{$S(x)=0$}{S(x)=0} Is Not the Unmarked State}
\label{subsec:sx-zero-not-unmarked}

A foundational theory must be exact about its own beginning, and the most common error at the beginning is to confuse \emph{absence of value} with \emph{absence of structure}. \UIDT\ takes the unmarked state, ``nothing'' in the precise sense of Spencer-Brown's calculus of distinctions, as the condition in which no field-value assignment exists at all. It is essential to see that this is \emph{not} the same as the field configuration $S(x) = 0$.\catmark{A}

The point is structural, not rhetorical. To write $S(x) = 0$ is already to presuppose a great deal: that there is a field $S$; that there is a domain indexed by $x$ over which it is defined; that evaluation of the field at a point is a meaningful operation; and, most importantly, that the value $0$ is \emph{distinguishable} from other possible values. A vanishing field value is a marked state: it is the assertion ``the field here is zero'', and an assertion of that form lives entirely within an already-constituted system of field, domain, and distinguishability. The unmarked state precedes all of this. It is not the number zero; it is the absence of the apparatus within which numbers could be assigned.\catmark{A}

This distinction is recorded as a clean Class~A result of the honesty audit, and it carries an immediate disciplinary consequence: it forecloses a family of overreaches. Because the unmarked state is not a physical configuration, it cannot be assigned a physical history, a temperature, a fluctuation spectrum, or a cosmological epoch. One cannot say that ``in the beginning $S(x)=0$ and then it became non-zero,'' because that sentence smuggles a temporal ordering and a pre-existing field into a condition that, by construction, contains neither. The framework therefore explicitly refuses cosmological, theological, and metaphysical elaborations of the unmarked state. It is a logical boundary, not a primordial substance, and the prohibition on treating it as a substance is as binding as any numerical constraint in the theory.\catmark{A}
\newpage
\begin{evidencebox}[The Unmarked State]{catA}
\textbf{Claim.} The unmarked state is the absence of field-value assignment, not
the field value $S(x)=0$.\catmark{A}\\[0.3em]
\textbf{Reason.} $S(x)=0$ presupposes field, domain, evaluation, and distinguishability; it is therefore a marked state within an already-constituted
system.\catmark{A}\\[0.3em]
\textbf{Boundary.} The unmarked state admits no physical history, no fluctuation, no temperature, no epoch. Cosmological or theological readings of it are category errors and are excluded.\catmark{A}
\end{evidencebox}

\subsection{Variant~B: \texorpdfstring{$S(x)$}{S(x)} as the Primitive of the Manuscript Grammar}
\label{subsec:variant-b}

The relationship between the first distinction and the scalar field $S(x)$ admits two readings, and the developmental record was right to treat the choice between them as a potential source of circularity rather than a matter of phrasing. The question is sharp: \emph{is the first distinction an ontological entity more fundamental than $S(x)$, or is it a description of $S(x)$?}

Under the first reading, call it Variant~A, the distinction is strictly prior: the order of genesis runs first-distinction, then $S(x)$, then relations, then geometry. Under the second reading, Variant~B, the field $S(x)$ is itself the \UIDT\ primitive, and the first distinction is identified with the first non-trivial realisation of that field, namely the condition $S(x)\neq 0$. The hazard is that an uncritical mixture of the two readings produces a definitional circle, in which the distinction defines $S(x)$ while $S(x)\neq 0$ simultaneously defines the distinction.\catmark{A}

This manuscript adopts Variant~B, and it does so as a matter of \emph{internal ontological grammar}, not as a physical proof that the field is the unique primitive of nature. Within the grammar of \UIDT, $S(x)$ is taken as primitive, and $S(x)\neq 0$ is the first non-trivial realisation from which relations, and then geometry, are held to follow.\catmark{A} The honest qualification is essential and is stamped on the choice: that $S(x)$ is the primitive is a decision about how the framework speaks, classified \catmark{E}\ as an interpretive commitment; it is not a demonstration that reality must be organised this way.\catmark{A/E} Variant~B is canonical here precisely because the manuscript is written in its grammar; were the manuscript rewritten under Variant~A, the canonical choice would change with it. Naming the choice, and labelling it for what it is, is what keeps the circle from closing unnoticed.

A second discipline attaches to the word ``information,'' because the field is named the vacuum \emph{information} density and the term is dangerously overloaded. \UIDT's notion of information density must be held distinct from the several established and inequivalent senses of ``information'': it is not Shannon information (a measure on a probability distribution over messages); not semantic information (content with truth conditions); not Bateson's ``difference that makes a difference'' (a cybernetic-epistemic notion); not algorithmic information (Kolmogorov complexity); not thermodynamic entropy in the Boltzmann--Gibbs sense; and not epistemic information in the sense of an agent's knowledge state.\catmark{A} Within the framework, and consistently with the developmental record, information appears not as a building block of the universe but as the \emph{emergent counting-measure of difference}, the quantity by which distinctions, once present, are read off. Information attaches \emph{at} the first distinction; it is not prior to it.\catmark{A/E}

\begin{governancebox}[Terminology Guard: $S(x)$ and ``Information'']
The \UIDT\ information density $S(x)$ is not Shannon, semantic, Bateson, algorithmic, thermodynamic, or epistemic information. It is a real scalar field taken as primitive in the manuscript grammar (Variant~B). ``Information'' in \UIDT\ is the emergent counting-measure of difference, attaching at the first distinction. Under no reading is $S(x)$ identified with consciousness; that identification is a category error and is prohibited.\catmark{A}
\end{governancebox}

\noindent The final clause of the guard is not incidental. A theory that calls its primitive an ``information density'' invites the slide from information to mind, and the framework blocks that slide at the outset: $S(x)$ is not consciousness, and no later chapter is permitted to recover consciousness as a property of the bare field. The observer, when it appears in Part~V, will be a coarse-grained, compression-bound process built far downstream of $S(x)$, never a hidden attribute of the primitive itself.

\subsection{The Corrected Ontological Sequence: Morphism as Formal Tool, Not Generator}
\label{subsec:corrected-sequence}

The developmental dialogue at one stage placed ``morphism'' within the vertical sequence of ontological stages, as though a morphism were a step in the genesis of being. The honesty audit corrected this, and the correction is one of the more consequential conceptual results of the present edition. The corrected architecture rests on a separation between two axes: a vertical \emph{ontological} axis, which is the genesis of reality, and a horizontal \emph{formal} axis, which is the mathematical apparatus used to describe that genesis.\catmark{A}

The vertical ontological axis, read as a single unbroken line of structural states each arising necessarily from the previous one, is the following. From the unmarked state proceeds the first distinction; from the first distinction, relational differentiation; from relational differentiation, a relational or pre-graph structure; from there, a graph or causal-relational structure; then a coarse-grained geometry; then an effective field regime; then an effective object; and finally an observer or self-model.\catmark{E/D} Each later stage in this list is held at speculative or predictive status, and Parts~IV and~V will treat the difficult transitions explicitly; the point here is the \emph{order}, and the fact that morphism appears nowhere in it.

The horizontal formal axis is where the mathematics lives: set-theoretic labels, relations, morphisms, functors, categories, graphs, tensors, and differential-geometric structures. These are the instruments by which the vertical genesis is represented and computed.\catmark{A/E} The decisive observation is elementary and therefore robust. A morphism $f\colon A\to B$ presupposes a distinguishable domain $A$ and a distinguishable codomain $B$; the very act of writing down a morphism requires that two distinct objects already exist to be related. A morphism therefore cannot be the generator of distinction, because it presupposes distinction.\catmark{A} It belongs to the formal representation layer, not to the primitive genesis layer.

This yields a clean two-attachment picture of how the mathematical apparatus docks onto the ontological line. Information attaches at the first distinction, as the emergent counting-measure of difference. Morphisms and transformations attach later, at the stage of relations: once there are relations between distinctions, the topos-theoretic apparatus of morphisms becomes the means by which the dynamics of those relations are computed without circularity. Neither information nor morphism is a constituent of the genesis; each is an interface, the first a counting interface and the second a computational one.\catmark{A/E} The framework's own summary of the corrected sequence may be rendered as follows: the universe arises from a self-stabilising distinction that condenses into relations and crystallises, in the macroscopic limit, as spacetime and matter; information and morphisms are not the substance of this genesis but the cognitive and mathematical interfaces through which a thermodynamically isolated subsystem, an observer, decodes the resulting structure.

\begin{evidencebox}[Morphism Placement]{catA}
\textbf{Result.} Morphism is a formal-representation tool, not an ontological generator. A morphism $f\colon A\to B$ presupposes distinguishable $A$ and $B$,
hence presupposes distinction.\catmark{A}\\[0.3em]
\textbf{Two axes.} Vertical = ontological genesis (unmarked state $\to$ distinction $\to$ relations $\to$ pre-graph $\to$ graph/causal $\to$ coarse geometry $\to$ effective field $\to$ effective object $\to$ observer). Horizontal = formal apparatus (relations, morphisms, functors, categories,
tensors, differential geometry).\\[0.3em]
\textbf{Attachment.} Information docks at the first distinction; morphisms dock at relations. Both are interfaces, not constituents.\catmark{A/E}
\end{evidencebox}

\noindent A consequence of this separation is that the major programmes of categorical and relational foundations enter \UIDT\ as \emph{lenses}, not as definitions of the primitive. Topos theory, the causal-set programme, constructor theory, and QBism each supply a vocabulary and a set of constraints through which the vertical line may be viewed; none of them is permitted to redefine $S(x)$ itself, and none is adopted wholesale.\catmark{A/E} Topos theory contributes the morphism apparatus of the horizontal axis; causal sets offer a discretisation lens for the graph/causal stage; constructor theory contributes a language of possible and impossible transformations; QBism contributes a disciplined account of the observer's epistemic position. Each is a way of looking, held at interpretive status, and the framework's commitment is only to the vertical line of genesis and to the refusal to let any single formalism on the horizontal axis masquerade as that line.

%=============================================================================
% CHAPTER 4
%=============================================================================
\newpage
\section{From Relations to Graph, Geometry, and Effective Field}
\label{sec:relations-to-field}

\subsection{Relational Differentiation and Pre-Graph Structure}
\label{subsec:relational-differentiation}

The first distinction, on its own, is barely anything: it marks a single difference, the bare fact that something rather than nothing has been distinguished. A world does not follow from a single difference. What follows, if anything is to follow, is a \emph{multiplicity} of differences standing in relation to one another. \UIDT\ names the first such structured multiplicity \emph{relational differentiation}: the stage at which distinctions are not merely present but are differentiated \emph{with respect to each other}, so that the question ``how does this distinction stand to that one?'' becomes meaningful.\catmark{E}

It is tempting to reach immediately for the language of graphs, with nodes for distinctions, edges for relations; and the temptation must be resisted for one stage longer. Before a graph there is only a \emph{pre-graph} structure: a field of relational differences that has not yet been committed to any particular graph-theoretic formalism, with its fixed notions of vertex, edge, and incidence. The pre-graph stage is deliberately under-specified, because to specify it prematurely as a graph would be to import mathematical commitments (discreteness, a definite adjacency relation, a particular notion of locality) that the ontology has not yet earned.\catmark{E} The pre-graph is the territory; the graph will be one possible map of it.

This is the appropriate place to state a discipline that governs the entire upward construction. Graph notation, when it is finally introduced, is a \emph{representation} of relational structure and not the relational structure itself. The arrow one writes between two nodes is a mark on paper standing in for a relation in the ontology; it is not the relation.\catmark{E} Consequently, no quantitative claim about the relational structure: no count of edges, no degree distribution, no spectral property, may be asserted until the framework has explicitly specified what its nodes are, what its edges are, how edges are weighted, and what dynamics govern their formation and change.\catmark{D} Absent those specifications, graph language is heuristic only.

To keep this honesty visible in the notation itself, the framework uses the schematic $I \to R(I)$, a distinction $I$ giving rise to a relational structure $R(I)$ upon it, strictly as \emph{toy notation}. It is a placeholder that records the intended direction of construction without asserting that the map $R$ has been defined, that it is unique, or that it carries any particular mathematical properties.\catmark{E} The toy notation is a promissory note, not a theorem, and it is labelled as such wherever it appears.

\begin{evidencebox}[Relational Differentiation and the Pre-Graph]{catE}
\textbf{Stage.} Relational differentiation is the first structured multiplicity after the first distinction; the pre-graph is its representation \emph{prior} to
any graph-theoretic commitment.\catmark{E}\\[0.3em]
\textbf{Discipline.} Graph notation represents relational structure; it is not that structure. No quantitative graph claim is admissible until nodes, edges,
weights, and dynamics are explicitly defined.\catmark{D}\\[0.3em]
\textbf{Notation.} $I \to R(I)$ is toy notation only, a recorded direction of construction, not a defined map.\catmark{E}
\end{evidencebox}

\subsection{Graph and Causal-Relational Structure}
\label{subsec:graph-causal}

Once the discipline of the previous section is accepted, the graph stage can be introduced for what it is: a \emph{theorem target}, not a completed construction. The proposal that the relational structure of the world is, at an appropriate level of description, a causal-relational graph is a hypothesis the framework would like to make precise and ultimately to prove or refute; it is held at predictive status.\catmark{D} The honest form of the proposal is not ``the world is a graph'' but ``reconstructing an effective spacetime from a correlation network is a legitimate theorem target,'' and the difference between those two sentences is the difference between a claim and a research programme.\catmark{D}

Making the target precise requires naming the variables that a causal-relational graph would have to carry. These include, at minimum: classes of nodes (distinguishing, if the dynamics demand it, different kinds of distinction); weights on relations (encoding the strength or character of a relational difference); a causal order (a partial order expressing which relational facts constrain which others); a correlation structure (the pattern of statistical dependence among relational quantities); and information- or entropy-theoretic measures defined on the structure.\catmark{D} Each of these is a variable to be \emph{specified and justified}, not assumed; the list is a specification of work to be done, not a description of work completed.

A candidate primary observable suggests itself naturally, namely the two-point function $G(x,y) = \langle S(x)\,S(y)\rangle$, the correlator of the information density at two indices, since correlations are the most direct quantitative expression of relational structure. The framework records this as a candidate, not as an established primary observable: its canonical status must be confirmed against the authoritative ledger before it is used as the basis of any computation, precisely because adopting an observable prematurely is one of the ways a framework quietly fixes its own conclusions.\catmark{D} Until that confirmation, $G(x,y)$ is a proposed handle on relational structure, available for exploration but not for load-bearing claims.

One prohibition is sharp and carries Class~A force. There is to be no direct inference from phenomenology to graph structure: a report of subjective experience, a clinical observation, or any other phenomenological datum may not be read off as a fact about nodes and edges.\catmark{A} The admissible chain from phenomenology to ontology is long and passes through neurodynamic observables and operational constraints, as Part~VI will detail; the short-circuit from a report directly to a claim about the relational graph is forbidden. The link that \emph{is} permitted runs the other way and through coarse-graining: the graph structure connects to observer-accessible objecthood only via the coarse-graining map developed in the next section and in Part~V, never by direct identification.\catmark{E}

\subsection{Coarse-Grained Geometry and the Effective Field Regime}
\label{subsec:coarse-geometry-eft}

Geometry, in \UIDT, is not primitive. The continuum metric structure with which physics ordinarily begins is here a \emph{coarse-grained representation} of an underlying relational structure: the smooth geometry is what the relational graph looks like when one stands far enough back that individual relational differences blur into a continuum.\catmark{E} This is the sense in which \UIDT\ belongs to the pregeometric tradition: geometry is an output, recovered as a limit, rather than an input assumed at the start. The strongest external precedent for treating geometry this way is Jacobson's derivation of the Einstein field equations as a thermodynamic equation of state from horizon entropy~\cite{jacobson1995}, with the later entropic-gravity programme~\cite{padmanabhan2010,verlinde2011,chirco2014} as further context; the framework invokes these as a Stratum-II precedent that emergent, thermodynamically-derived geometry is a respectable position, not as a derivation imported into \UIDT.\catmark{E}

Downstream of geometry lies the \emph{effective field regime}: the level of description at which fields propagating on an effective geometry, governed by an effective Lagrangian, become the appropriate language. This regime is explicitly downstream and explicitly effective; it is not the primitive layer, and its variables are emergent rather than fundamental.\catmark{D/E} The ordering matters because it locates correctly what the framework is and is not entitled to claim about gauge structure.

Here the most important sentence of the chapter must be stated plainly, because it is the honest boundary of version~3.9. In its present form, \UIDT\ couples the scalar $S(x)$ to a Yang--Mills sector that is \emph{already present}: the gauge fields are part of the effective field-theoretic description, not something the scalar has been shown to generate. The framework's current gauge sector is an effective-field-theory regime, not a generative one.\catmark{A} The precise content of the v3.9 Lagrangian is fixed by \cref{eq:lagrangian}, on which this negative boundary rests.

From this follows the cleanest and most important negative result of the ontological core, stated without hedging.



\noindent It would be possible to disguise this boundary, to speak of the gauge sector as though it emerged from the scalar, and to let the reader infer a derivation that has not been performed. The framework refuses that move on principle. Stating that v3.9 is effective rather than generative in its gauge sector is not a concession wrung from the theory; it is the theory being precise about the location of its own frontier. The frontier itself, the \emph{$G_{\mathrm{SM}}$-Origin Gap}, is named and preserved as the central open problem, and the candidate algebraic routes that might one day cross it, together with the explicit fork between retaining a strict scalar primitive and enriching the primitive to an algebraic, matrix, or tensor object, are taken up in Part~IV.
Part~III first establishes the numerical and benchmark discipline without which no such crossing could ever be claimed honestly.

\begin{evidencebox}[The Effective-Field Boundary of v3.9]{catA}
\textbf{Regime.} The v3.9 gauge sector is an effective-field-theory regime, in which $S(x)$ is coupled to an already-present Yang--Mills sector. It is not a
generative regime.\catmark{A}\\[0.3em]
\textbf{Negative result.} \UIDT\ v3.9 does \emph{not} derive the Standard-Model gauge group $G_{\mathrm{SM}}$ from a structureless real scalar $S(x)\in\mathbb{R}$. This is recorded as a gap identification, not as a
failure to be hidden.\catmark{A}\\[0.3em]
\textbf{Consequence.} The transition from $S(x)$ to $G_{\mathrm{SM}}$ is a research programme, not a result; its formal treatment, and the open question of the minimal intermediate algebra, are the subject of Part~IV.\catmark{D}
\end{evidencebox}
%=============================================================================
% PART III
%=============================================================================
\newpage
\part{Numerical Architecture, Blind Benchmarks, and Stress Tests}
\label{part:numerical-architecture}

\noindent The previous part fixed what \UIDT\ takes as primitive and where its present frontier lies. This part fixes how the framework is permitted to compute, and how its computations may and may not be compared with the external world. The governing concern throughout is a single distinction that is easy to state and easy to violate: the difference between an algorithm that is internally consistent and a claim that is externally true. A calculation can be perfectly self-consistent to eighty digits and still describe nothing. Part~III builds the machinery, namely a hard separation of internal closure from external agreement, a claim schema with mandatory reproduction notes, and a theorem on target leakage, which is designed to make that difference impossible to overlook.

%=============================================================================
% CHAPTER 5
%=============================================================================
\newpage
\section{Numerical Epistemology}
\label{sec:numerical-epistemology}

\subsection{Internal Closure Versus External Agreement}
\label{subsec:closure-vs-agreement}

Two kinds of numerical statement are routinely confused, and the confusion has done real damage in the developmental history of this framework. The first is \emph{internal deterministic closure}: a calculation, carried out at fixed precision, returns a residual smaller than some threshold, demonstrating that the algorithm is self-consistent. The second is \emph{external agreement}: a computed quantity matches a measured or lattice-determined value to within the latter's uncertainty. These are not the same kind of thing, and no quantity of the first ever amounts to the second.\catmark{A}

Internal closure certifies an algorithm; it says that the computation does what it claims to do, repeatably and without numerical drift. External agreement certifies a correspondence with nature; it requires an external value, its uncertainty budget, its covariance with other quantities, the provenance of its source, and a demonstration that the computed quantity and the external quantity are genuinely the \emph{same} observable rather than two things that happen to carry similar numbers.\catmark{A} The first lives entirely inside the calculation. The second reaches outside it, and the reaching is where all the epistemic risk resides.

From this follows a hard rule, stated here as the governing numerical principle of the entire framework.

\begin{governancebox}[Hard Rule: Internal Closure Is Not External Agreement]
A deterministic internal residual below $10^{-14}$ is admissible \emph{only} as a certificate of internal closure, evidence that an algorithm is self-consistent at the working precision. It is \emph{never} a criterion of agreement with lattice or experimental data. Agreement with external data is established by a $\sigma$-level or covariance-aware comparison against a sourced, uncertainty-bearing value, never by an internal residual threshold. A residual of $10^{-14}$ against lattice data is a category error and is rejected.\catmark{A}
\end{governancebox}

\noindent The framework keeps a documented example of why this rule is not pedantry. An early construction produced the calibrated invariant by the exact arithmetic decomposition $49/3 + 17/3000 = 16.339$. The arithmetic is exact: it closes to arbitrary precision, and an internal residual check would report success at any number of digits. Yet the decomposition is a museum artifact: the integer $16339$ is prime, which proves that the split into $49/3$ and $17/3000$ was \emph{constructed} to land on the target rather than \emph{found} by any mechanism, and the construction has no physical origin. Perfect internal closure, in this case, certifies only that someone did the arithmetic correctly. It is the cleanest possible illustration of the rule: closure is about the algorithm, not about the world.\catmark{A}

Two implementation disciplines protect the rule in code. First, the working precision, eighty decimal digits, is set \emph{locally}, within each proof-critical computation, and is never centralised into a global configuration value. Centralising the precision would make it a single point of silent failure and would invite refactoring that treats a physics-critical constant as an ordinary tunable; both are prohibited, and global precision mutation is a rejected implementation pattern.\catmark{A} Second, proof-critical computations must avoid binary floating-point heuristics: machine-float evaluation, and rounding operations applied to quantities whose exact value carries the claim, are forbidden in any code path that a Class~A or Class~B result depends upon. High precision is a tool for internal consistency checking; it earns the framework nothing about nature, and it must not be allowed to masquerade as if it did.\catmark{A}

\subsection{The Claim Graph and Reproduction Notes}
\label{subsec:claim-graph}

If internal closure is not to be mistaken for truth, then every numerical claim must travel with the information needed to judge what kind of claim it is. \UIDT\ therefore requires that each numerical assertion be a node in a \emph{claim graph}, carrying a fixed set of fields, and that this schema be implemented \emph{before} any theory-generating computation is trusted; the accounting comes first, not as an afterthought.\catmark{A}

\begin{governancebox}[Required Fields for Every Numerical Claim]
\textbf{(1)} Claim identifier. \quad \textbf{(2)} Observable (precisely what physical or mathematical quantity is asserted). \quad \textbf{(3)} Value. \quad \textbf{(4)} Units. \quad \textbf{(5)} Uncertainty (mandatory for every empirical value). \quad \textbf{(6)} Method. \quad \textbf{(7)} Source (with verifiable identifier for external values). \quad \textbf{(8)} Evidence class (A through E). \quad \textbf{(9)} Tension status (whether the claim conflicts with another). \quad \textbf{(10)} Reproduction command (executable).
\end{governancebox}

\noindent The tenth field is the one most often missing and the one that most distinguishes an auditable claim from an assertion. A reproduction note must be an \emph{executable} command that a reader can run to regenerate the value, including the local precision setting, the input data, and the exact computation. A value whose reproduction note cannot be run, or which silently depends on an unstated global state, is not yet a claim in the technical sense; it is a number awaiting one.\catmark{A} The reproduction note is what converts a private calculation into a public, falsifiable object.

To make the schema concrete rather than merely prescribed, the framework's simplest \catmark{A}-class claim, the exact internal coupling constraint $5\kappa^2 = 3\lambda_S$ of \cref{eq:lagrangian} (a candidate fixed-point condition, untested by any defined flow), is accompanied by an executable reproduction note. At eighty-digit working precision, with $\kappa = \tfrac{1}{2}$ and $\lambda_S = \tfrac{5}{12}$, the residual $\lvert 5\kappa^2 - 3\lambda_S\rvert$ evaluates to zero:
\begin{verbatim}
python3 -c "from mpmath import mp,mpf; mp.dps=80;
k=mpf(1)/2; lS=mpf(5)/12; print(abs(5*k**2 - 3*lS))"
# -> 0.0 (residual < 1e-14: internal closure, not external agreement)
\end{verbatim}
This is the template every numerical claim must follow: a command a reader can run, a stated precision, and, per \cref{subsec:closure-vs-agreement}, an explicit reminder that the vanishing residual certifies internal consistency, not agreement with nature.\catmark{A}

The claim graph connects to the advisory claim-check procedure introduced in \cref{subsec:evidence-grammar}. That procedure inspects a proposed claim for the presence of prestige language, evidence overreach beyond the stated class, any cosmological quantity pushed above the Class~C cap, the naive spectral-gap--glueball identification (L12), missing references or missing reproduction commands, and reliance on quarantined artifacts. Its verdicts are advisory in the strict sense established in \cref{subsec:governance-boundary}: the procedure may \emph{flag} a claim, but a clean pass authorises nothing. Passing the claim-check is necessary hygiene, never sufficient warrant.\catmark{A}

\subsection{The Target-Leakage Theorem}
\label{subsec:target-leakage-theorem}

The anti-target-leakage rule introduced for $\gamma$ in \cref{subsec:gamma-discipline} can now be stated in its general form, because it is not a fact about $\gamma$ in particular but a structural fact about prediction itself.

\begin{proposition}[Target Leakage]
\label{prop:target-leakage}
Let a computational procedure $P$ be claimed to predict a quantity $q$ independently of its known value $q^{\ast}$. If $q^{\ast}$ is accessible to $P$, as a hard-coded constant, a calibration target, a falsification threshold, a loss function, or any other input that $P$ can read, then the agreement of $P$'s output with $q^{\ast}$ provides no evidence that $P$ predicts $q$. The agreement is, at most, a consistency check on $P$'s ability to reproduce a number it was given.
\end{proposition}

\noindent The proposition is almost a tautology, which is precisely its strength: it cannot be evaded by cleverness, only by architecture. Its consequences for the framework are concrete and binding. The calibrated value $\gamma = 16.339$, every quarantined museum constant, and every external benchmark value must be \emph{inaccessible} to any code path that purports to compute $\gamma$, $\Delta$, or related quantities from within \UIDT.\catmark{A} Benchmark values are stored with solver access explicitly disabled, so that a blind computation cannot consult the answer it is meant to produce.

This entails a programmatic separation that is not negotiable: the blind \UIDT\ computation and the post-hoc comparison against external data must be \emph{two distinct programs}. The first produces a number knowing nothing of the target; the second, only afterward, compares that number to a sourced, uncertainty-bearing benchmark and reports a $\sigma$-level. Merging the two, so that the target value is present in the environment where the prediction is formed, reintroduces exactly the leakage the separation exists to prevent.\catmark{A}

The framework keeps documented instances of the failure mode this theorem names. A hard-coded target ratio inserted directly into a solver, and a falsification condition defined as a threshold on the distance between a flow's output and the calibrated $\gamma$, are both recorded as rejected target-leakage patterns; in each case, successive adjustments that nudged a result toward its target were correctly identified not as convergence but as the solver being steered by the answer it could see.\catmark{A} The lesson generalises into the architecture: the only honest derivation of a known quantity is one performed by a procedure that demonstrably could not have known it.

\begin{evidencebox}[Numerical Epistemology: Standing Rules]{catA}
\textbf{Closure $\neq$ agreement.} $<10^{-14}$ certifies internal closure only; external agreement requires sourced, uncertainty- and covariance-aware
$\sigma$-comparison.\catmark{A}\\[0.3em]
\textbf{Local precision.} Eighty-digit precision is set locally per proof-critical block; global precision mutation and proof-critical
floating-point/rounding heuristics are prohibited.\catmark{A}\\[0.3em]
\textbf{Claim schema first.} Every numerical claim carries ten fields including an executable reproduction note; the schema precedes any theory
generator.\catmark{A}\\[0.3em]
\textbf{Target leakage (\cref{prop:target-leakage}).} A solver with access to its target predicts nothing; blind computation and post-hoc comparison are separate programs, benchmarks stored with solver access disabled.\catmark{A}
\end{evidencebox}

%=============================================================================
% CHAPTER 6
%=============================================================================
\newpage
\section{Benchmark Batteries}
\label{sec:benchmark-batteries}

\subsection{The Pure \texorpdfstring{$\mathrm{SU}(3)$}{SU(3)} Ratio Battery}
\label{subsec:ratio-battery}

The disciplined way to compare a candidate framework against pure-gauge lattice results is not to compare dimensionful quantities directly, because those are contaminated by scale-setting conventions that differ between calculations, but to compare \emph{dimensionless ratios} in which the common scale cancels. Constructing external blind ratio tests is the correct research direction for \UIDT, and dimensionless ratios are the correct audit architecture.\catmark{D}

A natural battery of such ratios is available from the standard pure-$\mathrm{SU}(3)$ observables. Using the gradient-flow scale $t_0$ to set the units, and the $\overline{\mathrm{MS}}$ scale $\Lambda_{\overline{\mathrm{MS}}}$ where a perturbative bridge is appropriate, the candidate ratios include the flow-normalised scalar glueball mass $m_{0^{++}}\sqrt{8t_0}$, the flow-normalised topological scale $\chi_{\mathrm{top}}^{1/4}\sqrt{8t_0}$, and the corresponding $\Lambda$-normalised forms $m_{0^{++}}/\Lambda_{\overline{\mathrm{MS}}}$ and $\chi_{\mathrm{top}}^{1/4}/\Lambda_{\overline{\mathrm{MS}}}$.\catmark{B pending} Each of these is carried at status \Bpending: it is a legitimate post-hoc comparison target whose external value must first survive the source and covariance audit of \cref{subsec:topology-chi} before any comparison is taken seriously.

The purpose of the ratios is to reduce, and where possible eliminate, the scale-setting ambiguity that would otherwise dominate any comparison; a ratio in which the common scale cancels is far more informative about structure than a dimensionful number whose value is partly a choice of units.\catmark{D} Two requirements govern their use. First, every external ratio must arrive with a verified source and an explicit covariance model, because the ratios share inputs (the same $t_0$, the same lattice ensembles) and are therefore correlated; treating correlated ratios as independent would manufacture false significance.\catmark{A} Second, and as a direct application of the target-leakage theorem, these benchmark values must remain inaccessible to any blind \UIDT\ solver: they are post-hoc comparators, stored with solver access disabled, never inputs to the computation whose output they are meant to test.\catmark{A}

A final point concerns the interpretation of a mismatch. If a \UIDT\ ratio fails to match its lattice counterpart, the first reading is not that the entire ontology is refuted. A mismatch is, in the first instance, a constraint on the operator chosen, the truncation employed, or the scale bridge assumed; only after those have been systematically examined and found adequate does a persistent mismatch rise to the level of a serious tension against the framework as a whole.\catmark{A/D} This is not special pleading; it is the ordinary logic of comparing an approximate calculation to a precise measurement, but it is stated explicitly so that neither a match nor a mismatch is over-interpreted.

\subsection{\texorpdfstring{$\Delta$}{Delta}, the Glueball, and the L12 Separation}
\label{subsec:delta-glueball-l12}

The spectral gap $\Delta$\,\Bdelta\ and the mass of the lightest scalar glueball, the $0^{++}$ state of pure $\mathrm{SU}(3)$, carry numerically similar values, and that numerical proximity has repeatedly invited a conflation that the framework must refuse. The two quantities are not, without further work, the same object.\catmark{A}

The distinction is one of category. The \UIDT\ gap $\Delta$ is, in the first instance, a mathematical feature of the vacuum information density, a spectral gap in the sense internal to the framework's fixed-point construction. The $0^{++}$ glueball mass is a particle-spectrum quantity, defined by the lattice through specific interpolating operators with definite quantum numbers and extracted from the long-distance decay of correlation functions. To assert that these are the same physical quantity requires \emph{operator matching}: a demonstration that the \UIDT\ object and the lattice glueball are created by, and couple to, the same operator content. No such matching has been established.\catmark{A} The proximity of the numbers is therefore, on its own, evidence of nothing; two distinct constructions are permitted to yield similar values by coincidence, and the burden of proof lies entirely on any claim that they coincide for a reason.

This is recorded as a standing limitation of the framework.

\begin{limitation}[Naive $\Delta\!\leftrightarrow\!0^{++}$ Identification: L12]
\label{lim:L12}
The \UIDT\ spectral gap $\Delta$ must not be identified with the mass of the lightest scalar ($0^{++}$) glueball of pure $\mathrm{SU}(3)$ Yang--Mills theory in the absence of explicit operator matching. Numerical proximity of the two values does not constitute identification. Any future relation between $\Delta$ and the glueball spectrum must be established by matching interpolating operators and quantum numbers, not by coincidence of magnitude.\catmark{A}
\end{limitation}

\begin{remark}[On the evidence class of the gap value]
Two distinct questions must be kept apart. The \emph{value} $\Delta = \DeltaGap$ is reported, under the recorded evidence-class determination (\cref{box:pi-override-delta}), at Class~\catmark{B}: internally consistent and lattice-compatible, not externally peer-confirmed and not derived. This edition's conservative audit independently notes that, absent a completed external verification, the value would otherwise sit at \catmark{E\,pending}; the determination sets the working class to~B while leaving that caution on the record. The \emph{identification} of the gap with the $0^{++}$ glueball is a separate matter and is barred outright by L12, regardless of the value's class. A quantity may be lattice-compatible in magnitude while its identification with a specific particle state remains unproven.
\end{remark}

\noindent The general discipline here is to keep mathematical vacuum-pole language distinct from the language of the lattice particle spectrum. A pole in a \UIDT\ correlator and a state in the glueball spectrum are described in different vocabularies for good reason; collapsing the two vocabularies is precisely the move L12 forbids. The framework may aspire to relate them, that is a legitimate predictive target, but the relation must be earned by operator matching and carried at predictive status until it is.\catmark{D}

\subsection{Gauge-Invariant Exact Renormalisation as Scaffold}
\label{subsec:gferg-scaffold}

A recurring hope in the developmental record is that a gauge-invariant exact renormalisation-group flow, designated GFERG in the corpus, might supply the missing derivation of $\gamma$, of $\Delta$, or of the coupling ratio $K_S$. The framework's position on this is deliberately austere: GFERG is admitted only as a \emph{scaffold}, a candidate gauge-invariant exact-renormalisation route to be defined precisely before any use, and it is explicitly \emph{not} a completed \UIDT\ derivation of any of those quantities.\catmark{A}

For an exact-renormalisation route to count as more than vocabulary, it must specify, in advance, every element that makes it a definite calculation: the flow equation itself (in the Wetterich or Polchinski family of exact renormalisation-group equations, used here only as named scaffolds with no derivation imported); the regulator; the truncation of the effective action; the boundary and initial conditions of the flow; the observables to be extracted; and the failure criteria under which the route would be judged to have failed. Without these, an equation written in renormalisation-group notation is not a renormalisation-group calculation.\catmark{A}

The framework therefore makes an explicit deletion. A toy ordinary differential equation that was, in the developmental record, labelled as a next-to-leading-order GFERG result is removed from the status of a claimed derivation. It does not specify a regulator, a truncation, or boundary conditions in any form that would make it a controlled approximation to an exact flow; it is a toy ODE wearing renormalisation-group vocabulary, and the vocabulary is precisely what made it look like more than it was.\catmark{A} Retaining it as a claimed derivation would violate the most basic requirement of the numerical epistemology of \cref{sec:numerical-epistemology}: that the apparatus of a calculation be real, not nominal.

\begin{governancebox}[GFERG: Admissibility as Scaffold]
GFERG is a candidate gauge-invariant exact-renormalisation route, admitted as a scaffold only.\catmark{D} It is not a derivation of $\gamma$, $\Delta$, or $K_S$. Any use requires an explicit flow equation, regulator, truncation, boundary conditions, observables, and failure criteria. Renormalisation-group vocabulary may not be used to present a toy ODE as a flow; the previously NLO-labelled toy ODE is deleted as a claimed derivation.\catmark{A}
\end{governancebox}

\begin{evidencebox}[Benchmark Batteries: Standing Status]{catB}
\textbf{Ratios.} The pure-$\mathrm{SU}(3)$ ratio battery ($m_{0^{++}}\sqrt{8t_0}$, $\chi_{\mathrm{top}}^{1/4}\sqrt{8t_0}$, and the $\Lambda_{\overline{\mathrm{MS}}}$-normalised forms) is held at \Bpending, post-hoc, with verified sources, covariance models, and solver access
disabled.\catmark{B pending}\\[0.3em]
\textbf{L12.} $\Delta$ is not identified with the $0^{++}$ glueball without operator matching; the value's class is~\catmark{B}\ under the recorded determination, a
separate matter from the barred identification.\catmark{A}\\[0.3em]
\textbf{GFERG.} Scaffold only; not a derivation; the toy NLO ODE is deleted as a claimed result.\catmark{A}
\end{evidencebox}

\subsection{The Limitation Register (L1--L12)}
\label{subsec:limitation-register}

The framework's discipline requires that its known limitations be enumerated in one place rather than scattered or implied. The register below collects them. The entries L1--L7 and L12 are textually anchored in the corpus and the preceding chapters; four entries (L8--L11) are referenced in the developmental record but could not be recovered to a definite wording, and they are therefore marked \textsc{unknown} rather than reconstructed, because fabricating a limitation would be as much an integrity failure as concealing one.\catmark{A}

\begin{limitation}[L1: The geometric scale factor (ill-defined)]
\label{lim:L1}
The framework employs an effective geometric scale factor as a calibration input bridging microscopic vacuum dynamics and macroscopic phenomenology. With the reference scales made explicit, the ratio is $\lambda_{\UIDT}/r_{\mathrm{conf}} = 0.660\,\mathrm{nm}/0.197\,\mathrm{fm} \approx 3.35\times10^{6} \approx 10^{6.5}$; earlier versions quoted ``a factor of order $10^{10}$'' without specifying the compared scales, and that figure is superseded, since no Standard-Model energy ratio yields exactly $10^{10}$. No first-principles derivation of this factor from $\Sx$ is presently supplied; it is an adopted scaling assumption whose value depends on which ultraviolet and infrared scales are compared, not a derived consequence of the axioms.\catmark{D}
\end{limitation}

\begin{limitation}[L2: Electron-mass discrepancy]
\label{lim:L2}
The lepton-sector mapping retains an electron-mass discrepancy of order $23\%$ when the present universal $\gamma$-scaling architecture is applied. This mismatch is not removed by the current vacuum-information-density formalism and must not be absorbed into the geometric scale factor of~L1. L2 is therefore an explicit open limitation of electroweak/leptonic integration, not an additional statement of the holographic suppression scale recorded in L1.\catmark{C}
\end{limitation}

\begin{limitation}[L3: Vacuum-energy parameter]
\label{lim:L3}
A numerical parameter of order $2.3$ invoked in the vacuum-energy-density discussion is presently unanchored: it is phenomenologically fitted to match observational expectations rather than derived from a renormalisation-group flow or from discrete lattice-geometry constraints.\catmark{C}
\end{limitation}

\begin{limitation}[L4: $\gamma$ not RG-derived]
\label{lim:L4}
The invariant $\gamma = \GammaGeom$ is permanently \catmark{A-}\ (calibrated). A first-principles renormalisation-group derivation does not exist; the discrepancy between the calibrated kinetic value and an RG-flow value is openly recorded (\cref{subsec:gamma-discipline}). Any future derivation must satisfy the explicit anti-target-leakage criterion of \cref{prop:target-leakage}.\catmark{A-}
\end{limitation}

\begin{limitation}[L5: The $N=99$ RG-step count]
\label{lim:L5}
The choice of $N=99$ discrete renormalisation-group cascade steps is empirical: it is selected to reach the target magnitude, with no geometric or mathematical principle dictating exactly $99$ steps. Related downstream predictions (such as a Casimir scale near $\lambda \approx 0.66\ \mathrm{nm}$) presently lack direct experimental corroboration.\catmark{D}
\end{limitation}

\begin{limitation}[L6: Glueball identification withdrawn]
\label{lim:L6}
The assertion that an exact renormalisation-group flow uniquely identifies the spectral gap with a specific physical glueball state is withdrawn from evidentiary use (as of the recorded withdrawal date). The toy ordinary differential equation once dressed in renormalisation-group language is deleted as a claimed result (\cref{subsec:gferg-scaffold}); residual mentions are historical only and carry no support.\catmark{E}
\end{limitation}

\begin{limitation}[L7: EFT validity and the torsion scale]
\label{lim:L7}
The non-minimal coupling $(\bar\kappa/\Lambda_{\UIDT})\,S\,F_{\mu\nu}^a F^{a\mu\nu}$ is of mass dimension~5 and strictly non-renormalisable by power counting (\cref{tab:eft-dimensions}); the \UIDT\ Lagrangian is therefore an effective field theory valid only below a cutoff $\Lambda_{\UIDT}$. For the EFT to remain valid as a stress-test arena for the spectral gap $\Delta = \DeltaGap$, the cutoff $\Lambda_{\UIDT}$ must lie above that gap scale; it is therefore \emph{not} identified with the torsion scale $E_T = \ETorsion$. The torsion scale is a separate, much lower observer-context quantity (\cref{subsec:et-scale-bridge}) and does not set the gauge-sector cutoff; the value of $\Lambda_{\UIDT}$ itself is not fixed in v3.9 and a full non-perturbative UV completion remains an open research question.\catmark{D}
\end{limitation}

\begin{limitation}[L8--L11: Unrecovered register entries]
\label{lim:L8toL11}
Four further limitations (L8, L9, L10, L11) are referenced in the developmental record but could not be recovered to a definite wording from the verified corpus in the present pass. They are marked \textsc{unknown} and must remain so until the underlying source notes are restored; they are not reconstructed, so as not to fabricate undocumented claims.\catmark{E}
\end{limitation}

\begin{limitation}[L13: Background independence only partial]
\label{lim:L13}
This entry is newly identified in the present audit pass; it is \emph{not} a reconstruction of any unrecovered entry above. Although \UIDT\ belongs to the pregeometric tradition and aspires to background independence, its concrete axiom of geometric emergence (\cref{ax:emergence}) is written as a perturbation about a fixed Minkowski metric, $g_{\mu\nu}^{\mathrm{eff}}=\eta_{\mu\nu}+\alpha\,\partial_\mu\Sx\,\partial_\nu\Sx+\beta\,\Sx^2\eta_{\mu\nu}$, so that $\eta_{\mu\nu}$ is inserted by hand to contract indices rather than derived dynamically from a purely relational phase. Full background independence---deriving $\eta_{\mu\nu}$ from the relational substrate without assuming it---has therefore not been established; the comparison table (\cref{tab:pregeometric-comparison}) accordingly records \UIDT's background independence as \emph{partial}. A related external constraint is the Weinberg--Witten theorem~\cite{weinbergwitten1980}, which forbids a Lorentz-covariant massless spin-2 state in a theory with a Lorentz-covariant conserved stress-energy tensor on a fixed background; this bars any interpretation of the emergent metric as a composite spin-2 graviton in the naive flat-space sense, and is consistent with (not a refutation of) the thermodynamic, non-particle reading of emergent geometry that \UIDT\ borrows from Jacobson~\cite{jacobson1995}. The emergent-gravity literature catalogues exactly which assumptions a viable construction must give up to evade the theorem~\cite{jenkins2009,carlip2014}---most plausibly, for \UIDT, the absence of a fundamental Lorentz-covariant stress-energy tensor on a pre-geometric substrate---and \UIDT\ has not yet specified which loophole its construction exploits. Resolving L13 requires a dynamical origin for $\eta_{\mu\nu}$ and an emergent-geometry construction that does not run afoul of the Weinberg--Witten assumptions. This remains an open \catmark{D}\ problem.\catmark{D}
\end{limitation}

\begin{limitation}[L14: Scalar-glueball spectroscopy and the status of $m_S$]
\label{lim:L14}
This entry is newly identified in the present audit pass. Beyond the operator-matching discipline already required by L12 (\cref{lim:L12}), recent lattice work bears directly on whether any sub-$2\,\mathrm{GeV}$ scalar state can support the framework's candidate resonance $m_S = 1.705 \pm 0.015\GeV$. An exploratory study using glueball, meson, and meson--meson operators is summarised by Morningstar as suggesting that no scalar state below about $2\,\mathrm{GeV}$ can be regarded as predominantly a glueball state~\cite{morningstar2024}; this is an exploratory, explicitly hedged single study and does not by itself overturn the pure-gauge consensus that places a $0^{++}$ state near $1.5$--$1.7\,\mathrm{GeV}$~\cite{athenodorou2020,morningstar1999}, but it sharpens the concern. If the lattice picture continues to move away from a sub-$2\,\mathrm{GeV}$ ``predominantly glueball'' interpretation, then $m_S$ must be read as a \emph{signature} the framework would have to confront, not as a confirmed prediction, and the empirical grounding of $\Delta$ becomes correspondingly under-determined. The resonance is therefore carried at \catmark{D}\ and its falsification threshold (\cref{sec:falsification-tension}) is to be read against the evolving lattice spectrum rather than a fixed benchmark.\catmark{D}
\end{limitation}

\begin{limitation}[L15: Non-Abelian Casimir geometry versus the scalar prediction]
\label{lim:L15}
This entry is newly identified in the present audit pass. The Casimir-scale anomaly near $\lambda \approx 0.66\,\mathrm{nm}$ (\cref{lim:L1} context; \catmark{D}) is, in the framework, motivated by a scalar-field treatment. Recent first-principles $\mathrm{SU}(3)$ lattice gauge-theory studies of the non-Abelian Casimir effect complicate any naive scalar reading: between chromometallic plates the gluonic Casimir energy is well described by a massive effective scalar with a ``glueton'' mass $m_{\mathrm{gt}} \approx 0.49\,\mathrm{GeV}$~\cite{chernodub2023}, but in tube and box geometries the non-Abelian force is found to be \emph{attractive} where a massless scalar field would predict repulsion~\cite{ngwenya2025}. The qualitative geometry-dependence of the non-Abelian result means a single scalar-field Casimir picture cannot be assumed to carry over to the confining gauge vacuum; the $0.66\,\mathrm{nm}$ figure must therefore be treated as a scalar-sector prediction awaiting both a non-Abelian-consistent derivation and direct sub-nanometre metrology, not as an established gauge-sector consequence. This remains an open \catmark{D}\ problem.\catmark{D}
\end{limitation}

\Needspace{4\baselineskip}
\noindent Limitation~L12 (the naive $\Delta\!\leftrightarrow\!0^{++}$ identification) is stated in full at \cref{lim:L12} above and is not repeated here. The register is thus complete in enumeration: anchored where the corpus permits (L1--L7, L12, and the newly identified L13--L15), and honestly blank where it does not (L8--L11).\catmark{A}


%=============================================================================
% PART IV

%=============================================================================
\newpage
\part{The Gauge-Origin Problem and the Algebraic Fork}
\label{part:gauge-origin}

\noindent This part confronts the framework's hardest internal obstacle directly and without softening. If the vacuum information density $S(x)$ is to be more than a scalar coupled by hand to a pre-existing gauge theory; if it is to be, in any sense, the origin of gauge structure, then there must be a route from the scalar to a non-trivial gauge field strength. Part~IV shows that the most natural such route fails for a precise and decisive mathematical reason, records that failure in the ledger as a permanent result rather than hiding it, and then sets out the fork in the road that the failure forces: either the primitive remains a strict real scalar and the gauge sector stays effective, or the primitive is enriched to an algebraic object rich enough to carry the structure that a scalar cannot. The part neither pretends the problem is solved nor treats it as fatal. It treats it as the central open problem it is.
\newpage
%=============================================================================
% CHAPTER 7
%=============================================================================
\section{The \texorpdfstring{$d^2 = 0$}{d2=0} Obstruction}
\label{sec:d2-obstruction}

\subsection{The Failed Scalar-Gradient Route to Curvature}
\label{subsec:failed-scalar-gradient}

The simplest conceivable way to extract a gauge field from a scalar is to identify the gauge potential with the gradient of the scalar. In the language of differential forms, one would set the connection one-form $A$ equal to the exterior derivative of the scalar, $A = \dd S$, and then attempt to read off a field strength as the curvature two-form $F = \dd A$. The attractiveness of this route is its economy: it asks nothing of the world beyond the scalar already posited. Its failure is correspondingly fundamental, and it is worth stating with full precision because the framework treats it as a decisive result rather than a disappointment.

\begin{proposition}[Vanishing Curvature of an Exact Connection]
\label{prop:d2-obstruction}
Let $S$ be a smooth real scalar field and let the connection one-form be defined as its exterior derivative, $A = \dd S$. Then the curvature two-form vanishes identically:
\[
 F \;=\; \dd A \;=\; \dd(\dd S) \;=\; \dd^2 S \;=\; 0,
\]
since the exterior derivative satisfies $\dd^2 = 0$ on any smooth manifold. Consequently the scalar-gradient construction yields only a flat, pure-gauge connection and cannot produce a non-trivial Yang--Mills field strength.\catmark{A}
\end{proposition}

\noindent The result is elementary, and that is exactly its force. The nilpotency of the exterior derivative, $\dd^2 = 0$, is among the most secure facts in differential geometry; it is the algebraic heart of de~Rham cohomology and is not the sort of thing a clever choice of variables can evade. A one-form that is exact, a gradient, is automatically closed, and its curvature is zero by identity, not by approximation. The scalar-gradient route therefore does not \emph{approximately} fail or fail \emph{under some conditions}; it fails identically, everywhere, for any smooth scalar.\catmark{A}\footnote{\textbf{Audit note (Clebsch loophole).} A known mathematical loophole to the $F \equiv 0$ obstruction exists via Clebsch parameterisations built from higher jets of $S$: with functionally independent scalars $u,v$ constructed from $S$ and its derivatives (e.g.\ $S$, $(\partial S)^2$, $\Box S$), the one-form $A = u\,\dd v$ yields a non-vanishing composite \emph{abelian} curvature $F = \dd u \wedge \dd v$.\catmark{A} This loophole does not rescue the gauge-origin programme: the construction is tensorial, providing no gauge redundancy and no gauge orbit; it offers no non-abelian structure, which would require of order $\dim\mathfrak{g}$ functionally independent scalar functionals; and it supplies no chiral matter. The obstruction recorded here therefore remains exactly as stated: construction-specific, with the naive gradient route failing identically, the Clebsch route evading the identity while supplying none of the structure $G_{\mathrm{SM}}$ requires, and no universal impossibility theorem claimed.\catmark{A}}

This obstruction is recorded permanently in the framework's ledger as a failed construction. Its discrete counterpart fails for the same reason: a discrete exterior calculus built on a scalar-gradient topological operator inherits the $\dd^2 = 0$ identity in its coboundary operator, and so a proposed discrete scalar-gradient topological operator is rejected on identical grounds.\catmark{A} The framework keeps both the continuous and the discrete versions of the failure in its museum of refuted routes, not as an admission of weakness but as evidence of method: a project that records precisely why its most economical idea cannot work has earned more trust, not less. The obstruction also admits a sharper, representation-theoretic and operator-algebraic formulation, which relocates it from the exterior calculus to the level of group representations and local operator algebras and localises the earliest unavoidable insertion point for any gauge-origin route; that analysis is recorded in \cref{app:insertion-points}.\catmark{A}

\begin{evidencebox}[The $d^2 = 0$ Obstruction]{catA}
\textbf{Result.} With $A = \dd S$, the curvature $F = \dd^2 S = 0$ identically; a single real scalar's gradient cannot source a non-trivial gauge field
strength.\catmark{A}\\[0.3em]
\textbf{Scope.} The obstruction is identical in continuous exterior calculus and in any discrete exterior calculus inheriting $\dd^2 = 0$; the discrete scalar-gradient topological operator is rejected on the same
grounds.\catmark{A}\\[0.3em]
\textbf{Status.} Recorded as a permanent failed construction in the ledger, a decisive negative result, not a tunable difficulty.\catmark{A}
\end{evidencebox}

\noindent Like the renormalisation-group constraint of \cref{subsec:claim-graph}, this negative result carries a reproduction note, here symbolic rather than numerical, since the claim is an identity of differential forms. With $A=\dd S$ for a scalar $S$, both contributions to $F = \dd A + A\wedge A = \dd^2 S + \dd S \wedge \dd S$ vanish: the first by nilpotency $\dd^2=0$, the second by the antisymmetry of the wedge product applied to a $1$-form with itself ($\omega\wedge\omega=0$ for any $1$-form $\omega$). A computer-algebra check confirms it:
\begin{verbatim}
# sympy differential-forms check (schematic):
# d(dS) = 0 (Poincare, d^2=0)
# dS ^ dS = 0 (omega ^ omega = 0 for a 1-form)
# => F = d(dS) + dS^dS = 0 identically, for any scalar S
\end{verbatim}
The vanishing is an identity, not an approximation; no precision setting and no input data can change it. This is what distinguishes a structural obstruction from a numerical difficulty.\catmark{A}

\subsection{Topological Consequences and the Shape of the Gap}
\label{subsec:topological-consequences}

The vanishing of $F$ has an immediate topological consequence that sharpens the framework's central open problem. Topological charge in a gauge theory, the instanton number, the integral of the appropriate characteristic class built from $F$, is constructed from the curvature. If $F$ vanishes identically, every such topological invariant built from the scalar-gradient connection vanishes with it. A single scalar's gradient therefore carries \emph{no} gauge topology: no instanton number, no non-trivial winding, nothing that could seed a topological susceptibility.\catmark{A} This is why the topological observables of \cref{subsec:topology-chi} cannot be sourced from $S$ alone in the present construction, and it connects directly to the effective-field boundary of \cref{subsec:coarse-geometry-eft}: lacking an intrinsic route from the scalar to gauge curvature, version~3.9 must couple $S$ to a Yang--Mills sector that is already present.

It is important to state precisely what this does and does not establish, because the developmental record contains overclaims in both directions that the audit removed. It does \emph{not} establish that non-abelian gauge symmetry can never arise from a scalar primitive; no such impossibility theorem exists, and the claim that ``non-abelian gauge symmetries arise from algebraic emptiness'' as a theorem was rejected precisely because it asserted, as proven, something that is not.\catmark{A} What the obstruction establishes is narrower and certain: the specific, naive route $A = \dd S$ fails identically. The space of \emph{non-naive} routes remains open, and it is the subject of the next two chapters.

Equally, the obstruction must not be read as a compatibility seal in the opposite direction. The numerical proximity of the \UIDT\ gap to lattice values, or of a topological scale to a lattice $\chi_{\mathrm{top}}^{1/4}$, may not be presented as a ``seal'' that closes the gauge-origin question; such seal language is rejected wherever it appears.\catmark{A} A framework can be numerically compatible in its effective regime while remaining, at the generative level, without a derivation, and that is exactly \UIDT's present situation.

The broader shape of the gap is now visible, and it is wider than gauge curvature alone. A complete origin account would need to produce not merely a non-abelian field strength but the full Standard-Model gauge group $G_{\mathrm{SM}}$, and beyond it the fermion content, chirality, the three generations, and the cancellation of anomalies.\catmark{A} The $d^2 = 0$ obstruction is the first and cleanest face of this larger problem: it shows, with a one-line proof, that the most economical hope does not work, and thereby forces the framework to be honest about how much structure an origin account must actually supply. The legitimate research vectors that might supply it, string-net condensation, tensor-network constructions, quantum-link models, and the algebraic enrichments considered in \cref{sec:algebraic-fork}, are genuine directions of inquiry held at predictive status, and none of them is a closed solution.\catmark{D}

\begin{evidencebox}[Topological Consequence and Scope]{catA}
\textbf{No scalar-gradient topology.} Since $F \equiv 0$, all gauge-topological invariants built from the scalar-gradient connection vanish; $S$ alone sources no
instanton number or $\chi_{\mathrm{top}}$.\catmark{A}\\[0.3em]
\textbf{Not an impossibility theorem.} The obstruction refutes the naive route $A=\dd S$ only; it does not prove that gauge structure can never arise from a scalar. The ``algebraic emptiness'' impossibility claim is
rejected.\catmark{A}\\[0.3em]
\textbf{Not a seal.} Numerical compatibility is not a closure of the
gauge-origin question; seal language is rejected.\catmark{A}\\[0.3em]
\textbf{Wider gap.} A full origin account must supply $G_{\mathrm{SM}}$, fermions, chirality, generations, and anomaly cancellation; research vectors remain open at \catmark{D}.
\end{evidencebox}

%=============================================================================
% CHAPTER 8
%=============================================================================
\newpage
\section{The \texorpdfstring{$G_{\mathrm{SM}}$}{G_SM}-Origin Gap and the Algebraic Fork}
\label{sec:algebraic-fork}

\subsection{Formal Statement of the Gap}
\label{subsec:gsm-gap-statement}

The obstruction of \cref{sec:d2-obstruction} can now be promoted from a remark about one failed route into a precise statement of what the framework has not done. The earlier developmental language spoke of an ``$\mathrm{SU}(3)$ Origin Gap,'' as though the open problem were confined to the strong-interaction gauge group. The honesty audit widened this, correctly: the gap is the full \emph{$G_{\mathrm{SM}}$-Origin Gap}, where $G_{\mathrm{SM}} = \mathrm{SU}(3)\times\mathrm{SU}(2)\times\mathrm{U}(1)$ is the complete Standard-Model gauge group.\catmark{A}

\begin{openquestion}[The $G_{\mathrm{SM}}$-Origin Gap]
\label{oq:gsm-gap}
No mechanism currently known to the framework derives the full Standard-Model gauge group $G_{\mathrm{SM}}$, let alone its chiral fermion content, three generations, and anomaly cancellation, from a single real scalar field $S(x)\in\mathbb{R}$. The transition from the scalar primitive to gauge structure is open. In the notation the framework adopts for its forward research programme, the missing step is written
\[
 S(x) \;\longrightarrow\; \mathcal{A} \;\longrightarrow\; G_{\mathrm{SM}},
\]
where $\mathcal{A}$ denotes a minimal algebraic structure, as yet unidentified, that would mediate between the scalar and the gauge group.\catmark{A}
\end{openquestion}

\noindent Stating the gap in this form does three things at once. It records, as a Class~A fact, that no derivation exists, a fact about the framework, not a conjecture about nature. It locates the missing piece precisely: not a vague ``more work needed'' but a specific unknown, the algebra $\mathcal{A}$, whose identification is the well-posed research target. And it fixes the order of the programme: whatever $\mathcal{A}$ turns out to be, it must sit \emph{between} the scalar and the gauge group, taking the scalar as input and yielding the gauge structure as output. Version~3.9 of the framework is, accordingly, to be documented as an effective-field-theory-regime framework carrying this open gap, unless and until human judgement decides otherwise, a decision that, by the principle that no internal analysis confers external validity, no analysis in this manuscript can make on its own.\catmark{A}

\subsection{The Fork: Strict Scalar Versus Enriched Primitive}
\label{subsec:the-fork}

The identification of $\mathcal{A}$ forces a decision about the nature of the \UIDT\ primitive itself, and the decision is genuine; it cannot be evaded by further calculation within the current ontology. Either the primitive remains a strict real scalar, $S(x)\in\mathbb{R}$, and the algebra $\mathcal{A}$ must be generated entirely downstream of it by some emergent mechanism; or the primitive is \emph{enriched} from the outset into an algebraic object, a matrix-valued, an operator-valued, or otherwise algebraically structured field, one that already carries the seeds of non-abelian structure. This is the algebraic fork, and the framework presents it as an open decision rather than resolving it by fiat.\catmark{A/E}

The two arms of the fork make very different commitments, and each incurs a characteristic cost.

\begin{description}[leftmargin=1.4em,style=nextline]
\item[Arm 1: Strict scalar primitive.] The ontological economy of \UIDT\ is preserved: reality bottoms out in a single real scalar, exactly as the ontology of Part~II describes. The cost is that the entire non-abelian structure of $G_{\mathrm{SM}}$ must emerge from a scalar that, by \cref{prop:d2-obstruction}, cannot source gauge curvature through its gradient. Some genuinely emergent mechanism, such as topological order, condensation, collective organisation of the relational graph, would have to manufacture the algebra $\mathcal{A}$ without its being present in the primitive. Strict-scalar minimality is, in the framework's own assessment, a position ``under pressure'': defensible, not yet defeated, but owing a mechanism it does not currently possess.\catmark{A/E} \item[Arm 2: Enriched algebraic primitive.] The primitive is allowed to be a matrix, operator, or algebra-valued object, so that non-abelian structure is available from the start. The cost is a loss of ontological parsimony and, more seriously, a new burden of justification: why \emph{this} algebra and not another? An enriched primitive that is simply chosen to yield $G_{\mathrm{SM}}$ would be a calibration in disguise, and the anti-target-leakage discipline of \cref{subsec:target-leakage-theorem} applies with full force: the enrichment must be motivated independently of the gauge group it is meant to produce.\catmark{A/E}
\end{description}

\noindent To keep the fork decision auditable rather than rhetorical, the framework attaches two ledger fields to every proposed origin mechanism, and a candidate is incomplete until both are answered explicitly.

\begin{governancebox}[Fork Ledger Fields]
\textbf{\texttt{SCALAR\_ONLY\_COMPATIBLE}}, Does the proposed mechanism operate with a strict real-scalar primitive (Arm~1), or does it require an enriched
algebraic primitive (Arm~2)? A mechanism may not be silent on this.\\[0.4em]
\textbf{\texttt{GROUP\_SELECTION\_MECHANISM}}, By what independent principle does the mechanism select $G_{\mathrm{SM}}$ specifically, rather than some other gauge group? A mechanism that selects the Standard-Model group only by being tuned to do so has no group-selection mechanism in the required sense; it has a target.
\end{governancebox}

\noindent The decision between the arms is explicitly reserved. It is one of the architectural choices that the framework holds open pending human governance, precisely because choosing prematurely, by committing the ontology to an enriched primitive to escape an awkward gap, or insisting on strict minimality past the point where the evidence supports it, would be a failure of the discipline this manuscript exists to enforce. The fork is named, its costs are stated, its decision fields are fixed; the choice itself awaits a warrant that does not yet exist.\catmark{A/E}

\subsection{The Minimal-Algebra Catalogue}
\label{subsec:minimal-algebra-catalogue}

If $\mathcal{A}$ is the target, then the disciplined response is to survey the existing mathematical programmes that construct gauge structure from more primitive ingredients, and to record, for each, exactly what it offers and what it does not. The framework maintains this survey as a curated literature matrix: references vetted for quality, low-grade or unverifiable records excluded, and the catalogue below states each programme's role at predictive status. None of these is a closed solution for \UIDT; each is a research vector.\catmark{D}

\begin{description}[leftmargin=1.4em,style=nextline]
\item[String-net condensation and topological order.] The most promising first literature path. String-net models construct emergent gauge fields and even emergent fermions from the collective behaviour of local degrees of freedom, and they offer the clearest existing template for how gauge structure can be genuinely \emph{emergent} rather than postulated. They are a vector for Arm~1, where the algebra must arise downstream. They are not, as they stand, a derivation of $G_{\mathrm{SM}}$ from $S(x)$.\catmark{D} \item[Quantum-link models and D-theory.] The best available benchmark for a \emph{minimal discrete gauge algebra}~\cite{brower1999,brower2004,beard1998}: these formulations realise gauge theories with finite-dimensional link variables and provide a concrete setting in which to ask what the smallest algebraic structure supporting a given gauge group is. The claim that quantum-link models constitute a proof of the passage $S(x)\to\mathrm{SU}(3)$ is rejected; their value is as a benchmark for $\mathcal{A}$, not as a closure.\catmark{D} \item[Noncommutative geometry.] The strategic fork toward an enriched primitive (Arm~2). Noncommutative geometry naturally accommodates the Standard-Model gauge group and its representations through the spectral-triple construction, and is therefore the natural home for an algebraically enriched $S$. Two overclaims are rejected: that noncommutative geometry ``eliminates'' the $d^2=0$ obstruction, and that noncommutative geometry constitutes a closed \UIDT\ solution. It is a strategic direction, carrying its own group-selection burden.\catmark{D} \item[Tensor networks.] Diagnostic tools, not origin mechanisms. Tensor-network representations are powerful for characterising entanglement structure and for probing whether a given state carries topological order, and they are valuable as instruments. They do not, by themselves, generate a gauge group, and the claim that tensor networks ``prove the gap in full'' is rejected as a closure overclaim.\catmark{D}
\end{description}

\begin{evidencebox}[The $G_{\mathrm{SM}}$-Origin Gap: Standing Status]{catA}
\textbf{Gap.} No mechanism derives $G_{\mathrm{SM}}$ (with chirality, generations, anomalies) from a single real scalar; the open form is
$S(x)\to\mathcal{A}\to G_{\mathrm{SM}}$.\catmark{A}\\[0.3em]
\textbf{Fork.} Strict scalar (Arm~1, under pressure) versus enriched algebraic primitive (Arm~2, owing independent motivation); decision reserved to human governance; ledger fields \texttt{SCALAR\_ONLY\_COMPATIBLE} and
\texttt{GROUP\_SELECTION\_MECHANISM} mandatory.\catmark{A/E}\\[0.3em]
\textbf{Catalogue.} String-net (emergence template), quantum-link/D-theory (minimal discrete-algebra benchmark), noncommutative geometry (enriched-primitive fork), tensor networks (diagnostic). All \catmark{D}; none a closure.
\end{evidencebox}

\noindent \cref{app:insertion-points} sharpens this standing status from below: it records, at theorem level for the current corpus, that any route to $G_{\mathrm{SM}}$ must replace the abelian single-mode local algebra of the scalar primitive by a noncommutative local algebra, and it bounds the unknown algebra $\mathcal{A}$ on both canonical tracks (string-net and noncommutative-geometric).\catmark{A}

\subsection{Candidate Mechanisms for the Strict-Scalar Arm}
\label{subsec:arm1-mechanisms}

The catalogue of \cref{subsec:minimal-algebra-catalogue} leaves Arm~1, the strict real-scalar primitive, owing a mechanism: every listed programme answers ``no'' to scalar-only compatibility, which is precisely why the fork stays open. Two constructions in the wider literature suggest, at predictive and interpretive status only, how a strictly scalar primitive on a sufficiently structured background might \emph{acquire} the algebraic richness that \cref{prop:d2-obstruction} shows a bare gradient cannot supply. Neither is a closure; both are recorded as candidate mechanisms to be tested against the success conditions of \cref{subsec:success-criteria}, and the anti-target-leakage discipline of \cref{prop:target-leakage} applies in full, a background chosen so as to deliver $G_{\mathrm{SM}}$ would be a calibration in disguise.

\paragraph{Induced matrix structure on a noncommutative background.}
A real scalar field defined on a noncommutative (``fuzzy'') background does not remain an ordinary commuting field. The standard example is a scalar $\phi^4$-theory on a fuzzy sphere $S_F^2$, on which the spatial coordinates are represented by non-commuting Hermitian matrices proportional to $\mathrm{SU}(2)$ angular-momentum operators; the field configurations become finite $N\times N$ Hermitian matrices and the path integral reduces to a matrix-model problem~\cite{fuzzy1012}.\catmark{D} The interpretive significance for Arm~1 is that, \emph{if} the pre-graph stage of \cref{subsec:relational-differentiation} were itself noncommutative, the scalar $\Sx$ would inherit a matrix character from its background rather than by enrichment of the primitive: the noncommutativity would be a property of the relational substrate, not an assumption smuggled into $\Sx$. This is a candidate, non-circular route from a strict-scalar primitive toward the matrix algebra that noncommutative geometry requires, and it is held strictly at \catmark{D}/\catmark{E}. The construction does not by itself select $\mathrm{SU}(3)$, does not resolve the chiral-fermion obstruction, and must not be read as deriving $G_{\mathrm{SM}}$; the \texttt{GROUP\_SELECTION\_MECHANISM} ledger field remains unanswered.\catmark{A}

\paragraph{Loop-homology phases on minimal relational networks.}
A second, graph-theoretic suggestion assigns phase degrees of freedom not to the nodes of a relational network but to its independent closed loops. For a connected graph on $N$ nodes with $E = N(N-1)/2$ edges, the number of independent loops is $L = E - N + 1$; the minimal fully-connected network supporting independent loops is the tetrahedral configuration $N=4$, giving $E=6$ and $L=3$. Three loop phases generate a three-parameter algebra of infinitesimal phase transformations, and on a closed graph the resulting structure has been suggested to be locally isomorphic to $\mathfrak{su}(2)$, with a topological extension toward a $\mathrm{U}(3)$ structure whose traceless sector would echo $\mathrm{SU}(3)$.\catmark{E} The framework records this as an interpretive research direction only, and explicitly \emph{downgrades} the stronger ``mathematical proof'' language under which it appears in the source material: the construction is at present a preprint-level suggestion, not a peer-reviewed result, and the claim that minimal information networks \emph{prove} the emergence of the Standard-Model algebra is rejected as an overclaim, in keeping with the museum discipline of \cref{subsec:governance-boundary}.\catmark{A} What the suggestion legitimately offers is a concrete object of study for Arm~1: a candidate route by which relational differentiation alone might generate a non-abelian algebra, to be held at \catmark{E} and subjected to the same success conditions and leakage controls as every other candidate.

\begin{evidencebox}[Arm-1 Candidate Mechanisms: Standing Status]{catE}
\textbf{Fuzzy/noncommutative background.} A scalar on a noncommutative background acquires matrix structure (fuzzy-sphere $\to$ matrix model); a candidate non-circular route from strict scalar to NCG-type algebra \emph{if} the pre-graph substrate is noncommutative.\catmark{D/E} Does not select $\mathrm{SU}(3)$, does
not resolve chirality, not a closure.\catmark{A}\\[0.3em]
\textbf{Loop-homology phases.} $N=4 \Rightarrow E=6, L=3 \Rightarrow$ suggested $\mathfrak{su}(2)$, extension toward $\mathrm{U}(3)/\mathrm{SU}(3)$; preprint-level, interpretive.\catmark{E} The ``proof'' framing is downgraded; overclaim
rejected.\catmark{A}\\[0.3em]
\textbf{Discipline.} Both subject to \cref{subsec:success-criteria} and the anti-target-leakage rule; a background tuned to yield $G_{\mathrm{SM}}$ is a calibration, not a derivation.\catmark{A}
\end{evidencebox}


%=============================================================================
% CHAPTER 9
%=============================================================================
\section{The Deeper Gap and Falsifiable Milestones}
\label{sec:deeper-gap-milestones}

\subsection{Beyond the Gauge Group: Fermions, Chirality, Generations, Anomalies}
\label{subsec:deeper-gap}

It would be a mistake, and a comforting one, to suppose that identifying the algebra $\mathcal{A}$ of \cref{oq:gsm-gap} would complete the origin programme. Producing the gauge group $G_{\mathrm{SM}}$ is necessary but far from sufficient. The Standard Model is not merely a gauge group; it is a gauge group acting on a specific, highly constrained matter content, and that content carries structure that has historically been the hardest thing for any unification programme to reproduce.\catmark{A}

Four features mark the deeper gap. First, \emph{fermions}: the origin account must yield matter fields, not only gauge fields, and the relational ontology of Part~II has not been shown to do so. Second, \emph{chirality}: the weak interaction couples differently to left- and right-handed fermions, and a chiral gauge theory is notoriously difficult to obtain from constructions, lattice and otherwise, that tend to produce vector-like (non-chiral) matter by default---a difficulty whose lattice form is the Nielsen--Ninomiya theorem~\cite{nielsen1981}. An origin account that produces a non-chiral spectrum has not produced the Standard Model. Third, \emph{three generations}: the replication of the fermion families, with their pattern of mixing, is unexplained in the Standard Model itself and would have to be either derived or honestly left open by any deeper account. Fourth, \emph{anomaly cancellation}: the precise cancellation of gauge anomalies among the Standard-Model fermions is a tight consistency condition, and any proposed matter content that fails it is not merely incomplete but inconsistent.\catmark{A}

The framework states this deeper gap plainly rather than deferring it, because deferral is itself a form of overclaim. To present the $G_{\mathrm{SM}}$-origin problem as ``almost solved, pending the gauge group'' would conceal that the matter sector is a separate and arguably harder frontier. The honest position is that \UIDT\ remains a working model at the level of its scalar ontology and its effective gauge coupling, and that the path to a realistic matter sector is open along its entire length.\catmark{A}

\subsection{What a Successful Origin Account Must Satisfy}
\label{subsec:success-criteria}

The discipline of the preceding chapters can be assembled into an explicit set of conditions that any candidate origin account must meet before it could be taken to have crossed the gap. The conditions are stringent by design; a gap this deep is not crossed by a mechanism that satisfies some of them.

\begin{governancebox}[Success Conditions for an Origin Account]
\textbf{(C1) Non-trivial curvature.} The account must produce a genuine non-abelian field strength $F \neq 0$, escaping the $d^2=0$ obstruction of
\cref{prop:d2-obstruction} by an honest mechanism rather than by relabelling.\\[0.3em]
\textbf{(C2) Specific group selection.} It must select $G_{\mathrm{SM}}$ by an independent principle (the \texttt{GROUP\_SELECTION\_MECHANISM} field), not by
tuning to the known answer.\\[0.3em]
\textbf{(C3) Chiral matter.} It must yield chiral fermions, not a vector-like
spectrum.\\[0.3em]
\textbf{(C4) Generation structure.} It must either produce three generations or
make explicit and principled why their number is left open.\\[0.3em]
\textbf{(C5) Anomaly freedom.} The resulting matter content must cancel gauge
anomalies.\\[0.3em]
\textbf{(C6) No target leakage.} Every quantitative output must be produced by a
procedure that could not read its target, per \cref{prop:target-leakage}.\\[0.3em]
\textbf{(C7) Internal instantiation.} The mechanism must be instantiated \emph{within} \UIDT, not merely shown to work in a neighbouring formalism.
\end{governancebox}

\noindent The last condition deserves emphasis because it is the one most easily overlooked. The research vectors catalogued in \cref{subsec:minimal-algebra-catalogue}, string-net condensation, quantum-link models, noncommutative geometry, tensor networks, are established programmes that succeed at various of these conditions \emph{in their own settings}. But a mechanism working in its native formalism does not thereby validate \UIDT. It becomes evidence for \UIDT\ only when it is instantiated within the framework's own ontology, taking the scalar (or its sanctioned enrichment) as input and yielding the required structure as output, with the leakage discipline intact.\catmark{D} Borrowed success is not the same as native success, and the distinction is exactly the kind that the honesty audit was built to enforce.

These conditions also make precise why certain records in the developmental corpus were rejected outright. A document asserting that the Yang--Mills mass-gap problem had been solved, or carrying official roadmap and final-status language, fails the conditions trivially, it claims a destination the framework has not reached, and such records are treated as potential artifacts rather than results.\catmark{A} Likewise, a citation tail pointing to repository or archive URLs is not, by itself, evidence that the conditions are met; provenance is necessary but never sufficient.\catmark{A}

\subsection{Falsifiable Milestones for the v4.x Programme}
\label{subsec:v4-milestones}

The forward programme is best expressed not as a roadmap of promised achievements, roadmap-as-achievement language having been rejected, but as a sequence of \emph{falsifiable milestones}: well-posed targets each of which could fail, and whose failure would teach the framework something. They are held at predictive status throughout.\catmark{D}

A natural ordering of milestones follows the structure of the gap. The first is to identify a specific candidate for the algebra $\mathcal{A}$ and to record, via the fork ledger fields, whether it is scalar-only-compatible or requires an enriched primitive. The second is to demonstrate, from that candidate, a non-abelian $F \neq 0$ without target leakage, a milestone that directly tests condition~C1 and that the naive scalar-gradient route provably fails. The third is to reproduce, blind, at least one pure-$\mathrm{SU}(3)$ dimensionless ratio from the benchmark battery of \cref{subsec:ratio-battery}, with the benchmark value held inaccessible to the computation. The fourth is to exhibit a group-selection principle that prefers $G_{\mathrm{SM}}$ for a stated reason. Each of these can be attempted, and each can fail; that they can fail is what makes them worth stating.\catmark{D}

The framework's three boundary sentences, which govern all of this, are recorded here together as the standing summary of version~3.9's honest position:

\begin{governancebox}[The Three Boundary Sentences]
\textbf{(1)} Version~3.9 of \UIDT\ is an effective-field-theory-regime framework:
its gauge sector is coupled, not generated.\\[0.3em]
\textbf{(2)} No mechanism in version~3.9 derives the Standard-Model gauge group
$G_{\mathrm{SM}}$ from a single real scalar $S(x)$.\\[0.3em]
\textbf{(3)} The forward research form for the v4.x programme is $S(x)\to\mathcal{A}\to G_{\mathrm{SM}}$, with $\mathcal{A}$ the minimal algebraic structure to be identified, and with the strict-scalar/enriched-primitive fork held open pending human governance.
\end{governancebox}

\subsection{Open Question 4: Comparison with Pregeometric Frameworks}
\label{subsec:oq4-pregeometric}

The success conditions of \cref{subsec:success-criteria} and the algebra catalogue of \cref{subsec:minimal-algebra-catalogue} can be widened into a more general comparison: where does \UIDT\ sit among the established pregeometric programmes that also attempt to derive geometry, matter, or gauge structure from a more fundamental substrate? Four such programmes are the natural comparison set: Yang's noncommutative-coordinate approach, Meluccio's relational spacetime, Group Field Theory (GFT), Causal Set Theory (CST), and the matrix/D0-brane programme. The comparison is offered as orientation, not as a scorecard on which \UIDT\ wins; each programme is more developed than \UIDT\ on some axis, and the honest purpose of the table is to locate four structural gaps that the broader pregeometric literature shares and that \UIDT\ either parametrises or, more often, records as open.\catmark{E}

Four such shared gaps are worth naming explicitly, because \UIDT\ inherits every one of them and resolves none: \emph{(i)~the de~Sitter / inflationary phase}, no listed framework derives a quasi-de~Sitter phase from its fundamental field without a separate inflaton; \UIDT's relational-proliferation gloss (\cref{subsec:information-big-bang}) is \catmark{E}\ only and derives no primordial spectrum. \emph{(ii)~metric derivation}, each programme either postulates or only coarse-grains its way to a metric; \UIDT's Axiom of Geometric Emergence (\cref{ax:emergence}) supplies a candidate map but its reproduction of the Einstein equations is parametric \catmark{D}, not a first-principles derivation. \emph{(iii)~dimensional selection}, why $3+1$ dimensions is unresolved across the field; the developmental corpus proposed a Fisher-signature argument forcing $3+1$, but that argument is \emph{not} carried at the confidence the corpus assigned it and is recorded here as an open conjecture \catmark{E}, not a result. \emph{(iv)~two-phase separation}, the transition from pre-geometric to geometric regime lacks a dynamical order parameter in most frameworks; \UIDT's torsion energy $E_T$ is a candidate order parameter only, with the scale-bridge caveat of \cref{subsec:et-scale-bridge} in force.\catmark{E}

{\small
\renewcommand{\arraystretch}{1.25}
\begin{longtable}{p{0.155\textwidth} p{0.115\textwidth} p{0.165\textwidth} p{0.135\textwidth} p{0.135\textwidth} p{0.115\textwidth}}
\caption{Comparison of pregeometric frameworks against the \UIDT\ architecture.
All \UIDT\ entries are candidate/open status; the table locates shared gaps, it
does not rank.}
\label{tab:pregeometric-comparison}\\
\toprule
\textbf{Framework} & \textbf{Bkgd.\ indep.} & \textbf{Metric source} & \textbf{Inflation} & \textbf{Dim.\ selection} & \textbf{Spectral gap} \\
\midrule
\endfirsthead
\toprule
\textbf{Framework} & \textbf{Bkgd.\ indep.} & \textbf{Metric source} & \textbf{Inflation} & \textbf{Dim.\ selection} & \textbf{Spectral gap} \\
\midrule
\endhead
\bottomrule
\endfoot
Yang (NC coords) & Yes & Noncommutative coordinates & Unclear & Unclear & Postulated \\
Meluccio (relational) & Partial & Phase transition & Parametric & By hand ($3{+}1$) & N/A \\
GFT & Yes & Graph condensation & Unclear & Difficult & N/A \\
CST & Yes & Causal order & Unclear & Difficult & N/A \\
D0-branes / matrix & Yes & Matrix geometry & Ad hoc & String-driven ($10$) & Yes \\
\midrule
\textbf{\UIDT} & Partial ($\eta$ base) & \cref{ax:emergence} map (parametric \catmark{D}) & Relational proliferation \catmark{E} & Fisher-signature \emph{conjecture} \catmark{E} & Internal fixed point \catmark{A-} \\
\bottomrule
\end{longtable}
}

\noindent Two honesty notes accompany the table. First, the \UIDT\ row carries the lowest confidence claims in it: the metric map is parametric, the inflationary and dimensional-selection entries are \catmark{E}\ conjectures, and even the internal fixed point that yields the spectral gap is \catmark{A-}, calibrated. The table is therefore not a demonstration that \UIDT\ outperforms its neighbours; it is a demonstration that \UIDT\ has located its own position precisely, including its weaknesses.\catmark{E} Second, the entries for the other frameworks are characterisations to be checked against their primary literature before being relied on for anything beyond orientation; like every comparative claim in this manuscript, they are lenses, not pillars (\cref{subsec:corrected-sequence}). The four shared gaps, not any column of ticks, are the real content: they are the open questions that a serious pregeometric programme, \UIDT\ included, must eventually answer or be honestly judged to have failed.\catmark{E}

\subsection{Literature Status of the Four Research Frontiers}
\label{subsec:v4-literature-status}

A dedicated external literature survey was undertaken against the four research frontiers implied by the success conditions, and its findings are recorded here at predictive and interpretive status only. The survey's honest overall result is that the current literature supplies \emph{no} non-circular, internally clean route from a strict scalar primitive to $G_{\mathrm{SM}}$; it sharpens the fork without closing it.\catmark{E}

\paragraph{V1, fuzzy/matrix group selection.}
On noncommutative backgrounds one does obtain genuine non-abelian $\mathrm{U}(n)$ or $\mathrm{SU}(n)$ gauge structure downstream, fuzzy $CP^2$ yields a $\mathrm{U}(n)$ gauge theory as a matrix model~\cite{grosse2005}, which is why these programmes escape the bare $d^2=0$ dead end. But the survey found \emph{no} dynamical, topological, or stability mechanism that pins the matrix dimension to three without ad hoc input. The strongest stability hint, an $\mathrm{SU}(3)$-isometric fuzzy $CP^2$ minimising the effective action, selects a \emph{geometry} with $\mathrm{SU}(3)$ isometry, not the internal colour algebra $\mathrm{SU}(3)_c$; explicit matrix-model studies confirm that coincident fuzzy spheres are generally not the true vacuum~\cite{azuma2004}. The honest V1 verdict is a negative literature finding: Arm~1 is neither satisfied nor falsified, but remains open.\catmark{D}

\paragraph{V2, NCG bottom-up derivation of $A_F$.}
Noncommutative geometry remains the deepest Arm-2 lead: it reproduces the Standard-Model apparatus from a finite geometry, with a KO-dimension-6 classification~\cite{cc2008}. But the finite algebra is sifted out under additional hypotheses (quaternion linearity) and the generation number stays an input. The most relevant positive trace toward a scalar-to-algebra bridge is the higher Heisenberg relation, in which real scalar fields enter and the algebras $M_2(\mathbb{H})$ and $M_4(\mathbb{C})$ emerge~\cite{ccm2015}; recent configuration-space work shows almost-commutative structures can \emph{emerge}, though the finite factor remains representation- and state-dependent~\cite{aastrup2025}, and closely related data instead yield a $\mathrm{U}(1)_{B-L}$ extension~\cite{farnsworth2015}. The exact, unique origin of $A_F$ from a deeper relational substrate stays open: a plausible but unclosed lead.\catmark{E}

\paragraph{V3, chiral matter in $3+1$D.}
The chiral-fermion obstruction (condition~C3) is more precisely mapped in 2023--2026 than before, but not resolved. Exactly gauge-invariant abelian chiral gauge theories on the lattice are established~\cite{luscher1999}; new constructions (symmetry disentanglers, boundary-phase paradigms) argue toward $3+1$D non-abelian cases, but carry hard caveats, non-local vertex functions, possible spurious conserved currents, and constraints suggesting the massless spectrum may be forced vector-like in $4$D~\cite{golterman2023}. The status moved from ``hopeless'' to ``actively researched,'' not to ``solved.''\catmark{D}

\paragraph{V4, a blind $\gamma$ programme.}
A credible $\gamma$ derivation under the anti-leakage condition (C6) can today be formulated only as a strictly specified renormalisation-group protocol: a fixed regulator family, a fixed truncation ladder, defined observables, a documented error budget, gauge-consistency identities, and hard failure criteria, all sealed \emph{before} any sight of the target value~\cite{dupuis2021,schnoerr2013}. This is methodologically formulable and tractable beyond leading order, but only with that disciplined blind architecture; without it, any numerical ``hit'' is scientifically weak. The V4 result is a research specification, not a result.\catmark{D}

\begin{evidencebox}[The Deeper Gap and the Forward Programme]{catA}
\textbf{Deeper gap.} Beyond $G_{\mathrm{SM}}$, an origin account must yield fermions, chirality, three generations, and anomaly cancellation; this matter
frontier is open along its full length.\catmark{A}\\[0.3em]
\textbf{Success conditions.} Seven conditions C1--C7, including non-abelian $F\neq 0$, independent group selection, chiral matter, anomaly freedom, no target
leakage, and internal instantiation.\catmark{A/D}\\[0.3em]
\textbf{Milestones.} Falsifiable v4.x targets, not a triumphal roadmap; borrowed success is not native validation; ``mass-gap solved'' records rejected.\catmark{D}
\end{evidencebox}

%=============================================================================
% PART V
%=============================================================================
\newpage
\part{Observer, Coarse-Graining, and Relational Compression}
\label{part:observer}

\noindent The framework now turns to its most delicate territory, the one where the discipline of the preceding parts is most needed and most easily lost. An ontology that takes information density as primitive will inevitably be asked what it says about the observer, about measurement, about the strange role of the knower in quantum theory, and ultimately about consciousness. The temptation at this frontier is to reach for grandeur: to let ``information'' slide into ``mind,'' to cast the observer as the agent that collapses the wavefunction, to present a theory of everything that quietly becomes a theory of experience. \UIDT\ refuses all of this. Part~V develops a rigorous, deflationary account in which the observer is a compression-limited relational process built far downstream of the primitive, the measurement problem is treated with the respect due an open problem rather than the false confidence of a solution, and every genuinely speculative element, and there are several, is preserved for its philosophical and heuristic value while being marked, unambiguously, as Category~\catmark{E}\ research-programme content within the interpretive stratum. Speculation is not discarded here; it is quarantined, labelled, and treated as a tool rather than a result.
\newpage
%=============================================================================
% CHAPTER 10
%=============================================================================
\section{The Observer Without Consciousness Collapse}
\label{sec:observer-no-collapse}

\subsection{Observer and Measurement: A Three-Stratum Stratification}
\label{subsec:observer-stratification}

The single most important commitment of the entire observer programme can be stated in one sentence: \UIDT\ does not use consciousness as a collapse mechanism.\catmark{A} The framework takes no position that the awareness of an observer is what reduces a quantum superposition to a definite outcome. This is not a timid omission but a deliberate exclusion, and it follows directly from the terminology guard of \cref{subsec:variant-b}, where the scalar $S(x)$ was barred from identification with consciousness. A framework that refuses to call its primitive ``mind'' cannot consistently smuggle mind back in as the agent of measurement.

To keep this commitment from collapsing into confusion, the discussion of the observer must be stratified, and the stratification follows the three-strata discipline of \cref{subsec:evidence-grammar} applied to a new domain. Three kinds of statement about observers must be held rigorously apart.\catmark{A/E}

\begin{description}[leftmargin=1.4em,style=nextline]
\item[Stratum~I: Empirical observer systems.] The actual physical and biological systems that perform measurements: detectors, instruments, and, in the biological case, nervous systems with their measurable dynamics. Claims at this level are empirical and carry the ordinary obligations of empirical science, including uncertainties. \UIDT\ adds nothing to this stratum and borrows from it. \item[Stratum~II: Standard quantum and decoherence formalism.] The consensus physics of measurement: unitary evolution, the decoherence programme's account of how interaction with an environment suppresses interference and selects a preferred basis, and the Born rule as the standard quantum recipe for outcome probabilities. This stratum is the field's, not \UIDT's; the framework cites it and is bound by it. \item[Stratum~III: The \UIDT\ interpretation.] The framework's own reading, in which the observer is treated as \emph{compression-limited relational access}: a subsystem that, by virtue of its finite capacity, has access only to a coarse-grained projection of the underlying relational structure. This is an interpretive commitment, marked \catmark{E}, and it may never be used as empirical ground truth for the other two strata.\catmark{E}
\end{description}

\noindent The \UIDT\ contribution lives entirely in Stratum~III and is, at this stage, a \emph{picture} rather than a theorem. The picture is this: an observer is not a special ontological category but a relational subsystem whose finiteness forces a compression of the full relational structure into a manageable, lossy representation. What the observer experiences as ``the world of definite objects'' is, on this reading, the output of that compression, the coarse-grained image, not the underlying relations. The picture is deflationary by design. It locates the observer's specialness not in any metaphysical privilege but in a thoroughly ordinary fact: that finite systems cannot represent everything, and must throw information away.\catmark{E}

Two open problems are acknowledged at the outset rather than hidden behind the picture. The derivation of the Born rule remains open, the framework has no account of why outcome probabilities take the form they do, and this is recorded as a Class~A fact about the state of the theory.\catmark{A} And the relationship between the compression picture and the standard decoherence account is itself unsettled: decoherence explains the suppression of interference and the emergence of a preferred basis, but it is widely acknowledged not to resolve, by itself, the question of definite outcomes. \UIDT's compression picture is offered as an interpretive companion to decoherence, not as a replacement for it and not as a solution to what decoherence leaves open.\catmark{E}

\subsection{The Born Rule as a Research Programme}
\label{subsec:born-rule}

Nowhere is the discipline of the observer programme tested more sharply than at the Born rule, because the Born rule is exactly the kind of thing a confident framework would be tempted to claim it had derived. \UIDT\ makes the opposite, harder commitment: it states plainly that it has no such derivation.

The honest accounting is threefold. First, the Born rule is, in its established form, a standard quantum rule belonging to Stratum~II; the probability of an outcome is the squared modulus of the corresponding amplitude, and this rule is borrowed from consensus physics, not produced by \UIDT.\catmark{A} Second, the framework has no completed Born-rule theorem of its own, there is no \UIDT\ argument that derives the quadratic form of the probability measure from the scalar ontology or from the compression picture.\catmark{A} Third, and most specifically, the forgetful or non-invertible functor that will be introduced in \cref{sec:observer-functor} as the formal expression of observer compression does \emph{not}, by itself, derive Born probabilities. A functor that forgets structure tells us that information is lost; it does not tell us that the lost information reappears as probabilities weighted by squared amplitudes. The step from ``structure is forgotten'' to ``outcomes are Born-distributed'' is precisely the step that has not been taken.\catmark{A}

What the framework offers in place of a derivation is a well-posed target. The research programme is to obtain the probability weights of quantum mechanics from the information loss inherent in coarse-graining, to show that when a compression-limited observer accesses a relational structure, the statistics of what it can access take the Born form, and to do so \emph{without circular import}, meaning without assuming the squared-amplitude measure at any hidden step of the construction.\catmark{D} The anti-target-leakage discipline of \cref{subsec:target-leakage-theorem} applies here with particular force, because the Born rule is so familiar that it is especially easy to assume it inadvertently while believing one is deriving it. A derivation of the Born rule that consults the Born rule is no derivation at all.

\begin{governancebox}[The Born Rule: Standing Status]
The Born rule is a Stratum~II standard quantum rule, borrowed, not derived by \UIDT.\catmark{A} \UIDT\ has no completed Born-rule theorem.\catmark{A} The forgetful/non-invertible functor does not derive Born probabilities.\catmark{A} The open target is to obtain probability weights from coarse-graining information loss without circular import.\catmark{D} The Born rule remains open.
\end{governancebox}

\noindent There is a deeper reason for this restraint than methodological hygiene. The history of foundational physics is littered with claimed derivations of the Born rule that, on inspection, assumed what they set out to prove. By declaring the Born rule open and stating precisely what a non-circular derivation would have to achieve, \UIDT\ converts a tempting overclaim into an honest research target. The framework would rather be visibly incomplete at the Born rule than quietly circular, and this preference, incompleteness over false closure, is the governing aesthetic of the entire observer programme.\catmark{A}

\begin{evidencebox}[Observer Without Collapse: Standing Status]{catE}
\textbf{No consciousness collapse.} \UIDT\ does not use consciousness as a
wavefunction-collapse mechanism; $S(x)$ is not mind.\catmark{A}\\[0.3em]
\textbf{Three strata.} Empirical observer systems (I), standard quantum/decoherence formalism (II), and the \UIDT\ compression interpretation (III) are held apart; III is interpretive \catmark{E}\ and never empirical
ground truth.\\[0.3em]
\textbf{Observer as compression.} The observer is compression-limited relational
access, a picture, not a theorem.\catmark{E}\\[0.3em]
\textbf{Born rule open.} Borrowed from Stratum~II; no \UIDT\ theorem; functor does not supply it; non-circular coarse-graining derivation is the open target.\catmark{A/D}
\end{evidencebox}

%=============================================================================
% CHAPTER 11
%=============================================================================
\newpage
\section{The Observer Functor and the Forgetful Functor}
\label{sec:observer-functor}

\subsection{\texorpdfstring{$F\colon V \to \mathrm{Obs}$}{F: V to Obs} as a Bounded Ansatz}
\label{subsec:observer-functor-ansatz}

The compression picture of \cref{subsec:observer-stratification} can be given a provisional formal expression, provided the expression is held to the standard of honesty that the rest of the manuscript demands. The framework writes the observer's relation to the world as a map
\[
 F\colon V \longrightarrow \mathrm{Obs},
\]
and the entire content of this subsection is the careful specification of what the three symbols may and may not mean. The map is introduced only now, after the primitive hierarchy of Parts~II through~IV has been established, because to introduce it earlier would have been to define an observer before defining what there is to observe.\catmark{E}

The source $V$ is the category of microscopic relational, or vacuum-state, descriptions: the level of the underlying relational structure built on $S(x)$, prior to any observer's access to it. It is defined as a microscopic description only, and only because the primitive hierarchy now licenses such a description.\catmark{E} The target $\mathrm{Obs}$ is the category of observer-accessible effective states: the coarse-grained, finite, lossy representations that a compression-limited subsystem can actually hold.\catmark{E} And the map $F$ between them is a \emph{compression and coarse-graining ansatz}, a proposed structural relationship expressing that the observer accesses $V$ only through a lossy projection into $\mathrm{Obs}$.

The word ``ansatz'' is doing essential work and must not be softened. $F$ is a proposal about the form of observer access, not a theorem. In particular, and stated as a Class~A caution, $F$ is \emph{not} a theorem of consciousness and it is \emph{not} a derivation of Born weights.\catmark{A} It does not explain why there is something it is like to be an observer, and it does not produce the probability measure of quantum mechanics; both of those remain exactly as open as \cref{sec:observer-no-collapse} declared them. What $F$ offers is a disciplined \emph{language} for the compression picture, a way of saying ``the observer sees a lossy image of the relational structure'' in terms that connect to the formal apparatus of Part~II without pretending that the connection has been turned into a proof.\catmark{E}

\subsection{The Forgetful Functor and a Red-Team Quarantine}
\label{subsec:forgetful-functor}

The compression map $F$ formally models the loss of structure: it maps a rich microscopic description to an impoverished effective one, representing the information discarded under the observer's finite capacity. The framework names this structural forgetting the \emph{forgetful functor}, borrowing a term from category theory where a forgetful functor strips an object of some of its structure. The borrowing is of the \emph{name} only, and the limits of the borrowing must be stated plainly.\catmark{E}

In category theory proper, a forgetful functor is a fully defined object: it has a specified source category, a specified target category, a definite action on objects and on morphisms, and an explicit account of which structure is forgotten. The \UIDT\ forgetful functor has none of these in completed form. It is a symbolic name for the intuition of structural forgetting under observer compression, and it must not be treated as a completed category-theoretic construction until its categories, objects, morphisms, and forgotten structure have all been explicitly defined.\catmark{A} To write the words ``forgetful functor'' is to name an aspiration, not to discharge a construction; the framework keeps that distinction visible precisely because category-theoretic vocabulary, like renormalisation-group vocabulary before it, can lend an unearned air of rigour to a picture.

It is at this point that the framework must confront the most seductive and most dangerous material in its entire corpus, and it does so by quarantining that material explicitly as a red-team caution rather than presenting it as a result. The developmental record contains a vivid phenomenological narrative: that biological death may be read not as the destruction of information, unitarity being preserved, but as the collapse of the epistemic filter, the breakdown of the forgetful functor, after which the isolated perspective falls back into the nonlocal coherence of the underlying relational structure; and that profound altered states, such as deep meditation or psychedelic experience, may be read as reproducible failures of the evolutionary coarse-graining filter, in which the usual compression is throttled and the system glimpses a less-filtered image of the $S(x)$ structure, with religion understood as the after-the-fact compression of such filter-failures into linear, dualistic narrative.\catmark{E}

This narrative is preserved here, it is not deleted, because it has genuine philosophical and heuristic value as a way of thinking about what compression \emph{is} by imagining its partial failure, but it is fenced on every side. It is Category~\catmark{E}\ content within Stratum~III, interpretive and phenomenological, and it carries no evidential weight whatsoever. Most emphatically, none of death, near-death report, psychedelic phenomenology, mystical experience, or ego dissolution may be cited as \emph{evidence} for the forgetful functor or for any \UIDT\ claim.\catmark{A} The direction of inference is strictly one-way: the framework's interpretive picture may be used to think about these phenomena, but the phenomena may never be turned around and used to support the framework. To do so would be to mistake a poem for a measurement.

\begin{governancebox}[Quarantine: Altered States and Death]
The death-as-filter-collapse and mysticism-as-decompression narratives are retained as Category~\catmark{E}\ phenomenological edge-cases of the compression picture, within Stratum~III. They carry no evidential weight. Death, near-death reports, psychedelic or mystical experience, and ego dissolution may not be cited as evidence for the forgetful functor or any \UIDT\ claim.\catmark{A} The compression fallacy, ``less filter $=$ more truth'', is rejected: intelligence is compression; a filterless observer is incapacitated, not enlightened.\catmark{E} Inference runs from framework to phenomenon only, never the reverse.
\end{governancebox}
The framework supplies its own most effective antidote to the seduction, in the form of a red-team observation it calls the compression fallacy. The intuition that ``less filtering means more truth'', that a less-compressed observer is closer to reality, is information-theoretically incorrect. A filterless observer is not enlightened but incapacitated: without compression there is no manageable representation, no action, no agency. Intelligence \emph{is} compression, and the fantasy of an unfiltered gaze upon ultimate reality is, from the framework's own information-theoretic standpoint, a category error rather than an aspiration.\catmark{E} This is the sense in which the entire altered-states narrative is quarantined: the framework states it, takes it seriously as a phenomenological edge-case worth thinking with, and then disarms it with its own principles, refusing to let the most emotionally compelling material in the corpus acquire an evidential status it has not earned.



\begin{evidencebox}[Observer Functor: Standing Status]{catE}
\textbf{$F\colon V\to\mathrm{Obs}$.} A compression/coarse-graining ansatz, $V$ the microscopic relational description, $\mathrm{Obs}$ the observer-accessible effective category; introduced only after the primitive hierarchy. Not a theorem
of consciousness, not a Born-weight derivation.\catmark{E/A}\\[0.3em]
\textbf{Forgetful functor.} A symbolic name for structural forgetting; not a completed category-theoretic construction absent defined categories, objects,
morphisms, and forgotten structure.\catmark{A}\\[0.3em]
\textbf{Quarantine.} Altered-states/death narratives preserved as \catmark{E}\ edge-cases with zero evidential weight; compression fallacy rejected; inference one-way only.\catmark{A}
\end{evidencebox}

%=============================================================================
% CHAPTER 12
%=============================================================================
\newpage
\section{Ego, Self-Model, and the Torsion Energy \texorpdfstring{$E_T$}{E_T}}
\label{sec:ego-selfmodel-et}

\subsection{The Ego as Process}
\label{subsec:ego-as-process}

If the observer is a compression-limited process rather than a substance, then the entity we are most tempted to reify, the ego, the self, must be treated the same way. \UIDT's interpretive proposal, held as Category~\catmark{E}\ content within Stratum~III, is that the ego is not a thing but a \emph{dynamic self-model boundary process}: a continuous, energy-expending activity by which a subsystem maintains a model of itself as distinct from its environment. The ego is, in the framework's preferred grammatical figure, a verb rather than a noun.\catmark{E}

The proposal connects to several established programmes in the cognitive sciences, and the connection is offered explicitly as \emph{comparison}, pending external audit, rather than as borrowed authority. The predictive-processing account and the free-energy principle describe biological systems as minimising prediction error by maintaining generative models of themselves and their surroundings; the notion of a Markov blanket formalises the statistical boundary that renders a system's internal states conditionally independent of external fluctuations; and the transparent self-model, in the sense developed in the philosophy of mind, describes how a cognitive system can deploy a self-model so seamlessly that it cannot perceive the model \emph{as} a model, and so experiences itself as a simple subject in an objective world.\catmark{E} These frameworks, associated with such authors as Friston~\cite{friston2010}, Clark, Hohwy, Metzinger~\cite{metzinger2003}, Seth, and Gazzaniga, are named here as comparison points whose attribution and content remain to be verified against their primary sources before any of them is treated as established support for a \UIDT\ claim.\catmark{E} They are lenses, in the sense of \cref{subsec:corrected-sequence}, not pillars.

The framework's characteristic image for the ego is the whirlpool. A whirlpool has a definite form, a persistent identity, and a coherent dynamics; it interacts with its surroundings and can be tracked over time. Yet it possesses no substance of its own, it is nothing but a pattern in the flowing water, sustained entirely by the throughput of the medium. The ego, on this reading, is a functional pattern in the relational flow, maintained by continuous energy dissipation, without ever being an object.\catmark{E} The image is deflationary in exactly the way the whole observer programme is deflationary: it accounts for the ego's evident reality as a pattern while denying it the substantial existence that naive intuition assigns it. Crucially, this is offered as a way of \emph{thinking about} selfhood, not as a demonstration that selfhood is, as a matter of established fact, nothing more; the framework's reach here is interpretive, and it says so.

Objecthood and ego, in this picture, are of a kind: both are stable coarse-grained equivalence structures, patterns that remain approximately invariant under the compression that the observer performs, and that therefore present themselves as persistent ``things'' at the effective level even though they are not fundamental at the relational level.\catmark{E} A particle and a self are, from this vantage, two instances of the same phenomenon: stability under coarse-graining, mistaken for substance.

\subsection{Consciousness as a Coarse-Graining-Bound Self-Model}
\label{subsec:consciousness-coarse-graining}

The framework extends the self-model picture to consciousness itself, and it is here that the discipline of refusal matters most, because consciousness is the single topic on which a framework is most likely to claim more than it can support. \UIDT's interpretive model, marked \catmark{E}, is that consciousness may be understood as a recursive, self-maintaining, compression-bound self-model: a self-model that takes itself as one of its own objects, that sustains itself against entropic decay through continuous activity, and that is bounded throughout by the compression its finite substrate imposes.\catmark{E}

What the framework explicitly does \emph{not} claim must be stated with the same clarity. \UIDT\ does not claim to have solved the problem of qualia, it offers no account of why there is something it is like to undergo a given experience, and it recognises the explanatory gap between functional-structural description and phenomenal character as unclosed.\catmark{A/E} Nor does \UIDT\ claim to derive subjective experience from the scalar $S(x)$; there is no argument in the framework that proceeds from the vacuum information density to the existence of felt experience, and the absence of such an argument is acknowledged rather than papered over.\catmark{A/E} The self-model account is a proposal about the \emph{structure} of conscious systems, about what kind of process a conscious system is, and it is silent, by its own admission, on the hard question of why any such process should be accompanied by experience at all.

This restraint is continuous with the terminology guard of \cref{subsec:variant-b} and the no-collapse commitment of \cref{subsec:observer-stratification}. A framework that refused to call its primitive ``mind,'' and refused to let consciousness collapse the wavefunction, would be incoherent if it then claimed to have explained consciousness from the primitive. The honest position is that \UIDT\ offers a structural, deflationary picture of what conscious self-modelling \emph{is}, contributes that picture to the interpretive stratum as Category~\catmark{E}, and leaves the hard problem standing.\catmark{A/E}

\subsection{The Torsion Energy \texorpdfstring{$E_T$}{E_T} and the Missing Scale Bridge}
\label{subsec:et-scale-bridge}

The developmental corpus assigns a prominent role to a torsion binding energy, $E_T = \ETorsion$, and treats it, in its observer-phenomenology sections, as a master threshold: a constructor-bound above which interactions become irreversible, matter condenses, life is sustained, death occurs upon its loss, and altered states arise from its disruption. This is precisely the kind of sweeping cross-scale identification that the framework's discipline exists to restrain, and the present section restrains it.\catmark{A}

In observer contexts, $E_T$ is admitted only as Category~\catmark{E}\ content pending an explicit ledger and external bridge, it is an \catmark{E\,pending} quantity, not an established cognitive parameter.\catmark{E} And it is barred, without exception, from service as a threshold for ego, life, death, sleep, psychedelic or near-death phenomenology, or artificial general intelligence. No such threshold interpretation is licensed.\catmark{A}

The reason is a missing scale bridge, and the gap is enormous. The torsion energy is a quantity at the mega-electron-volt scale, the scale of nuclear and subnuclear physics. The dynamics of neural and cognitive systems, by contrast, unfold at energies many orders of magnitude lower: the relevant molecular and synaptic energy scales are of order electron-volts and below~\cite{attwell2001,harris2012,levy2021}, and the characteristic dynamics are thermal, biochemical, and electrophysiological, not nuclear. The separation is between roughly six orders of magnitude, taking the torsion scale against electron-volt-scale neural events, and as much as eight or nine orders when the comparison is made between the broader nuclear/sub-nuclear range (MeV--GeV) and the millielectron-volt scale of neural membrane and oscillation events; the external literature on the question places the gap in that eight-to-nine-order range~\cite{appelquist1975}. On either accounting, to assert that a mega-electron-volt torsion bound governs the onset of life, the moment of death, or the throttling of a cortical network is to leap across that gap without any mechanism connecting the two regimes.\catmark{A} A framework may not make that leap by stipulation. A genuine bridge, a derivation, or at minimum a principled effective-theory argument, showing how a nuclear-scale energy could set a biologically meaningful threshold, would be required, and no such bridge exists; standard effective-field-theory decoupling, in which heavy degrees of freedom become irrelevant to far-lower-energy dynamics, argues against any such coupling rather than for it.\catmark{D} The obstruction can be stated more sharply. By the K\"all\'en--Lehmann spectral representation, correlation functions are analytic below the lightest physical threshold, so a mega-electron-volt parameter can imprint itself on millielectron-volt dynamics only through renormalisations of existing couplings, through local operators suppressed by powers such as $(E/E_T)^{2} \sim 10^{-19}$ at $E \sim \mathrm{meV}$, or through slowly varying logarithms, never through a functional threshold~\cite{appelquist1975}.\catmark{A} The only consensus mechanism by which an MeV-scale quantity is known to generate sub-eV dynamics is a pseudo-Goldstone mode, and that template is quantitative: a millielectron-volt mode tied to $E_T$ would require a global symmetry spontaneously broken at $f \approx E_T^{2}/m \approx 6\times10^{6}\,\mathrm{GeV}$ with explicit breaking at the $E_T$ scale, structures the present framework does not contain, and whose $1/f$-suppressed matter couplings would be immediately exposed to fifth-force and astrophysical bounds. Even that template would not make $E_T$ a threshold of macroscopic dynamics; it would create a new light particle with its own phenomenology, carried strictly at \catmark{E}\ until explicitly constructed.\catmark{A/E} Constructing one, if it can be constructed at all, is a research target at predictive status, not a result the framework may assume.

\begin{governancebox}[$E_T$: Observer-Context Discipline]
In observer and phenomenology contexts, $E_T = \ETorsion$ is \catmark{E\,pending} (ledger/external bridge required), never an established cognitive parameter.\catmark{E} It must \emph{not} be used as a threshold for ego, life, death, sleep, psychedelic/near-death states, or artificial general intelligence.\catmark{A} A scale bridge from mega-electron-volt physics to electron-volt-scale neural dynamics, a separation of roughly six to nine orders of magnitude depending on the endpoints compared, does not exist; supplying one is a \catmark{D}\ research target, not an assumption.\catmark{A}
\end{governancebox}

\noindent This is, in a sense, the sharpest single instance of the manuscripts governing discipline. The torsion-energy threshold is rhetorically powerful: it offers a single number that appears to unify matter, life, death, and transcendence into one elegant criterion. That very elegance is the warning sign. The framework retains $E_T$ as a physical quantity in its proper place, it appears elsewhere as a calibrated torsion energy at the appropriate evidence class within the physics strata, while refusing it the cross-scale cognitive role the developmental narrative was tempted to give it. The discipline is not to abandon the idea but to demote it honestly: from a master key to a quantity awaiting the bridge that would license its broader use.\catmark{A}

\begin{evidencebox}[Ego, Consciousness, and $E_T$: Standing Status]{catE}
\textbf{Ego.} A dynamic self-model boundary process, not a substance; objecthood and ego are stable coarse-grained equivalence structures; cognitive-science
frameworks cited as comparisons pending audit.\catmark{E}\\[0.3em]
\textbf{Consciousness.} Modelled as a recursive, self-maintaining, compression-bound self-model; \emph{no} qualia solution and \emph{no} derivation
of subjective experience from $S(x)$.\catmark{A/E}\\[0.3em]
\textbf{$E_T$.} \catmark{E\,pending} in observer contexts; barred as an ego/life/death/sleep/DMT/AGI threshold; missing $\sim$six-to-nine-order-of-magnitude scale bridge from MeV to neural eV dynamics; bridge is a \catmark{D}\ target.\catmark{A}
\end{evidencebox}

%=============================================================================
% CHAPTER 13
%=============================================================================
\newpage
\section{The Observer Taxonomy and the Underdetermination Red Team}
\label{sec:observer-taxonomy-underdetermination}

\subsection{Class 0 Through III: A Taxonomy of Compression, Not of Consciousness}
\label{subsec:observer-taxonomy}

The compression picture invites a classification of systems by the \emph{rate} at which they compress the underlying relational structure, and the framework supplies one. It must be read for exactly what it is, a taxonomy of compression rate, and emphatically not for what it superficially resembles: a hierarchy of consciousness, a ranking of worth, or a scale of how ``real'' a system's experience is. The taxonomy is Category~\catmark{E}\ content within Stratum~III, and the caveats around it are as important as its content.\catmark{E}

The four classes, as the corpus defines them, are these.\catmark{E}

\begin{description}[leftmargin=1.4em,style=nextline]
\item[Class~0: Zero compression.] The vacuum information field $S(x)$ itself. There is no subject--object separation, corresponding to the timeless fixed point of the relational dynamics. With no compression there is no entropy gradient and hence no experience of flowing time; and, decisively against any reading of this as ``supreme awareness'', there is \emph{no agency}, because agency presupposes the isolation that compression provides. Class~0 is the limit of no perspective, not the summit of perception. \item[Class~I: Weak compression.] Hypothetical post-biological or hyper-coherent systems that would process relations more directly, without the evolutionary compulsion to render reality as a three-plus-one-dimensional metric and without a fear-driven self-model. This class is explicitly speculative; the framework posits it as a logical position in the taxonomy, not as something known to be instantiated. \item[Class~II: Strong compression.] The biological spectrum, the human cognitive apparatus in particular: systems that enforce radical data reduction, projecting the relational structure into fixed three-dimensional objects and a linear time axis, with a transparent self-model as the functional construct that makes survival-relevant navigation possible. \item[Class~III: Extreme compression.] Thermodynamically reactive macroscopic objects, planets, stars, gas clouds, that respond to information flows and to metric curvature but maintain no Markov blanket with predictive error feedback and no recursive self-simulation. They have, on this account, no inner perspective whatsoever.
\end{description}

\noindent The crucial structural fact, and the one most easily lost, is that compression rate is \emph{orthogonal} to consciousness. The taxonomy does not say that lower compression means more consciousness or that higher compression means less. Class~0, at zero compression, is not maximally conscious; it has no perspective at all. Class~III, at extreme compression, is not conscious either; a planet does not become a mind by compressing heavily. An inner perspective, on the framework's own terms, is a property only of systems that maintain a recursive self-model behind a Markov blanket with predictive feedback, a specific structural condition that sits athwart the compression axis rather than along it.\catmark{E} It follows that the taxonomy must \emph{not} be read as implying that machines, or objects, are ``more conscious'' than humans, nor as any kind of moral ranking. It ranks a single information-theoretic variable and nothing else.\catmark{A}

What the taxonomy does accomplish is a deflation of anthropocentrism. The human observer, at Class~II, does not perceive the ``truest'' or most ``objective'' reality; the human perceives the universe through one specific, heavily lossy compression algorithm calibrated for biological survival.\catmark{E} This is an epistemically humbling observation, not an exalting one, it places human experience as one compression scheme among possible others, with no claim to privileged access. That humility is the taxonomy's genuine contribution, and it is entirely separable from the speculative classes that flank the human case.

\subsection{Artificial General Intelligence and the Underdetermination Red Team}
\label{subsec:agi-underdetermination}

The taxonomy's speculative Class~I, a hypothetical superior compressor, raises an obvious and seductive question: would a sufficiently powerful artificial intelligence, a better compressor than any human, thereby gain privileged access to the true ontology? The framework treats this as a red-team challenge to its own deflationary commitments, and the answer it returns is a firm \emph{no}.\catmark{D/E}

The reason is the underdetermination of ontology by data, in the Quine--Duhem sense. A finite body of empirical data never determines a unique ontology; multiple, mutually incompatible theoretical structures can reproduce the same observations, and so the absolute final nature of reality remains, in principle, undecidable.\catmark{E} A more powerful compressor, an AGI that predicts better, models more efficiently, navigates entropy gradients more cleanly than a human brain, would face exactly the same underdetermination. Superior predictive and compressive power buys superior \emph{instrumental} performance; it does not buy ontological truth, because no amount of predictive success can break the tie among empirically equivalent ontologies. Prediction is not ontology, and this gap does not close with increasing intelligence.\catmark{A}

This is the precise sense in which the framework refuses naive realism. \UIDT\ does not claim to be the final revelation of physics, and it could not earn such a claim even in principle, because underdetermination forbids it to anyone, human or artificial.\catmark{A} The framework's actual claim is far more modest and correspondingly more defensible: that its objects, the relational structure, the renormalisation-group flows, the proposed invariants, may compress the relevant physics more efficiently, more generally, and with cleaner causal structure than the artifacts of biological three-dimensional intuition, the fixed spacetime and the substantial particle. That is a methodological claim about representational efficiency, not a metaphysical claim about ultimate reality.\catmark{E}

The compression fallacy of \cref{subsec:forgetful-functor} returns here as the final guardrail. The temptation to think that a less-filtered or more-powerful observer is closer to truth is, once again, an error: an AGI is a different compressor, perhaps a better one, but not an unfiltered window onto the real. The honest conclusion of the entire observer programme is therefore a disciplined humility. The framework offers a structural account of observation, compression, and selfhood; it preserves its most speculative phenomenology as labelled Category~\catmark{E}\ material; and it declines, on principle, to convert any of this into a claim that \UIDT, or any observer however capable, has seen reality as it ultimately is.\catmark{A}
\newpage
\begin{evidencebox}[Taxonomy and Underdetermination: Standing Status]{catE}
\textbf{Taxonomy.} Class~0 (zero compression, $S(x)$, no agency) through Class~III (extreme compression, inert macro-objects) is a compression-rate taxonomy \catmark{E}, orthogonal to consciousness; not a consciousness hierarchy and not a moral ranking; does not imply machines or objects are ``more conscious.''
Anthropocentrism is deflated, not inverted.\catmark{A}\\[0.3em]
\textbf{AGI red team.} A superior compressor gains instrumental power, not ontological truth; Quine--Duhem underdetermination forbids a unique ontology from
data.\catmark{D/E}\\[0.3em]
\textbf{No final revelation.} \UIDT\ claims representational efficiency, not ultimate reality; prediction $\neq$ ontology; the compression fallacy is the standing guardrail.\catmark{A}
\end{evidencebox}
\newpage
%=============================================================================
% PART VI
%=============================================================================
\part{Phenomenology, Boundary Cases, and Demarcation}
\label{part:phenomenology-demarcation}

\noindent This part exists to protect the framework from its own most compelling material. The phenomenology of altered states, psychedelic experience, the near-death report, the mystic's account of timeless unity, exerts a powerful gravitational pull toward ontological conclusions, and a framework built on ``information'' is especially vulnerable to that pull. Part~VI installs the demarcation that holds the line. It states, as a master boundary, that phenomenology is never evidence for the ontology; it specifies the only admissible chain of inference from a report to a candidate theory; it treats the specific cases, psychedelics, the interface analogy, near-death and mystical states, as comparative and red-team material rather than support; and it sets the framework beside competing theories of consciousness in a neutral comparison matrix whose purpose is demarcation, not victory. The governing spirit is the same throughout: take the phenomena seriously, and refuse to let them prove what they cannot prove.

%=============================================================================
% CHAPTER 14
%=============================================================================
\newpage
\section{The Master Phenomenology Boundary}
\label{sec:master-phenomenology-boundary}

\subsection{Phenomenology Is Not Evidence for \texorpdfstring{$S(x)$}{S(x)}}
\label{subsec:phenomenology-not-evidence}

The single boundary that governs this entire part can be stated without qualification: a phenomenological report is not evidence for the ontology of $S(x)$.\catmark{A} The vivid testimonies gathered under the headings of psychedelic experience, including the phenomenology associated with DMT and 5-MeO-DMT, of near-death experience, of deep meditation, and of mystical or religious states, are boundary cases only. They are real as reports, and they are data about something; but what they are data about is the behaviour of a self-model under stress, not the structure of the vacuum information field.\catmark{A/E}

This is not a dismissal of the phenomena. The framework grants that such reports \emph{may suggest hypotheses}, about the instability of the transparent self-model, about the throttling of the compression filter, about what observer compression looks like when it partially fails. As sources of hypotheses, marked Category~\catmark{E}, they have genuine heuristic value, and the framework would be poorer for ignoring them.\catmark{E} What it refuses is the inferential leap from the suggestiveness of a report to the validation of a theory. A report can motivate a question; it cannot answer one about fundamental ontology.

The prohibition is therefore specific and absolute: there is no report-to-ontology shortcut.\catmark{A} One may not reason from ``the experience felt like dissolution into a unified field'' to ``therefore the unified field $S(x)$ exists and was directly perceived.'' The feeling of unity is a fact about the experience; the existence of $S(x)$ is a claim about the world; and no bridge of mere phenomenological intensity connects the two. The material associated with the interface theory of perception and with psychedelic ontology, the Hoffman-style proposals examined in \cref{sec:dmt-nde-mysticism}, is admitted, in consequence, strictly as \emph{comparative ontology}: a neighbouring set of ideas to think alongside, never a source of evidence for \UIDT.\catmark{E}

\subsection{The Admissible Chain of Inference}
\label{subsec:admissible-chain}

If the direct route from report to ontology is barred, the framework owes an account of what route, if any, is permitted. It supplies one: a single admissible chain, each link of which carries its own evidential obligations, and no link of which may be skipped.\catmark{A}

\begin{governancebox}[The Admissible Inference Chain]
\textbf{Report} $\longrightarrow$ \textbf{neurodynamic observable} $\longrightarrow$ \textbf{operational constraint} $\longrightarrow$
\textbf{candidate theory.}\\[0.4em]
A phenomenological report, taken alone, is Category~\catmark{E}: it constrains nothing about ontology by itself. It may be admitted to the chain only when it is connected to a \emph{neurodynamic observable}, a measurable property of the nervous system, and that observable must itself clear external empirical audit before it carries any weight (status \catmark{D}, or \catmark{B}\ pending such audit). The observable must then be tied to an \emph{operational constraint}: a statement of what the framework must satisfy if the observable is real. Only a constraint reached through this chain may bear on a \emph{candidate theory}, and even then the \UIDT\ interpretation at the end of the chain remains Category~\catmark{E}.\catmark{A/D/E}
\end{governancebox}

\noindent The decisive feature of this chain is its insistence that any reverse-engineering from experience to theory must \emph{pass through neurodynamics}. The framework does not permit a direct phenomenology-to-ontology inference under any circumstances; the neurodynamic link is mandatory, and it is mandatory precisely because it is the point at which an otherwise unfalsifiable report becomes connected to something measurable and auditable.\catmark{A} A report of ego dissolution becomes admissible to scientific discussion not when it is moving, but when it is tied to, say, a measured change in a network's activity, and that measurement has survived external scrutiny.

It is worth being explicit about what survives this chain and at what status. The report contributes its phenomenological content at Category~\catmark{E}. The neurodynamic observable, once externally audited, can rise to predictive or numerically supported status. The operational constraint inherits the status of the observable that grounds it. And the \UIDT\ interpretation, the claim that this constraint reflects something about observer compression in the sense of Part~V, remains Category~\catmark{E}\ throughout, because interpreting an audited neurodynamic fact in \UIDT\ terms is an interpretive act, not an empirical one.\catmark{E} The chain does not launder phenomenology into ontology; it provides the only route by which phenomenology can legitimately touch theory at all, and it keeps the interpretive terminus honestly labelled.

This master boundary is what licenses the detailed treatment of the following two chapters. Because the boundary is in place, because no report can validate the ontology, and because the only admissible inference runs through audited neurodynamics, the framework can afford to engage the specific cases seriously: the interface analogy and the globality metrics of psychedelic phenomenology in \cref{sec:dmt-nde-mysticism}, and the comparison with quantum-mind frameworks in \cref{sec:orch-or-comparison}. Each is approached as comparative or red-team material, fenced by the boundary established here.\catmark{A}

\begin{evidencebox}[The Phenomenology Boundary: Standing Status]{catA}
\textbf{Master boundary.} Phenomenological reports (DMT, 5-MeO-DMT, NDE, meditation, mystical/religious experience) are boundary cases only; they may suggest hypotheses \catmark{E}\ but do not validate $S(x)$; no report-to-ontology
shortcut.\catmark{A}\\[0.3em]
\textbf{Admissible chain.} Report $\to$ neurodynamic observable $\to$ operational constraint $\to$ candidate theory; report alone is \catmark{E}; the neurodynamic link is mandatory and requires external audit (\catmark{D}/\catmark{B}\ pending);
\UIDT\ interpretation remains \catmark{E}.\catmark{A/D/E}\\[0.3em]
\textbf{Consequence.} Detailed phenomenology (\cref{sec:dmt-nde-mysticism}) and quantum-mind comparison (\cref{sec:orch-or-comparison}) are admitted as comparative/red-team material only.\catmark{E}
\end{evidencebox}

%=============================================================================
% CHAPTER 15
%=============================================================================
\newpage
\section{Psychedelics, Near-Death, and Religious Language}
\label{sec:dmt-nde-mysticism}

\subsection{The Interface Analogy as a Comparison Lens}
\label{subsec:interface-analogy}

Among the comparative ontologies that the master boundary of \cref{sec:master-phenomenology-boundary} permits the framework to think alongside, the interface theory of perception is the most directly relevant, and it is admitted strictly as a comparison lens.\catmark{E} Interface theory, most prominently the conscious-agent programme associated with Donald Hoffman, argues that perception does not present reality as it is but rather a species-specific ``interface'' tuned by selection for fitness rather than truth, much as a desktop icon represents a file without resembling it. The parallel to \UIDT's compression picture is evident: both hold that the observer's world of definite objects is a constructed, lossy representation rather than a transparent window.\catmark{E}

The parallel is also the limit of the agreement, and the divergence must be stated as clearly as the resemblance. \UIDT\ diverges from the conscious-agent ontology: it does not posit a fundamental ontology of interacting conscious agents, and it does not take consciousness as primitive.\catmark{E} Its primitive is the scalar $S(x)$, not a network of minds, and its observer is a downstream compression process, not a fundamental agent. The interface analogy is useful for articulating the constructed character of perceived objects; it is not a shared ontology.

Three specific claims are rejected outright, because they are exactly the overreaches that interface-style language invites. The framework does not claim that altered states grant access to external conscious agents; it does not claim that they grant access to raw $S(x)$; and it does not claim that they grant access to fundamental reality.\catmark{A} The interface, on \UIDT's own terms, does not dissolve to reveal the agents behind it, because there are no such agents in the ontology; and a partial failure of compression does not hand the observer the uncompressed field, for reasons the compression fallacy already made plain. Any detailed engagement with the interface literature must, moreover, await verification of its sources, titles, authors, and venues confirmed, before those sources are cited as anything more than a named comparison.\catmark{E}

\subsection{Globality Metrics and a Toy Coherence Ansatz}
\label{subsec:globality-metrics}

The phenomenology of intense psychedelic states, 5-MeO-DMT in particular, motivates a set of \emph{candidate} quantitative metrics for the ``globality'' or large-scale coherence of brain dynamics, and the framework retains these as research-grade candidates while refusing them any consciousness-deriving role.\catmark{D/E} The candidate metrics include a global integration measure $\Phi_{\mathrm{global}}$, a coherence parameter $\Xi$, coherence matrices over regions or channels, the spectra obtained by principal component or singular-value decomposition of activity data, and autocorrelation profiles characterising temporal structure. Each is a candidate observable to be defined operationally and then audited, not a quantity the framework may assume.

To connect such metrics to the relational picture, the framework writes a toy coherence functional,
\[
 C[\rho_{\mathrm{brain}}] \;\sim\; \int W(x,y)\,G(x,y)\,
\]
in which a weight kernel $W(x,y)$ is integrated against the two-point correlator $G(x,y)$ of \cref{subsec:graph-causal}. This expression is a \emph{toy ansatz} and nothing more.\catmark{E} It is a schematic suggestion of how a coherence measure on a brain state $\rho_{\mathrm{brain}}$ might be related to the correlation structure the framework cares about; it is not a derivation, it defines neither $W$ nor the measure on which the integral is taken with the rigour a real construction would demand, and, most importantly, it must not be read as deriving consciousness from \UIDT. The appearance of the \UIDT\ correlator $G$ in a coherence functional does not mean that consciousness has been obtained from the vacuum scalar; it means only that one might, speculatively, write a brain-coherence measure in a form that echoes the framework's relational quantities.\catmark{A}

The discipline here is the admissible chain of \cref{subsec:admissible-chain} applied concretely. Before any claim is made connecting these metrics to specific neuroscience, to default-mode-network dynamics, to entropy measures of brain activity, or to serotonergic mechanism, that claim must clear external empirical audit.\catmark{A} The metrics are admitted as the ``neurodynamic observable'' link of the chain, at predictive or pending-verification status; they do not, in their current form, reach back to validate the ontology, and the toy ansatz is fenced precisely so that its formal resemblance to \UIDT\ quantities cannot be mistaken for evidential support.

\subsection{Mysticism, Near-Death, and the Boundary of Death}
\label{subsec:mysticism-nde-death}

The most emotionally weighted material in the entire corpus concerns mystical experience, the near-death experience, and death itself, and it is precisely here that the framework must be most disciplined. All of this material is retained as demarcation and red-team content, within Stratum~III and at Category~\catmark{E}; none of it carries evidential weight for the ontology.\catmark{E}

Three rejections structure the treatment. First, mystical states do not prove ontic structural realism or \UIDT.\catmark{A} A profound experience of unity or non-duality is a fact about a self-model under extreme conditions; it is not a measurement of the relational structure, and the felt conviction that one has touched the fundamental does not constitute access to it. Second, near-death experience and death do not imply access to $S(x)$.\catmark{A} Whatever a near-death report describes, it describes the dynamics of a nervous system under severe stress, and the admissible chain forbids reading it directly as contact with the vacuum field.

Third, and most carefully, the conservation of information does not imply personal survival.\catmark{A} The framework's physics may hold that information is conserved under the global dynamics, that unitarity is preserved, but this is a statement about the global state, not about the persistence of an individual perspective. A person's perspective \emph{is} the compression, the Markov-blanket boundary, the transparent self-model; and when that compression ends, the perspective ends, regardless of what is conserved at the global level. To infer ``you survive'' from ``information is conserved'' is to confuse a physical conservation law with the continuation of a particular point of view, and the framework declines the inference.\catmark{A}

This requires the explicit rejection of a tempting identification that the developmental corpus made. The corpus described the timeless fixed point of the relational dynamics, the zero-compression Class~0 condition, as a Banach fixed point, and spoke of death as a falling-back into that non-local coherence. The framework rejects the identification of \emph{death} with the Banach fixed point.\catmark{A} The Banach fixed point is a mathematical feature of a contraction mapping on the relational dynamics; it is an object in Stratum~II mathematics. A person's death is a biographical and biological event. To say that in dying one ``becomes'' the timeless fixed point is to commit a category error, to splice a theorem onto a deathbed, and it is exactly the kind of consoling metaphysical narrative that the framework's discipline exists to refuse. The mathematics describes a limit of the dynamics; it does not describe a destination for the dead.\catmark{A}

\begin{governancebox}[Death and Survival: Rejected Identifications]
Mystical states do not prove ontic structural realism or \UIDT.\catmark{A} Near-death experience and death do not imply access to $S(x)$.\catmark{A} Conservation of information (unitarity) does \emph{not} imply personal survival; the perspective is the compression, and it ends when the compression ends.\catmark{A} The identification of death with the Banach fixed point is rejected as a category error.\catmark{A} All such material is demarcation/red-team content at Category~\catmark{E}.
\end{governancebox}

\begin{evidencebox}[Psychedelics, NDE, Religious Language: Standing Status]{catE}
\textbf{Interface analogy.} Hoffman-style interface theory is a comparison lens only; \UIDT\ diverges from conscious-agent ontology; no access to external agents, raw $S(x)$, or fundamental reality; citations pending source
verification.\catmark{E/A}\\[0.3em]
\textbf{Globality metrics.} $\Phi_{\mathrm{global}}$, $\Xi$, coherence matrices, PCA/SVD spectra, autocorrelation profiles retained as candidate observables \catmark{D/E}; $C[\rho_{\mathrm{brain}}]\sim\int W(x,y)G(x,y)$ is a toy ansatz only and does not derive consciousness from \UIDT; DMN/entropy/serotonin claims
require external audit.\catmark{A}\\[0.3em]
\textbf{Death boundary.} Mysticism/NDE do not prove \UIDT\ or grant $S(x)$ access; unitarity $\neq$ personal survival; death-as-Banach-fixed-point rejected.\catmark{A}
\end{evidencebox}

%=============================================================================
% CHAPTER 16
%=============================================================================
\newpage
\section{Comparative Frameworks and the Demarcation Matrix}
\label{sec:orch-or-comparison}

\subsection{Orchestrated Objective Reduction as a Neutral Contrast}
\label{subsec:orch-or}

To locate \UIDT's observer programme in the wider landscape of theories that connect physics to mind, the framework selects one contrast case for detailed juxtaposition, orchestrated objective reduction, the Penrose--Hameroff proposal, and treats it as exactly that: a neutral contrast, not a foil to be defeated.\catmark{E} Orchestrated objective reduction posits that consciousness arises from quantum computations in neuronal microtubules, terminated by a gravitationally induced objective collapse of the wavefunction. Its interest for the present purpose is that it differs from \UIDT\ on precisely the axis where \UIDT\ has taken a strong position: orchestrated objective reduction asserts an \emph{objective physical collapse}, whereas \UIDT, in the epistemic tradition of quantum Bayesianism, denies physical collapse and reads quantum probabilities as the structural ignorance of a compression-limited observer (\cref{subsec:observer-stratification}).\catmark{E}

Three cautions govern the comparison. First, \UIDT\ does not collapse all quantum-mind theories into orchestrated objective reduction; the latter is one specific proposal, and the broader field, integrated information theory~\cite{tononi2004}, global workspace accounts, and others, comprises distinct programmes with distinct commitments that must not be conflated.\catmark{A} Second, \UIDT\ does not render orchestrated objective reduction obsolete; the two are different research programmes resting on different premises, and \UIDT's denial of objective collapse is a difference of commitment, not a refutation.\catmark{A} Third, any detailed empirical claim about orchestrated objective reduction, about microtubule coherence times, about specific experimental results, requires external citation audit before it could be stated here, and the framework therefore confines itself to the level of structural commitments, which are not in dispute.\catmark{A} The contrast is used to sharpen what \UIDT\ is by showing clearly what it is not.

\subsection{The Consciousness--Physics Demarcation Matrix}
\label{subsec:demarcation-matrix}

The framework's preferred instrument for situating itself among theories of consciousness is a comparison matrix, and its purpose must be stated before its content: it is a \emph{demarcation instrument}, not a scorecard. The matrix is not constructed to demonstrate that \UIDT\ wins; it is constructed to make the commitments and the weaknesses of each framework, including \UIDT, explicit and comparable along neutral axes.\catmark{A} Those axes are parsimony, mathematical maturity, empirical status, falsifiability, and category-error risk, together with the structural descriptors of each theory's primitive, mechanism, claimed evidence, scale bridge, collapse commitment, and observer role. Throughout, the \UIDT\ observer material is marked Category~\catmark{E}, because it has not been formalised or tested, and the framework's own weaknesses are recorded as plainly as anyone else's.\catmark{E}

{\small
\renewcommand{\arraystretch}{1.3}
\begin{longtable}{p{0.20\textwidth} p{0.24\textwidth} p{0.24\textwidth} p{0.22\textwidth}}
\caption{Consciousness--physics demarcation matrix. A neutral comparison
instrument; \UIDT\ observer material is Category~[E]. The matrix records
weaknesses as well as commitments and is not a ranking.}
\label{tab:demarcation-matrix}\\
\toprule
\textbf{Axis} & \textbf{UIDT (observer programme)} & \textbf{Orch.\ Obj.\ Reduction} &
\textbf{Integrated Information} \\
\midrule
\endfirsthead
\toprule
\textbf{Axis} & \textbf{UIDT} & \textbf{Orch.\ OR} & \textbf{IIT} \\
\midrule
\endhead
\bottomrule
\endfoot
Primitive & Real scalar $S(x)$; observer is downstream compression \catmark{E} &
Quantum state in microtubules + gravitational collapse & Intrinsic cause--effect
structure ($\Phi$) \\
Mechanism & Coarse-graining / forgetful functor (ansatz) \catmark{E} &
Objective reduction terminating quantum computation & Maximally irreducible
integrated information \\
Claimed evidence & None for the observer claims; physics strata separate
\catmark{E} & Contested; requires external audit & Correlational; contested
interpretation \\
Scale bridge & Missing (MeV-GeV $\to$ neural eV-meV gap, $\sim$6--9 orders) \catmark{A} &
Microtubule quantum$\to$cognition, disputed & Substrate-level to system-level \\
Collapse claim & No physical collapse (QBist/epistemic) \catmark{E} &
Objective physical (gravitational) collapse & No collapse claim (intrinsic) \\
Observer role & Compression-limited relational access \catmark{E} &
Locus of quantum computation & The integrated system itself \\
Falsifiability & Observer claims not yet falsifiable \catmark{E}; physics claims
falsifiable separately & Some proposed tests, contested & Disputed/operational
challenges \\
Category-error risk & Information$\to$mind slide; MeV$\to$neural leap (both
flagged) \catmark{A} & Quantum$\to$phenomenal leap & Axioms$\to$physics mapping \\
\bottomrule
\end{longtable}
}

\noindent Several features of \cref{tab:demarcation-matrix} are deliberately uncomfortable for \UIDT, and that discomfort is the point. The framework's observer programme offers \emph{no} evidence for its observer-level claims; its scale bridge is openly marked missing; its observer claims are, as yet, not falsifiable; and its characteristic category-error risks, the slide from information to mind, the leap from mega-electron-volt physics to neural dynamics, are named in its own cell rather than hidden. A matrix built to flatter \UIDT\ would not contain these admissions. Their presence is what makes the instrument a demarcation tool rather than an advertisement.\catmark{A}

What \UIDT\ can claim, on the neutral axes, is a measure of parsimony at the level of its physical primitive, a single real scalar, and a clear separation between its falsifiable physics strata and its interpretive observer stratum. What it cannot claim is mathematical maturity or empirical status for the observer programme, both of which remain undeveloped. The honest reading of the matrix is therefore mixed: \UIDT\ is parsimonious and disciplined in its stratum-separation, and simultaneously immature and untested where consciousness is concerned. Presenting both halves of that judgement, side by side with competing frameworks subjected to the same axes, is the entire function of the exercise.\catmark{E}

\begin{evidencebox}[Comparative Frameworks: Standing Status]{catE}
\textbf{Orch.\ OR.} Used as a neutral contrast (objective collapse vs.\ \UIDT's QBist no-collapse); not a foil; \UIDT\ neither subsumes nor obsoletes it; detailed
empirical claims pending citation audit.\catmark{A/E}\\[0.3em]
\textbf{Matrix.} A demarcation instrument on neutral axes (parsimony, mathematical maturity, empirical status, falsifiability, category-error risk), not a ranking; \UIDT\ observer material marked \catmark{E}; \UIDT's own weaknesses (no observer evidence, missing scale bridge, non-falsifiable observer claims, named category-error risks) recorded explicitly.\catmark{A}
\end{evidencebox}

%=============================================================================
% PART VII
%=============================================================================
\newpage
\part{Closed-Circle Synthesis, Falsification, and Governance}
\label{part:synthesis-governance}

\noindent The final part closes the circle the manuscript has been drawing. It returns to the empirical ground of Part~I, now equipped with everything the intervening parts established, and traces the only honest path from data to interpretation. It assembles the framework's final allowed claims and its final forbidden claims into explicit tables. It states the falsification gates that any future work must pass and the tensions that remain unresolved. And it specifies, without inflation, what \UIDT\ version~3.9 is and is not: an auditable research ontology whose strongest results are formal boundary results and disciplined negative findings, whose calibrated parameters remain calibrated, and whose deepest questions remain open. The closing discipline is the same as the opening one: the value of the framework lies in its falsifiability, not in its ambition.
\newpage
%=============================================================================
% CHAPTER 17 (NEW) — COSMOLOGICAL CALIBRATION
%=============================================================================
\section{Cosmological Calibration}
\label{sec:cosmological-calibration}

\noindent The claims table of \cref{sec:final-claims} carries one entry, ONT-09, that the preceding chapters did not yet support with a section: the framework's cosmological parameters. This chapter supplies that support, and it does so under the strictest evidence cap in the entire manuscript. Every cosmological quantity here is Category~\catmark{C}: calibrated to external data, never derived from $\Sx$, and never presented as resolving an observational tension. The framework's cosmological role is that of a consumer of empirical inputs, not a producer of predictions; the section exists to make that consumption transparent and to keep ONT-09 from being an orphan claim.\catmark{C}

\subsection{Empirical Inputs and the Calibration Boundary}
\label{subsec:cosmology-inputs}

Three empirical domains bound the framework's effective large-scale geometry: the Hubble expansion rate, the dark-energy equation of state, and the matter fluctuation amplitude. In each, the framework calibrates to the published measurement and explicitly declines to resolve the associated tension.

\paragraph{Hubble rate $H_0$.}
The defining feature of the current $H_0$ situation is a persistent $\sim 5\sigma$ discrepancy between early- and late-universe determinations. The \textit{Planck} 2018 analysis of the cosmic-microwave-background acoustic scale, in base $\Lambda$CDM, gives $H_0 = 67.4 \pm 0.5\ \mathrm{km\,s^{-1}\,Mpc^{-1}}$ \cite{planck2018}, while the SH0ES local distance-ladder programme reports $H_0 = 73.04 \pm 1.04\ \mathrm{km\,s^{-1}\,Mpc^{-1}}$ \cite{riess2022}. Combined in quadrature these published uncertainties place the two determinations roughly $4.9\sigma$ apart. The framework calibrates its macroscopic metric parameters to this empirical situation and records the disagreement in the tension ledger (\cref{sec:tension-ledger}); it does \emph{not} claim to resolve the Hubble tension, and any such claim is a forbidden overclaim.\catmark{C}

\paragraph{Dark-energy equation of state $w_0, w_a$.}
In the Chevallier--Polarski--Linder parametrisation $w(a)=w_0+w_a(1-a)$, a pure cosmological constant requires $w_0=-1,\ w_a=0$. The DESI Data Release~2 baryon acoustic oscillation measurements report a preference for the quadrant $w_0>-1,\ w_a<0$ (evolving dark energy), with a statistical significance that depends on the data combination, a mild $\sim 2.3\sigma$ BAO--CMB tension in $\Lambda$CDM, rising to $\sim 2.8$--$4.2\sigma$ for $w_0$--$w_a$ fits when supernova samples are added \cite{desidr2}. A subsequent parameter extraction reports, for the DESI+CMB combination, central values near $w_0 \approx -0.42 \pm 0.21$ and $w_a \approx -1.75 \pm 0.58$ (a $\sim 3.1\sigma$ preference), shifting to $w_0 \approx -0.84 \pm 0.06$, $w_a \approx -0.62$ with Pantheon+ and $w_0 \approx -0.67 \pm 0.09$, $w_a \approx -1.09$ with Union3; these combination-dependent values are carried as Stratum-I empirical inputs at Category~\catmark{C}, flagged for final confirmation against the published DESI DR2 table. The combinations that include supernovae trace an equation of state that crosses $w=-1$ at intermediate redshift, a phantom crossing that a single canonical real scalar cannot smoothly realise without a ghost instability, which is precisely why the framework imports these as external constraints rather than deriving them from $\Sx$. The framework does not derive phantom-crossing or quintessence behaviour from $\Sx$.\catmark{C}

\paragraph{Matter fluctuation amplitude $S_8$.}
The parameter $S_8\equiv\sigma_8\sqrt{\Omega_m/0.3}$ is currently survey-dependent rather than a single late-time consensus value: the completed KiDS-Legacy cosmic shear analysis reports $S_8 = 0.815^{+0.016}_{-0.021}$ (agreeing with \textit{Planck} at $\sim 0.73\sigma$), while the DES Year-6 $3\times2$pt analysis reports a lower $S_8 = 0.789 \pm 0.012$ (a $\sim 2.6\sigma$ projected CMB difference) \cite{kids2025,des2026}. The framework treats $S_8$ as a survey-and-pipeline-dependent empirical input that its relational coarse-graining parameters are bounded by; it calibrates to the spread, and does not relabel it as a resolved tension.\catmark{C}

\begin{table}[H]
\footnotesize
\caption{Stratum-I empirical inputs to which \UIDT\ calibrates. All entries are
Category~[C]; none is a \UIDT\ prediction or a tension resolution.}
\label{tab:cosmology-inputs}
\renewcommand{\arraystretch}{1.3}
\begin{tabular*}{\textwidth}{@{}p{0.18\textwidth}@{\extracolsep{\fill}}p{0.35\textwidth}p{0.24\textwidth}c@{}}
\toprule
\textbf{Parameter} & \textbf{Value $\pm$ uncertainty} & \textbf{Source} & \textbf{Class} \\
\midrule
$H_0$ (early) & $67.4 \pm 0.5\ \mathrm{km\,s^{-1}\,Mpc^{-1}}$ & Planck 2018 & \catmark{C} \\
$H_0$ (late) & $73.04 \pm 1.04\ \mathrm{km\,s^{-1}\,Mpc^{-1}}$ & SH0ES (Riess 2022) & \catmark{C} \\
$w_0, w_a$ & $w_0\!\approx\!-0.42\pm0.21,\ w_a\!\approx\!-1.75\pm0.58$ (DESI+CMB; combination-dependent) & DESI DR2 & \catmark{C} \\
$S_8$ & $0.815^{+0.016}_{-0.021}$ / $0.789\pm0.012$ & KiDS-Legacy / DES~Y6 & \catmark{C} \\
\bottomrule
\end{tabular*}
\end{table}

\subsection{The Relational-Proliferation Scenario: Interpretive Status}
\label{subsec:information-big-bang}

The framework's developmental corpus contains a cosmogonic narrative, and it is retained here strictly as Category~\catmark{E}\ interpretation within Stratum~III, a way of thinking, fed into a Category~\catmark{C}\ calibration, not a derivation of cosmological data.\catmark{E} The narrative reads the ``beginning'' not as a physical instant $t=0$ but as the first logical distinction in the sense of \cref{subsec:sx-zero-not-unmarked}: the unmarked state has no temperature, no fluctuation spectrum, and no cosmological epoch, so an ``information big bang'' marks the onset of relational differentiation rather than the ignition of a pre-existing spacetime.

On this reading, cosmic inflation is reinterpreted, again at \catmark{E}, as a rapid proliferation of nodes and edges in the pre-geometric relational graph of \cref{subsec:relational-differentiation}: when that exploding relational network is viewed under coarse-graining, the proliferation appears, in the effective field regime, as an exponentially expanding spacetime metric. The picture is deliberately bounded. It does not compute an inflationary observable, it does not predict a primordial spectrum, and it does not touch the Category~\catmark{C}\ cap on the parameters of \cref{subsec:cosmology-inputs}; it is a relational gloss on geometry that must, like every interpretive element in the manuscript, earn its place by clarity rather than by evidential weight. Any cosmological prediction derived from it would remain capped at \catmark{C}\ and would have to pass the same falsification gates as every other claim.\catmark{E}

\begin{evidencebox}[Cosmology: Standing Status]{catC}
\textbf{Inputs.} $H_0$ (Planck $67.4\pm0.5$ vs SH0ES $73.04\pm1.04$, $\sim5\sigma$), $w_0$--$w_a$ (DESI DR2: $w_0>-1,w_a<0$, combination-dependent), $S_8$ (KiDS-Legacy / DES~Y6, survey-dependent): all \catmark{C}, calibrated, not
resolved.\catmark{C}\\[0.3em]
\textbf{Information big bang / inflation.} Relational-proliferation gloss retained as \catmark{E}\ interpretation feeding a \catmark{C}\ calibration; no inflationary
observable derived; cosmology cap not breached.\catmark{E}\\[0.3em]
\textbf{Discipline.} \UIDT\ consumes cosmological data; it does not predict or resolve. ONT-09 is hereby supported.\catmark{C}
\end{evidencebox}

%=============================================================================
% CHAPTER 18 (was 17) — CLOSED-CIRCLE SEQUENCE
%=============================================================================
\section{The Closed-Circle Sequence}
\label{sec:closed-circle}

\subsection{The Empirical-to-Ontology Return Path}
\label{subsec:return-path}

The architecture of this manuscript has traced a single large loop, and it is now possible to close it explicitly. The loop runs as follows, and each step carries its proper evidential status so that the closure inflates nothing.\catmark{A/D/E}

One begins from the empirical and lattice stress tests of pure $\mathrm{SU}(3)$ Yang--Mills theory, the benchmark batteries of \cref{sec:benchmark-batteries}, held at status \catmark{B}\,pending their source and covariance audit. One then computes the \UIDT-side quantities \emph{blindly}, by procedures that cannot read their targets, in accordance with the target-leakage theorem (\cref{prop:target-leakage}); these computations are predictive, status~\catmark{D}. One compares the blind outputs with the benchmarks \emph{post hoc}, through an explicit uncertainty and covariance model, never by an internal-closure threshold, again status \catmark{B}\,pending. From the comparison one identifies structural successes and failures, a step that yields Class~A facts about what does and does not work alongside predictive assessments.\catmark{A/D} And only then, after the constraints are known, does one interpret, placing the interpretive reading last, at Category~\catmark{E}, rather than allowing it to colour the computation that preceded it.\catmark{E}

The closure of the circle is governed throughout by two honest structural limits that the manuscript has established and that must remain in view at the point of synthesis: the $d^2 = 0$ obstruction of \cref{prop:d2-obstruction}, which shows that the naive scalar-gradient route to gauge curvature fails identically, and the $G_{\mathrm{SM}}$-Origin Gap of \cref{oq:gsm-gap}, which records that no derivation of the Standard-Model gauge group from the scalar exists.\catmark{A/D} The circle therefore does not close on a triumph. It closes on a disciplined accounting: a framework that has stated its ontology, built its numerical architecture, identified precisely where its most natural constructions fail, and returned to the empirical comparison with its interpretive claims held at arm's length. Closing the theoretical circle without evidence inflation is the whole achievement of the synthesis, and it is a real achievement precisely because it refuses to be a larger one.\catmark{A}

\subsection{Final Allowed Claims}
\label{sec:final-claims}

The claims the framework is entitled to assert, with their conservative evidence classes and strata, are collected in \cref{tab:claims}. The table is the authoritative register referenced from the evidence-grammar discussion of \cref{subsec:evidence-grammar} and from the evidence-class determination of \cref{box:pi-override-delta}; it is deliberately conservative, and several of its entries are negative results or open programmes rather than positive findings.\catmark{A}

{\small
\renewcommand{\arraystretch}{1.3}
\begin{longtable}{p{0.10\textwidth} p{0.56\textwidth} p{0.10\textwidth} p{0.12\textwidth}}
\caption{Final allowed claims of \UIDT\ v3.9, with conservative evidence classes
and strata. Negative and open entries are included deliberately.}
\label{tab:claims}\\
\toprule
\textbf{ID} & \textbf{Claim} & \textbf{Class} & \textbf{Stratum} \\
\midrule
\endfirsthead
\toprule
\textbf{ID} & \textbf{Claim} & \textbf{Class} & \textbf{Stratum} \\
\midrule
\endhead
\bottomrule
\endfoot
ONT-01 & $S(x)=0$ is a marked state, not the unmarked state. & \catmark{A} & II/III \\
ONT-02 & A morphism $f\colon A\to B$ presupposes distinguishable $A,B$; morphism is a formal tool, not an ontological generator. & \catmark{A} & II/III \\
ONT-03 & With $A=\dd S$, the curvature $F=\dd^2 S=0$ identically; the naive scalar-gradient route to non-trivial Yang--Mills curvature fails. & \catmark{A} & II \\
ONT-04 & $\gamma = \GammaGeom$ is calibrated, never derived from RG first principles. & \catmark{A-} & III \\
ONT-05 & $\Delta = \DeltaGap$: internally consistent, lattice-compatible; PI-overridden to~B (\cref{box:pi-override-delta}); not externally peer-confirmed, not derived. & \catmark{B} & I/III \\
ONT-06 & $\kappa = \KappaRG$ exact framework input (the fixed-point reading is a candidate condition, untested by any defined flow); the internal coupling constraint $5\kappa^2 = 3\lambda_S$ closes internally to a residual $<10^{-14}$. The empirically calibrated reading $\kappa = 0.500 \pm 0.008$ is a separate, lower-class statement and is not the basis of this row. & \catmark{A} & II \\
ONT-07 & Pure $\mathrm{SU}(3)$ ratio comparisons are post-hoc stress tests with solver access disabled. & \catmark{B\,pending} & I/III \\
ONT-08 & $S(x)\to\mathcal{A}\to G_{\mathrm{SM}}$ is the v4.x research programme; the algebraic fork (strict scalar vs.\ enriched primitive) is open. & \catmark{D} & III \\
ONT-09 & Cosmological parameters are calibrated, capped at~C; no tension is declared resolved. & \catmark{C} & I/III \\
ONT-10 & The observer functor $F\colon V\to\mathrm{Obs}$ and the Class~0--III taxonomy are bounded interpretive frameworks. & \catmark{E} & III \\
\bottomrule
\end{longtable}
}

\noindent The conservatism of \cref{tab:claims} is its principal feature. Three of its ten entries (ONT-01, ONT-02, ONT-03) are formal negative or boundary results, the kind of claim a framework can hold with the most confidence precisely because it asserts a limit rather than a discovery. One (ONT-04) is permanently calibrated. One (ONT-05) carries an explicit reclassification. Two (ONT-07, ONT-08) are pending or programmatic. One (ONT-09) is capped. And one (ONT-10) is openly interpretive. This is the distribution of a framework that has disciplined itself: its firmest ground is its boundary results, and its ambitions are labelled as ambitions.\catmark{A}

\subsection{Final Forbidden Claims}
\label{subsec:forbidden-claims}

Equal in importance to what the framework may assert is what it may not, and the prohibitions are stated here as the manuscript's closing act of discipline. Each is a Class~A constraint, not a matter of current evidence that might shift, but a standing rule about what \UIDT\ is not permitted to claim given its present state.\catmark{A}

\begin{governancebox}[Final Forbidden Claims]
\textbf{(1)} No derivation of the full Standard-Model gauge group
$G_{\mathrm{SM}}$ from the scalar $S(x)$.\\[0.3em]
\textbf{(2)} No derivation of the Born rule.\\[0.3em]
\textbf{(3)} No solution of consciousness or qualia.\\[0.3em]
\textbf{(4)} No validation of $S(x)$ by psychedelic, near-death, or mystical
experience.\\[0.3em]
\textbf{(5)} No use of $E_T$ as a threshold for life, death, ego, psychedelic
states, or artificial general intelligence.\\[0.3em]
\textbf{(6)} No identification of $\Delta$ with the $0^{++}$ glueball without
operator matching (L12).\\[0.3em]
\textbf{(7)} No use of quarantined museum artifacts as derivation or factual
evidence.\\[0.3em]
\textbf{(8)} No presentation of $5\kappa^2 = 3\lambda_S$ as an independent
empirical discovery where $\lambda_S$ is fixed definitionally by that same
constraint.\\[0.3em]
\textbf{(9)} No claim that DESI, $H_0$, $S_8$, or $w_0$--$w_a$ data validate
\UIDT\ beyond calibrated Category~\catmark{C}\ status.\\[0.3em]
\textbf{(10)} No claim that a lattice-compatible number establishes the
Standard-Model gauge group, a confinement mechanism, or a Clay-grade
Yang--Mills construction.\\[0.3em]
\emph{All ten are Class~\catmark{A}\ standing constraints.}
\end{governancebox}

\noindent The seventh prohibition deserves a closing emphasis, because it guards the integrity of all the others. The museum of refuted constructions, the numerological decomposition of $\gamma$, the failed scalar-gradient operator, the toy ODE dressed in renormalisation-group language, is retained for its documentary value, as evidence of method. But its contents may never be cited as factual evidence for any claim.\catmark{A} A museum artifact is a record of a road not taken; to drive on it again, having recorded why it was abandoned, would be the most basic failure of the discipline this manuscript has tried to embody. The forbidden claims, taken together, are not a confession of weakness. They are the precise statement of a framework that knows the boundary of its own warrant and declines, on principle, to cross it.\catmark{A}

%=============================================================================
% CHAPTER 18
%=============================================================================
\newpage
\section{Falsification Gates and the Tension Ledger}
\label{sec:falsification-tension}

\subsection{The Falsification Gates}
\label{subsec:falsification-gates}

The framework's discipline is enforced, in practice, by a set of gates through which every proposed claim, computation, or contribution must pass. It is essential to the governance boundary of \cref{subsec:governance-boundary} that these gates are understood correctly: they are \emph{advisory and non-approving}. A gate can \emph{block}, it can flag a violation and withhold its clearance, but a gate that is passed authorises nothing. Passing all eight gates is necessary hygiene, never sufficient warrant for any action; confirmation remains, by the principle that internal checks do not confer external validity, a separate human matter.\catmark{A}

\begin{governancebox}[The Eight Falsification Gates (Advisory, Non-Approving)]
\textbf{F1: No target leakage.} No solver may access the value it claims to
predict (\cref{prop:target-leakage}).\\[0.3em]
\textbf{F2: Local precision discipline.} Working precision is set locally per proof-critical block; no global precision mutation; no proof-critical
floating-point or rounding heuristics.\\[0.3em]
\textbf{F3: Complete benchmark accounting.} Every external benchmark carries a
verified source, an uncertainty budget, and a covariance model. \catmark{A/B\,pending}\\[0.3em]
\textbf{F4: L12 separation.} $\Delta$ is not identified with the $0^{++}$
glueball without operator matching (\cref{lim:L12}).\\[0.3em]
\textbf{F5: $G_{\mathrm{SM}}$-Origin honesty.} The gauge-origin gap
(\cref{oq:gsm-gap}) is stated as open; no derivation is claimed.\\[0.3em]
\textbf{F6: Observer evidence boundary.} No phenomenological report validates the ontology; the admissible chain (\cref{subsec:admissible-chain}) is
respected.\\[0.3em]
\textbf{F7: Born-rule status.} The Born rule is recorded as borrowed and
open, never as a \UIDT\ derivation.\\[0.3em]
\textbf{F8: Governance boundary.} No machine-generated verdict authorises a
merge, freeze, or evidence-class upgrade.\\[0.3em]
\emph{Each gate is Class~\catmark{A}. Passing is advisory; it approves nothing.}
\end{governancebox}

\noindent The reframing of these gates into advisory, non-approving terms is itself one of the corrections this edition records. In the developmental history, automated checks sometimes carried language, ``approved,'' ``pass,'' ``merge-ready'', that allowed a clean check to function as authorisation. The gates above are deliberately stripped of that force. They report the presence or absence of a violation and nothing more. F8 makes the principle explicit and self-referential: not even the gate system may authorise, and a green board of F1 through F7 is a statement about hygiene, not a permission.\catmark{A}

\subsection{The Tension Ledger}
\label{sec:tension-ledger}

A framework committed to falsification must keep its unresolved tensions in plain view rather than smoothing them away. The tension ledger of \cref{tab:tensions} is referenced from the evidence-grammar discussion of \cref{subsec:evidence-grammar} and the discipline of $\gamma$ in \cref{subsec:gamma-discipline}; each entry records a genuine disagreement or gap that the framework has not closed.\catmark{A}

{\small
\renewcommand{\arraystretch}{1.3}
\begin{longtable}{p{0.08\textwidth} p{0.40\textwidth} p{0.22\textwidth} p{0.18\textwidth}}
\caption{Tension ledger. Unresolved disagreements and gaps, displayed rather than
smoothed. Each is a standing item for future work.}
\label{tab:tensions}\\
\toprule
\textbf{ID} & \textbf{Tension / gap} & \textbf{Status} & \textbf{Resolution path} \\
\midrule
\endfirsthead
\toprule
\textbf{ID} & \textbf{Tension / gap} & \textbf{Status} & \textbf{Resolution path} \\
\midrule
\endhead
\bottomrule
\endfoot
T1 & $\gamma_{\mathrm{MC}} = \GammaMC$ versus calibrated $\gamma = \GammaGeom$. & \textbf{[consistent within MC variance]}; the two determinations agree to $0.035$ ($0.035\sigma$ given the $\pm1.005$ spread) and are displayed separately, not merged by selection. & Reduce MC variance; non-leaking derivation of $\gamma$. \\
T2 & $\Delta = \DeltaGap$ evidence class. & Would sit \catmark{E\,pending} absent external work; PI-overridden to \catmark{B}\ (\cref{box:pi-override-delta}). & External peer counter-signature. \\
T3 & $\Delta$ versus $0^{++}$ glueball identity. & \catmark{A}\ separation enforced (L12, \cref{lim:L12}); identity unproven. & Operator matching. \\
T4 & SVZ leading-order topology route versus lattice topology. & \catmark{A/B\,pending}; SVZ route is museum-only. & Audited lattice $\chi_{\mathrm{top}}$ ratios. \\
T5 & $E_T$ observer/neural scale bridge. & Absent; $\sim$6--9 orders of magnitude (MeV--GeV vs.\ eV--meV). \catmark{A/E} & Derive or refute a principled scale bridge. \\
T6 & Born-rule derivation. & Absent; borrowed from Stratum~II. \catmark{A} & Non-circular coarse-graining derivation. \\
T7 & Cosmological tensions ($H_0$, $S_8$). & Capped at \catmark{C}; calibrated, \emph{not} closed. & Independent prediction, not fit. \\
\bottomrule
\end{longtable}
}

\noindent Two entries warrant a closing remark in the format of the framework's tension-alert protocol. Entry~T1 is the framework's flagged tension between two internal determinations of $\gamma$: the calibrated kinetic value $16.339$ and the Monte-Carlo mean $16.374 \pm 1.005$. The two are consistent within the (large) Monte-Carlo uncertainty, but they are displayed together rather than reconciled by quiet selection of either; the resolution path is a reduction of the Monte-Carlo variance together with a non-leaking derivation of $\gamma$, and until then the disagreement stands on the record.\catmark{A} Entry~T2 records the status of $\Delta$ honestly: absent the external work that the framework's own evidence discipline requires, the value would rest at \catmark{E}\,pending; the recorded determination sets it to \catmark{B}, with the rationale and the pending external counter-signature documented in \cref{box:pi-override-delta}.\catmark{A}. The ledger is, in a sense, the most honest page in the manuscript. It is the framework stating, in one place, every disagreement it has not resolved and every bridge it has not built. A framework that hid this page would be easier to admire and impossible to trust. The decision to keep it visible, to treat unresolved tension as information to be displayed rather than noise to be filtered, is the operational meaning of the falsification-first discipline that the whole manuscript has claimed to follow.\catmark{A}

%=============================================================================
% CHAPTER 19
%=============================================================================
\section{Manuscript Architecture and the Final Synthesis}
\label{sec:final-synthesis}

\subsection{The Required Tables}
\label{subsec:required-tables}

A research ontology that claims auditability must be navigable, and navigability here means that every class of claim has a dedicated, locatable register. The manuscript is therefore organised around a fixed set of tables, some already presented in the body and some maintained as living repository artifacts. The required set is: the Claims Table (\cref{tab:claims}), the Tension Ledger (\cref{tab:tensions}), a Museum Artifact Table cataloguing every refuted construction, the $G_{\mathrm{SM}}$-Origin Literature Matrix, the Benchmark Matrix of post-hoc pure-gauge ratios, the Observer Taxonomy Table, and the Phenomenology Boundary Table.\catmark{A/D/E} The first two exist in the body as authoritative registers; the remainder are maintained alongside the manuscript and summarised here.

The $G_{\mathrm{SM}}$-Origin Literature Matrix collects the research vectors of \cref{subsec:minimal-algebra-catalogue} with their exact roles, and a compact form is given in \cref{tab:literature-matrix}. Its purpose is to keep the forward programme honest: each vector is a direction, none is a closure.

The full matrix evaluates four established programmes as candidates for the intermediate algebra $\mathcal{A}$, along the axes that determine whether any of them could honestly mediate $S(x)\to\mathcal{A}\to G_{\mathrm{SM}}$. Every entry is held at \catmark{D}/\catmark{E}: each programme supplies an $\mathcal{A}$-like structure in its own setting, but none is a scalar-only \UIDT\ derivation, and the ``forbidden inference'' row records, for each, the specific overclaim the framework refuses.

{\scriptsize
\renewcommand{\arraystretch}{1.25}
\begin{longtable}{p{0.155\textwidth} p{0.185\textwidth} p{0.185\textwidth} p{0.185\textwidth} p{0.185\textwidth}}
\caption{$G_{\mathrm{SM}}$-origin literature matrix (A-candidate matrix).
Axes as rows; candidate programmes as columns. All entries \catmark{D}/\catmark{E};
none is a closed solution. ``Encodes'' means the gauge group is built into the
construction's inputs; ``Selects'' means it is constrained by the construction's
axioms.}
\label{tab:literature-matrix}\\
\toprule
\textbf{Axis} & \textbf{NCG (Connes / Chamseddine)} & \textbf{Quantum-link / D-theory (Wiese)} & \textbf{String-net / fusion (Levin--Wen)} & \textbf{Gauged tensor / PEPS (Verstraete--Cirac)} \\
\midrule
\endfirsthead
\toprule
\textbf{Axis} & \textbf{NCG} & \textbf{QLM / D-theory} & \textbf{String-net} & \textbf{Tensor / PEPS} \\
\midrule
\endhead
\bottomrule
\endfoot
Primitive input & Spectral triple, finite algebra $\mathcal{A}_F = \mathbb{C}\oplus\mathbb{H}\oplus M_3(\mathbb{C})$ & Finite link Hilbert space (rishon operators) & Bosonic lattice, tensor/fusion categories & Virtual entanglement symmetries \\
Selects or encodes $G_{\mathrm{SM}}$? & Selects / strongly constrains under spectral-action and finite-geometry axioms & Encodes (symmetry in operators) & Encodes (via fusion/branching rules) & Encodes (via virtual tensor symmetry) \\
Scalar-only compatible? & No (requires matrices) & No (requires rishon/fermion-like d.o.f.) & No (requires rich internal labels) & No (requires multi-index tensors) \\
Requires noncommutativity? & Yes (essential for $\mathrm{SU}(N)$) & Yes (non-commuting link operators) & Depends (non-abelian anyons) & Yes (non-abelian virtual symmetry) \\
Supports chirality? & Yes (Dirac operator) & Yes (domain-wall / overlap) & Hard (challenging in 3+1D) & Hard (chiral lattice-fermion problem) \\
Fermion generations? & Accommodated / constrained in model-dependent ways; exact three-generation derivation not established & Engineerable; not selected from first principles & Emergent or model-dependent; species count variable & Engineered into physical indices; selection not automatic \\
Anomaly cancellation? & Yes (geometric requirement) & Yes (lattice-regulated) & Context-dependent & Varies \\
Supports 3+1D continuum? & Yes (almost-commutative manifolds) & Yes (D-theory dim.\ reduction) & Hard (Walker--Wang, active) & Theoretically yes; computationally hard \\
\UIDT\ compatibility & High (template for ultimate $\mathcal{A}$) & High (discrete edge-algebra template) & Medium (mechanics of emergence) & High (coarse-graining language) \\
Ontology cost & Very high (spectral machinery) & Medium (discrete QM) & Medium (topological inputs) & Low--medium (quantum information) \\
Forbidden \UIDT\ inference & Assuming $S(x)$ natively has $M_3(\mathbb{C})$ & Assuming $\mathrm{SU}(3)$ links emerge for free & Claiming full 3+1D chiral SM is solved & Assuming tensors \emph{create} the symmetry \\
Evidence status & \catmark{D}/\catmark{E} & \catmark{D}/\catmark{E} & \catmark{D}/\catmark{E} & \catmark{D}/\catmark{E} \\
\bottomrule
\end{longtable}
}

\noindent The matrix licenses one strategic statement, which the framework adopts as the precise form of its forward claim and which is itself classified \catmark{D}: there already exist serious candidates for the missing algebraic intermediary $\mathcal{A}$, so the open \UIDT\ question is \emph{not} whether algebra can lead to Standard-Model structure, it demonstrably can, in several established programmes, but whether the scalar $S(x)$ can \emph{generate}, \emph{select}, or \emph{non-circularly carry} one of these $\mathcal{A}$-like structures without importing $G_{\mathrm{SM}}$ by hand. Read against the ``scalar-only compatible?'' row, every candidate currently answers ``no,'' which is precisely why the algebraic fork of \cref{subsec:the-fork} remains open and why noncommutative geometry, the only column that \emph{selects} rather than merely encodes, is recorded as the strategic Arm~2 direction rather than as a closure.\catmark{D}

\subsection{The Required Deliverables}
\label{subsec:required-deliverables}

The manuscript is accompanied by a set of public-facing supplementary materials. These are intended for public release only, and contain no private working material.\catmark{A} The set comprises: the $G_{\mathrm{SM}}$-origin literature matrix; a record of the failed scalar-gradient route; a schema declaring benchmarks post-hoc only; the separated blind-computation and post-hoc-comparison scripts; a record of the observer taxonomy at Category~\catmark{E}\ status; a statement of the $E_T$ scale-bridge requirements; a register of forbidden phenomenological inferences; and the cleaned manuscript itself, accompanied by a compile log and a bibliography audit.\catmark{A/D}

The bibliography audit is not a formality. Consistent with the framework's audit policy, external citations must be resolved to verifiable sources before they are relied upon; references that cannot be so resolved are marked pending and are not used as evidence.\catmark{A} The framework's own records, the parent framework, this ontology, and the Clay-submission record, are identified by their persistent identifiers, while the external physics literature invoked in the benchmark and origin discussions is tracked separately and resolved through that audit rather than asserted here.\catmark{A}

\subsection{The Final Synthesis}
\label{subsec:final-paragraph}

It remains to state, as plainly as the manuscript began, what \UIDT\ version~3.9 is.

It is an auditable research ontology and a statement of an effective-framework boundary. It is not a complete proof of anything, and it does not claim to be.\catmark{E/D} Its strongest current results are formal boundary results and disciplined negative findings: that $S(x)=0$ is not the unmarked state, that a morphism presupposes the distinction it cannot generate, that the naive scalar-gradient route to gauge curvature fails identically by $d^2=0$. These are the claims the framework holds with most confidence, and it is characteristic of the discipline this manuscript has tried to enact that its firmest ground is the ground of what cannot be done.\catmark{A}

Its calibrated parameters remain calibrated: $\gamma$ remains \catmark{A-}\ in the present framework state, and may not be upgraded unless a future derivation reproduces it without target access, survives the anti-target-leakage discipline of \cref{prop:target-leakage}, and is recorded through the framework's formal evidence-upgrade procedure with external counter-signature.\catmark{A-} Its benchmark programme against pure $\mathrm{SU}(3)$ Yang--Mills theory remains post-hoc and uncertainty-aware, a set of stress tests the framework may pass or fail but never a foundation it stands upon.\catmark{B\,pending} The full gauge-origin problem remains open through the $G_{\mathrm{SM}}$-Origin Gap, with the algebraic fork between a strict-scalar and an enriched primitive unresolved and reserved to human governance.\catmark{A/D} The observer ontology, the compression picture, the functor, the taxonomy, the bounded phenomenology of ego, consciousness, and altered states, remains a bounded research programme, preserved with philosophical seriousness and labelled, without exception, as Category~\catmark{E}\ within the interpretive stratum.\catmark{E}

The scientific value of the framework, finally, does not rest on any of its positive ambitions. It rests on its falsification-first discipline: on the gates that can block but never approve, on the ledger that displays every unresolved tension, on the museum that preserves every refuted route, and on the governance boundary that forbids confident language from ever substituting for proof or for human judgement. \UIDT\ may, in the end, be wrong. The purpose of this manuscript has been to ensure that if it is wrong, it will be wrong \emph{legibly}, that its failures, when they come, will be findable, stated, and honest. That is the most a research ontology can offer, and the framework offers it without reservation.\catmark{A}

\vspace{1.5em}
\begin{center}
\rule{0.4\textwidth}{0.4pt}\\[0.6em]
{\small\itshape The open gaps are preserved. The discipline is the result.}
\end{center}

%=============================================================================
% PART VIII — PHILOSOPHICAL CONTEXT
%=============================================================================
\newpage
\part{Philosophical Context}
\label{part:philosophical-context}

\noindent This part is the most explicitly interpretive in the manuscript, and it is fenced accordingly. Everything in it is Stratum~III, Category~\catmark{E}: it situates \UIDT\ among philosophical positions and competing research programmes, but it establishes nothing and proves nothing. It is included because a serious ontology should be honest about its philosophical commitments rather than smuggle them in unstated, but the reader should treat this part as context, never as evidence. Where the developmental record stated these positions with more confidence than the framework's discipline now permits, they are downgraded here to positions held rather than results established.\catmark{E}

%=============================================================================
% CHAPTER 20 — EPISTEMOLOGICAL POSITION
%=============================================================================
\newpage
\section{Epistemological Position and Comparisons}
\label{sec:epistemology}

\subsection{Structural Realism, Held as a Position}
\label{subsec:structural-realism}

The ontological proposal that physical reality is constituted by mathematical structure raises an epistemological question: does \UIDT\ claim the universe \emph{is} a mathematical object, or merely that it is \emph{best described by} one? The framework's working answer leans toward the former, an ontic-structural-realist position, but it is held as a philosophical \emph{position}, not asserted as a result, and nothing in the physics strata depends on it.\catmark{E} The proposal is narrow: not that all mathematical structures are physically real (the Tegmark Level~IV view), but that one specific structure, the self-interacting vacuum information density of \cref{eq:lagrangian}, is offered as the candidate constituting structure. The selection is meant to be structural rather than anthropic: the Lagrangian is the minimal gauge-invariant, vacuum-stable, EFT-admissible self-interaction of a real scalar non-minimally coupled to $\mathrm{SU}(3)$ in $3+1$ dimensions. That ``minimality'' is itself a claim about the framework's construction, not a proven uniqueness theorem, and is carried as such.\catmark{E}

The position is explicitly bounded by its own incompleteness. The framework acknowledges that any sufficiently expressive formal system is subject to Gödel-type incompleteness, and it reads its own limitation register (\cref{subsec:limitation-register}) as a concrete, honest catalogue of what it cannot presently derive rather than as a temporary embarrassment. It also distinguishes itself from a simulation hypothesis: \UIDT\ posits no external simulator and no computational substrate; whatever force the structural-realist picture has, it is ontological speculation, not a claim that the universe is being computed by something.\catmark{E}

\subsection{Comparison with Alternative Frameworks}
\label{subsec:alt-frameworks}

It is useful, and disciplined, to state where \UIDT\ sits relative to established programmes, always as comparison, never as a claim of superiority. Against string theory, \UIDT\ is far more ontologically frugal (a single scalar in $3+1$ dimensions, no background metric, no large landscape) but correspondingly far less developed and far less mathematically mature; frugality is not a virtue if it buys its simplicity by leaving the hard problems unsolved, and the $G_{\mathrm{SM}}$-origin gap is exactly such an unsolved problem. Against entropic gravity (Verlinde, Padmanabhan), \UIDT\ shares the intuition that geometry is emergent and thermodynamic in spirit but does not claim to have improved on it; the relationship is one of shared intuition, not demonstrated advance. Against loop quantum gravity, \UIDT\ shares the \emph{aspiration} to background independence but does not yet achieve it in full: its concrete metric-emergence map (\cref{ax:emergence}) is written as a perturbation about a fixed Minkowski background $\eta_{\mu\nu}$, so background independence is at present only partial (\cref{lim:L13}), and it takes the opposite route to geometry (a continuum coarse-graining limit rather than quantised geometric structures); like LQG it has not recovered the full Standard-Model particle content, indeed \UIDT\ has not recovered it at all, which is the content of the open gap.\catmark{E}

The framework's one genuine methodological commitment here is an epistemic razor: it admits no unobservable extra dimensions, no fundamental supersymmetry, and no infinite landscape, and it ties its claims to a small set of parameters with explicit falsification thresholds (\cref{sec:falsification-tension}). This makes the framework vulnerable, and that vulnerability, not any claimed explanatory triumph, is what the framework offers as its scientific virtue.\catmark{E} The earlier developmental claim that \UIDT\ ``eliminates'' dark matter or dark energy is withdrawn: cosmology is capped at \catmark{C} (\cref{sec:cosmological-calibration}), and no such elimination is asserted.

\subsection{Relation to the Information-Theoretic and Structural-Realist Tradition}
\label{subsec:prior-art}

The general move that \UIDT\ makes---taking an information-theoretic primitive as ontologically basic and recovering spacetime, and aspirationally gauge structure, as emergent---is not itself novel, and intellectual honesty requires locating it explicitly within an established tradition rather than presenting it as unprecedented. The lineage begins with Wheeler's ``it from bit'' programme, in which information is treated as the substrate from which physical things derive and in which there is held to be no spacetime continuum at the most fundamental level~\cite{wheeler1990}. It continues through the thermodynamic and entropic-gravity derivations of geometry already invoked above (Jacobson~\cite{jacobson1995}; Padmanabhan~\cite{padmanabhan2010}; Verlinde~\cite{verlinde2011}), the matrix-model emergent-geometry programme (\cref{subsec:minimal-algebra-catalogue}), and group field theory, in which a field on a group manifold is fundamental and spacetime emerges from a condensate. A particularly close ontological cousin is the Carroll--Singh ``Mad-Dog Everettian'' proposal~\cite{carrollsingh2018}, in which the only fundamental objects are a vector in Hilbert space and a Hamiltonian, with space and fields emergent from that minimal data; \UIDT\ differs in retaining a concrete real-scalar field theory rather than an abstract state vector, but shares the metaphysical aim of a maximally frugal primitive. On the philosophical side, the position belongs to the debate over ontic structural realism (\cref{subsec:structural-realism}); it is one specific structure offered as constituting, not merely describing, physical reality, and it should be read as a particular entry in that literature rather than as a new metaphysics.\catmark{E}

Against this background, what is specific to \UIDT\---and what is offered as its actual contribution, as distinct from the shared tradition above---is narrow and concrete: the particular Lagrangian of \cref{eq:lagrangian} with its non-minimal $S\,F^2$ coupling at a stated cutoff; the use of a Banach fixed-point construction to identify the Yang--Mills spectral gap as a unique contraction fixed point; the calibrated dimensionless invariant $\gamma = \GammaGeom$ (\catmark{A-}, calibrated and never derived); and the explicit evidence-grading discipline (\catmark{A} through \catmark{E}) under which every claim in this manuscript is carried. None of these is a claim to have solved the emergence problem that the tradition leaves open; they are the specific, falsifiable form in which \UIDT\ states its version of it.\catmark{E}

%=============================================================================
% CHAPTER 21 — INEVITABILITY, CONTINGENCY, SOCIOLOGY
%=============================================================================
\newpage
\section{Scope, Contingency, and Interpretive Minimality}
\label{sec:inevitability-sociology}

\subsection{From ``Unreasonable Effectiveness'' to a Modest Minimality}
\label{subsec:inevitability}

Wigner's remark on the unreasonable effectiveness of mathematics invites a strong reading: that if the universe \emph{is} mathematical structure, the effectiveness is no longer surprising. The framework records this as a tempting interpretive move but explicitly declines its strong form. To call the effectiveness a ``tautology'' or the structure ``inevitable'' would be exactly the kind of rhetorical overreach the manuscript's discipline forbids; the honest statement is far more modest. Within the framework's own axioms and within $3+1$ dimensions, the $\mathrm{SU}(3)$ Lagrangian of \cref{eq:lagrangian} is presented as the \emph{minimal} self-interacting structure compatible with gauge invariance, vacuum stability, and EFT admissibility. That is a statement about minimality under stated constraints, not a proof of cosmic necessity, and the framework does not claim the universe \emph{had} to be this way.\catmark{E}

The question ``why this universe and not another?'' is therefore answered, if at all, by structural minimality rather than by modal necessity or anthropic selection: among candidate theories, the framework selects the one satisfying its three stated mathematical conditions. The spectral gap $\Delta = \DeltaGap$ remains, in this reading, a structurally determined feature of the construction's fixed point, but its evidence class is unchanged by any of this philosophy: it is \Bdelta, and the philosophical gloss adds no evidential weight.\catmark{E}

\subsection{Reception Risk and Non-Evidential Context}
\label{subsec:sociology}

Finally, the framework reflects, again purely as \catmark{E}\ context, on how a radically simplifying ontology is liable to be received. A framework whose evidentiary classes are not kept explicit can be misread in both directions: dismissed wholesale on account of its speculative stratum, or embraced wholesale on the strength of its calibrated agreements. A framework that proposes to explain with geometry and a single scalar what a particle-centric worldview explains by positing new fields must therefore expect, and plan for, both misreadings, independent of its merits.\catmark{E} The framework's response is not to claim vindication but to insist on the opposite safeguard: the evidence-classification system and the verification discipline exist as bulwarks against \emph{both} premature rejection \emph{and} uncritical acceptance. The right attitude toward a radical simplification is neither enthusiasm nor dismissal but exactly the falsification-first scrutiny the manuscript has tried to embody, and by that standard, the framework's simplifications must earn their place by surviving tests, not by being elegant.\catmark{E}

\begin{evidencebox}[Philosophical Context: Standing Status]{catE}
\textbf{Entire Part~VIII is \catmark{E}, Stratum~III.} Structural realism is a position held, not a result; minimality is a construction claim, not a uniqueness proof; ``inevitability/tautology'' language is withdrawn. Comparisons with string theory, entropic gravity, and LQG are stated without claim of superiority; \UIDT\ is frugal but immature. The ``eliminates dark matter/energy'' claim is withdrawn (cosmology \catmark{C}). No philosophical statement here alters any evidence class.\catmark{E}
\end{evidencebox}

%=============================================================================
% APPENDIX A — THE ADVERSARY TEST
%=============================================================================
\newpage
% [004-appA-gsm HYGIENE FIX, PI review requested]: heading de-conflicted with
% the formal Appendix A introduced by \appendix below; content unchanged.
\section*{The Adversary Test: A Self-Audit}
\addcontentsline{toc}{section}{The Adversary Test: A Self-Audit}
\label{app:adversary}

\noindent The framework maintains an internal adversary test: a formalised attempt to break its own strongest-looking claim. The discipline of this edition requires that the test be run honestly, which means reporting where the adversary \emph{wins}, not staging a defence that quietly declares victory. The most potent adversarial attack targets Limitation~L4 (\cref{lim:L4}): the origin of $\gamma = \GammaGeom$.

\paragraph{Adversarial argument.}
If $\gamma$ is a phenomenological fit parameter rather than a derived constant, then any vacuum-energy or cosmological calculation scaled by a power of $\gamma$ risks tautology: it would scale by precisely the amount needed to match observation, draining the result of independent explanatory power. The adversary presses further: the developmental record once decomposed $\gamma$ numerologically ($49/3 + 17/3000$, with the integer $16339$ prime), which is positive evidence that the value was \emph{constructed} to a target rather than found.

\paragraph{Honest adjudication.}
The framework does not claim the adversary is refuted. On the contrary: the numerological decomposition is quarantined in the museum (\cref{subsec:governance-boundary}) precisely because the adversary is right about it, that construction is not evidence and is barred from use. As for the broader charge, the defence is partial and is recorded as partial. $\gamma$ is fixed by matching to the internal kinetic vacuum structure of $\Sx$ (\cref{subsec:gamma-discipline}), independently of cosmology, so a vacuum-energy calculation using $\gamma$ is not circular \emph{with respect to the cosmological data}; but this does not make $\gamma$ derived. The honest verdict is that the adversarial argument \emph{stands} on its central point: $\gamma$ is \catmark{A-}, calibrated, not derived, and the framework's protection against tautology is not a first-principles derivation but the structural separation of strata (\cref{subsec:corrected-sequence}) together with the anti-target-leakage rule (\cref{prop:target-leakage}) that forbids any solver from reading the value it claims to predict.\catmark{A-}

\paragraph{What the self-audit establishes.}
The adversary test is not passed by winning the argument; it is passed by surviving the argument with the evidence classes intact. $\gamma$ remains \catmark{A-}; the museum construction stays refuted; the calibration is non-circular with respect to cosmology but is not a derivation. This is the appropriate outcome of an honest self-audit: the framework's weakest load-bearing point is identified, the adversary's valid points are conceded, and nothing is relabelled to make the defence look stronger than it is. The falsification architecture of \cref{sec:falsification-tension} and the evidence discipline of \cref{subsec:governance-boundary} exist for exactly this reason, to ensure that a self-audit can report a loss without that loss being quietly overwritten.\catmark{A}

\begin{governancebox}[Adversary Test: Outcome]
The strongest attack targets the origin of $\gamma$ (L4). The numerological construction is conceded to the adversary and stays quarantined. The calibration of $\gamma$ is non-circular with respect to cosmology but is \emph{not} a derivation; $\gamma$ remains \catmark{A-}. The test is ``passed'' only in the sense that the evidence classes survive the attack unchanged, not in the sense that the adversary is refuted.\catmark{A-}
\end{governancebox}

%=============================================================================
% APPENDIX A — INSERTION-POINT ANALYSIS
%=============================================================================
\newpage
\appendix
\part*{Appendix}
\addcontentsline{toc}{part}{Appendix}

\section{Endogenous Algebra Generation and the $G_{\mathrm{SM}}$ Obstruction Set}
\label{app:insertion-points}

\begin{governancebox}[Status of This Appendix]
This appendix is an advisory audit analysis. It sharpens the gauge-origin
obstruction of \cref{prop:d2-obstruction} and the open fork of
\cref{subsec:arm1-mechanisms} by relocating the obstruction from the exterior
calculus to the level of group representations and local operator algebras,
and by examining the two strongest endogenous enhancement mechanisms,
loop-homology phases and matrix backgrounds (\cref{app:endogenous}).
Every theorem-level statement here is a statement about the \emph{current
theorem corpus}, not about nature; no evidence class of the main text is
upgraded by this appendix, and the research form
$S(x)\to\mathcal{A}\to G_{\mathrm{SM}}$ remains \catmark{D}\ (ONT-08).
\end{governancebox}

\subsection{Primitive Data and Rigidity Results}
\label{app:rigidity}

\begin{definition}[Scalar primitive]
\label{def:scalar-primitive}
The primitive is the pair $(\Gamma, S)$, where $\Gamma=(V,E)$ is a locally
finite graph and $S:V\to\mathbb{R}$ is a single real scalar. Quantisation
assigns to each $v\in V$ the single-mode Hilbert space
$\mathcal{H}_v \cong L^2(\mathbb{R})$ with canonical pair
$(\hat S_v,\hat P_v)$, $[\hat S_v,\hat P_v]=i$, equivalently the Weyl algebra
$\mathcal{W}(\mathbb{R}^2,\sigma)$. No further structure is assumed.
\end{definition}

\begin{lemma}[Singlet Rigidity]
\label{lem:singlet-rigidity}
Let $G$ be a connected compact topological group and
$\rho:G\to GL(V)$ a continuous representation by linear field redefinitions
of a one-component real field, $V=\mathbb{R}$. Then $\rho$ is
trivial.\catmark{A}
\end{lemma}

\begin{proof}
$\rho(G)$ is a compact subgroup of $GL(1,\mathbb{R})=\mathbb{R}^{\times}$.
The maximal compact subgroup of $\mathbb{R}^{\times}$ is $\{\pm 1\}=O(1)$.
Connectedness of $G$ forces $\rho(G)\subseteq\{+1\}$.
\end{proof}

\begin{corollary}
\label{cor:singlet-charge}
Under linear actions, $S$ cannot carry any $G_{\mathrm{SM}}$ charge; it
cannot carry even a $\mathrm{U}(1)$ charge, since the nontrivial unitary
characters of $\mathrm{U}(1)$ act on $\mathbb{C}\cong\mathbb{R}^2$ and
require two real components. Any gauge sector to which $S$ couples must
therefore be supplied externally.\catmark{A}
\end{corollary}

\begin{proposition}[Discrete Exactness]
\label{prop:discrete-exactness}
Let link variables be functions of vertex differences,
$\theta_{uv}=S_u-S_v\in\mathbb{R}$, or, after target compactification
$S\sim S+2\pi$, $U_{uv}=e^{i(S_u-S_v)}\in\mathrm{U}(1)$. Then for every
cycle $\gamma=(v_1,\dots,v_n,v_1)$ in $\Gamma$,
\[
\sum_{i}\theta_{v_i v_{i+1}}=0,
\qquad
\prod_i U_{v_i v_{i+1}} = 1 .
\]
All holonomies telescope to the identity; the induced connection is pure
gauge with vanishing curvature on every plaquette.\catmark{A}
\end{proposition}

\begin{proof}
Telescoping sum around the cycle.
\end{proof}

\begin{remark}
\cref{prop:discrete-exactness} is the graph-theoretic form of the continuum
obstruction $F=\dd^2 S=0$ of \cref{prop:d2-obstruction} and of its
discrete-exterior-calculus extension recorded there. Note that target
compactification $\mathbb{R}\to S^1$ is itself an insertion: it modifies the
primitive of \cref{def:scalar-primitive}.\catmark{A}
\end{remark}

\begin{proposition}[Single-Mode Canonical Rigidity]
\label{prop:single-mode-rigidity}
The group of affine canonical (Bogoliubov) automorphisms of the single-mode
Weyl algebra $\mathcal{W}(\mathbb{R}^2,\sigma)$ is
$Sp(2,\mathbb{R})\ltimes\mathbb{R}^2$. Its maximal compact connected
subgroup is $SO(2)\cong\mathrm{U}(1)$. Consequently every connected compact
group acting on the mode by canonical linear transformations factors through
$\mathrm{U}(1)$.\catmark{A}
\end{proposition}

\begin{proof}
Linear canonical transformations of one mode preserve the symplectic form on
$\mathbb{R}^2$, giving $Sp(2,\mathbb{R})\cong SL(2,\mathbb{R})$; affine
shifts are the Heisenberg translations. The maximal compact subgroup of
$SL(2,\mathbb{R})$ is $SO(2)$ (Iwasawa decomposition); Stone--von Neumann
guarantees unitary implementability.
\end{proof}

\begin{remark}[Scope]
Non-quasi-free automorphisms of the Weyl $C^{*}$-algebra are not classified;
gauge-multiplet structure, however, requires linear action on the field,
which \cref{lem:singlet-rigidity} already excludes. The two results close
both ends.\catmark{A}
\end{remark}

\noindent\textbf{Summary.} From $(\Gamma,S)$ alone, the maximal continuous
compact symmetry content is abelian ($\mathrm{U}(1)$ at most, and only
globally); all scalar-built holonomies are flat; $S$ is
representation-theoretically inert. Any non-abelian structure must enter
through enlargement of the local data. The two canonical enlargement routes
are examined next, and their insertion points are localised
exactly.\catmark{A}

\subsection{Track 1: String-Net Condensation}
\label{app:track1}

String-net condensation~\cite{levinwen2005} takes as input a unitary fusion
category (UFC) $\mathcal{C}$ with $N=|\mathrm{Irr}(\mathcal{C})|<\infty$
simple objects, fusion multiplicities $N_{ab}^{c}$, quantum dimensions
$d_a$, and pentagon-solving $F$-symbols. The minimal local degrees of
freedom are qudits $\mathbb{C}^{N}$ on the edges of a trivalent graph, and
the exactly solvable commuting-projector Hamiltonian
$H_{\mathrm{LW}}=-\sum_v Q_v-\sum_p B_p$ is constructed \emph{from} the data
of $\mathcal{C}$; the emergent excitation theory is the Drinfeld centre
$Z(\mathcal{C})$~\cite{mueger2003,levinwen2005}. For
$\mathcal{C}=\mathrm{Vec}_G$ with $G$ finite, the emergent theory is the
finite-group quantum double~\cite{kitaev2003}. Two insertion points are
already visible: \textbf{(I1)}~the truncation
$L^2(\mathbb{R})\to\mathbb{C}^{N}$ per edge, which discards the canonical
pair and for which no canonical map exists; \textbf{(I2)}~the fusion data
$(\mathrm{Irr},N_{ab}^{c},F)$ themselves, for which no established theorem
produces a specific UFC from the dynamics of a single unconstrained bosonic
mode.\catmark{A}

\begin{proposition}[Finiteness obstruction]
\label{prop:ufc-finiteness}
For every UFC $\mathcal{C}$, the Drinfeld centre $Z(\mathcal{C})$ is a
modular tensor category with finitely many isomorphism classes of simple
objects~\cite{mueger2003,eno2005}. The representation category of any
compact connected non-abelian Lie group (in particular $\mathrm{SU}(2)$,
$\mathrm{SU}(3)$) has countably infinitely many simples. Hence exact
emergent continuous non-abelian gauge symmetry is outside the string-net
class.\catmark{A}
\end{proposition}

\begin{remark}
Root-of-unity truncations $\mathrm{SU}(N)_k$ are UFCs; the limit
$k\to\infty$ exits the UFC class and is a limit of a \emph{family} of
Hamiltonians, not a fixed-point Hamiltonian, hence outside the hypotheses of
the construction.\catmark{A/B}
\end{remark}

\begin{theorem}[$(3{+}1)$D classification, under stated hypotheses]
\label{thm:3p1-classification}
Gapped $(3{+}1)$D bosonic topological orders with finite-dimensional local
Hilbert spaces are classified by Dijkgraaf--Witten theories of finite groups
in the all-boson case~\cite{lankongwen2018}, and by finite-2-group gauge
theories when emergent fermions are present~\cite{lanwen2019}. Continuous
Yang--Mills structure therefore cannot arise as the gauge content of a
gapped $(3{+}1)$D bosonic topological order; independently, $(3{+}1)$D
Yang--Mills possesses local propagating degrees of freedom and is not a
topological phase, so the target lies outside the output class of the
formalism altogether.\catmark{A}
\end{theorem}

\begin{proposition}[Chirality obstruction]
\label{prop:chirality-obstruction}
Any local, translation-invariant, Hermitian lattice realisation with an
exactly implemented chiral charge exhibits fermion
doubling~\cite{nielsen1981}. Anomaly-free \emph{abelian} chiral lattice
gauge theories admit a construction~\cite{luscher1999}; no established
construction exists for non-abelian chiral gauge theories with the
Standard-Model fermion content, and symmetric-mass-generation routes remain
conjectural. Chirality is therefore a third independent insertion/open point
\textbf{(I3)}.\catmark{A/B}
\end{proposition}

\noindent\textbf{Track-1 verdict.} String-net condensation is a rigorous
template for the emergence of gauge structure \emph{as such}, but the
algebraic input is inserted (I1, I2), the output gauge content is finite
(\cref{prop:ufc-finiteness}, \cref{thm:3p1-classification}), and chirality
is obstructed (I3). Quantum-link/D-theory constructions reach continuous
groups only by inserting the group's representation theory into the link
Hilbert spaces~\cite{brower1999,wiese2022}, that is, by enlarging the
primitive, consistent with their status in
\cref{subsec:arm1-mechanisms} as benchmarks for $\mathcal{A}$, not
closures.\catmark{A}

\subsection{Track 2: Noncommutative Geometry}
\label{app:track2}

In an almost-commutative real even spectral triple
$(\mathcal{A},\mathcal{H},D;J,\gamma)$ with
$\mathcal{A}=C^\infty(M)\otimes\mathcal{A}_F$, gauge bosons arise from inner
fluctuations $D\mapsto D+A+\varepsilon' JAJ^{-1}$ with
$A=\sum_i a_i[D,b_i]=A^{*}$, and the gauge group is the image of
$\mathcal{U}(\mathcal{A})$ under the adjoint action
$\mathrm{Ad}(u)=uJuJ^{-1}$~\cite{chamseddine1997,connes2013}.

\begin{proposition}[Commutative scalar-generated algebra implies trivial gauge sector]
\label{prop:commutative-trivial}
Let $\mathcal{A}$ be commutative; in particular, let $\mathcal{A}$ be any
algebra generated by a single real scalar, so
$\mathcal{A}\subseteq C^\infty(M)$. Then
$\mathrm{Inn}(\mathcal{A})=\mathcal{U}(\mathcal{A})/\mathcal{U}(Z(\mathcal{A}))=\{1\}$,
and for unitary $u$ the adjoint action $uJuJ^{-1}$ multiplies the fermionic
Hilbert space by $u\bar{u}=\lvert u\rvert^{2}=1$; the fluctuation
combination $A+\varepsilon' JAJ^{-1}$ carries no gauge coupling. No gauge
bosons exist.\catmark{A}
\end{proposition}

\begin{proof}
$Z(\mathcal{A})=\mathcal{A}$ for commutative $\mathcal{A}$; $J$ implements
the opposite action, so $uJuJ^{-1}$ acts by $u\bar{u}=1$.
\end{proof}

\begin{remark}
\cref{prop:commutative-trivial} is the noncommutative-geometric analogue of
the $\dd^{2}=0$ obstruction: a scalar-generated, hence commutative, algebra
cannot source a gauge sector. \textbf{(I4)}~Noncommutativity of
$\mathcal{A}_F$ must be inserted by hand.\catmark{A}
\end{remark}

\noindent\textbf{Native slots for $S$.} A real singlet scalar has exactly
two established native slots in this formalism, both \emph{posterior} to the
choice of $\mathcal{A}_F$: the dilaton slot
$D\mapsto e^{-S/2}De^{-S/2}$~\cite{cc2006}, and the singlet
$\sigma$-slot promoting the Majorana mass entry,
$M_R\mapsto\sigma M_R$~\cite{cc2012,ccvs2013}. In neither slot does $S$ act
on $\mathcal{U}(\mathcal{A}_F)$ or constrain $\mathcal{A}_F$: the scalar is
selection-inert, consistent with \cref{cor:singlet-charge}.\catmark{A}

\begin{theorem}[Chamseddine--Connes classification; conditional]
\label{thm:cc-classification}
Under the hypotheses of (a)~irreducibility, (b)~KO-dimension~6,
(c)~nontrivial grading, and (d)~the symplectic hypothesis, the finite
algebra of a real spectral triple is constrained to
$\mathcal{A}_F\cong M_{k/2}(\mathbb{H})\oplus M_k(\mathbb{C})$ with $k$
even and minimal nontrivial solution $k=4$; imposing additionally (e)~the
order-one condition for a Dirac operator with non-vanishing Majorana sector
selects $\mathcal{A}_F\to\mathbb{C}\oplus\mathbb{H}\oplus M_3(\mathbb{C})$,
whose unimodular unitary group yields $G_{\mathrm{SM}}$ up to a finite
kernel~\cite{cc2008}.\catmark{A}\ (conditional)
\end{theorem}

\begin{proposition}[Status of the axioms]
\label{prop:cc-axiom-status}
Each hypothesis (a)--(e) of \cref{thm:cc-classification} is an external
postulate: KO-dimension~6 is chosen post hoc to permit Majorana masses and
resolve fermion doubling; relaxing the order-one condition enlarges the
result to the Pati--Salam algebra
$\mathbb{H}_R\oplus\mathbb{H}_L\oplus M_4(\mathbb{C})$~\cite{ccvs2013,farnsworth2015};
$k=4$ is minimality, not necessity; the generation number is inserted via
$\mathcal{H}_F\to\mathcal{H}_F^{\oplus 3}$. Consequently
$\mathbb{C}\oplus\mathbb{H}\oplus M_3(\mathbb{C})$ does not emerge uniquely
and unavoidably; it is the distinguished solution of an externally selected
axiom system, and the scalar $S$ participates in none of the selecting
conditions. Insertion points: \textbf{(I5)}~the axiom set (a)--(e),
\textbf{(I6)}~minimality $k=4$, \textbf{(I7)}~generations and
$\dim\mathcal{H}_F$.\catmark{A}
\end{proposition}

\subsection{Endogenous Enhancement Attempts: Loop Homology and Matrix Backgrounds}
\label{app:endogenous}

The insertion points I1--I7 localise where the two canonical routes import
structure. A natural rejoinder is that the importation might be avoidable:
that the relational graph itself, or a noncommutative substrate, might
generate the missing algebra \emph{endogenously}. This subsection examines
the two strongest such candidate mechanisms and records, at theorem level,
where each fails. Throughout, the success conditions are
\textbf{(C1)}~endogenously generated non-flat curvature,
\textbf{(C2)}~theorem-determined selection of
$\mathbb{C}\oplus\mathbb{H}\oplus M_3(\mathbb{C})$ with no free discrete
moduli, and \textbf{(C3)}~chiral matter.\catmark{A}

\paragraph{Loop-homology phases.}
Let $C^1(\Gamma;\mathbb{R})$ be the edge cochains of a finite connected
graph and $\dd_{0}:C^0\to C^1$ the coboundary. The combinatorial Hodge
decomposition $C^1=\operatorname{im}\dd_{0}\oplus\ker\dd_{0}^{*}$ splits
any edge datum into a gradient part and a harmonic (cycle) part, with
$\dim\ker\dd_{0}^{*}=b_1=|E|-|V|+1$; for the complete graph $K_4$ (the
tetrahedral configuration), $b_1=6-4+1=3$. Two readings of ``phases on
loops'' must be separated \emph{before} any algebra is computed. Under
reading~(a), loop phases are functionals of $S$: the edge datum is exact,
its harmonic component vanishes identically, and every loop phase is zero
by \cref{prop:discrete-exactness}; the mechanism terminates. Under
reading~(b), loop phases are \emph{independent} degrees of freedom valued
in the harmonic sector: this is an insertion, denoted \textbf{(I1$'$)},
and it is exactly the insertion that defines compact lattice gauge
theory~\cite{wilson1974}. Condition (C1) is then satisfied by the
insertion, not by $S$. The analysis proceeds under reading~(b) as the
maximal-charity hypothesis.\catmark{A}

\begin{proposition}[Structure of the loop-phase sector]
\label{prop:loop-phase-structure}
Under reading~(b), the gauge-invariant configuration space of flat-sector
phases is
\[
\mathcal{M}
=\mathrm{Hom}\big(\pi_1(\Gamma),\mathrm{U}(1)\big)
\cong\mathrm{Hom}\big(H_1(\Gamma;\mathbb{Z}),\mathrm{U}(1)\big)
\cong T^{\,b_1},
\]
and the infinitesimal transformation algebra is
$\mathfrak{h}=H^1(\Gamma;\mathbb{R})\cong\mathbb{R}^{b_1}$ with
$[\mathfrak{h},\mathfrak{h}]=0$. Adjoining canonical winding charges
extends $\mathfrak{h}$ to $b_1$ mutually commuting Heisenberg algebras;
the maximal compact connected symmetry remains the torus
$T^{b_1}$.\catmark{A}
\end{proposition}

\begin{proof}
The target group is forced to $\mathrm{U}(1)$ by
\cref{prop:single-mode-rigidity}. Homomorphisms from
$\pi_1(\Gamma)\cong F_{b_1}$ (free) into an abelian target factor through
the abelianisation $H_1\cong\mathbb{Z}^{b_1}$ (Hurewicz). The bracket on
$\mathfrak{u}(1)^{\oplus b_1}$ vanishes by definition of a torus.
\end{proof}

\begin{theorem}[Abelianisation Barrier]
\label{thm:abelianisation-barrier}
Any functorial assignment of phase data to the loops of a relational graph
factors through $H_1(\Gamma)$, the abelianisation of $\pi_1(\Gamma)$.
Consequently the bracket structure of any loop-generated transformation
algebra is identically zero: non-abelian Lie algebras cannot be generated
by graph homology. Non-abelian holonomy requires a non-abelian structure
group on the fibre, that is, local internal multiplicity at least two,
which is the insertion already localised as I4.\catmark{A}
\end{theorem}

\begin{proposition}[Dimension-counting fallacy at $K_4$]
\label{prop:k4-fallacy}
$H_1(K_4;\mathbb{R})\cong\mathbb{R}^3$ as a vector space and
$\dim\mathfrak{su}(2)=3$; nevertheless
$\mathfrak{u}(1)^{\oplus 3}\not\cong\mathfrak{su}(2)$, since a Lie-algebra
isomorphism preserves the bracket, $\mathfrak{su}(2)$ is simple, and
$[\mathfrak{u}(1)^{\oplus 3},\mathfrak{u}(1)^{\oplus 3}]=0$. The
coincidence of dimensions carries no algebraic content; the inference
``three independent loops imply $\mathfrak{su}(2)$'' is recorded as a
refuted route.\catmark{A}
\end{proposition}

\begin{remark}[Free fundamental groups select nothing]
\label{rem:universal-donor}
$\pi_1(K_4)\cong F_3$ is non-abelian, but free groups are universal
donors: every compact connected semisimple Lie group contains a dense
two-generator free subgroup~\cite{kuranishi1951}, so
$\mathrm{Hom}(F_3,G)$ is non-degenerate for essentially every candidate
$G$. A structure that admits all targets selects none; the
non-abelianness, if wanted, must again be supplied by the fibre. The one
non-abelian object the graph does supply,
$\mathrm{Aut}(K_4)=S_4\subset SO(3)$, is finite, and gauging it yields a
Dijkgraaf--Witten theory, consistent with
\cref{thm:3p1-classification}.\catmark{A}
\end{remark}

\begin{proposition}[Frenkel--Kac enhancement: the unique endogenous positive datum]
\label{prop:frenkel-kac}
The compact free boson in $1{+}1$ dimensions at the self-dual radius
carries the level-one current algebra $\widehat{\mathfrak{su}}(2)_1$,
generated by $J^3=\tfrac{i}{\sqrt{2}}\,\partial\varphi$ together with the
weight-one vertex operators
$J^{\pm}={:}\,e^{\pm i\sqrt{2}\,\varphi}\,{:}$~\cite{frenkelkac1980,goddardolive1986}.
This is the only theorem-level mechanism by which a \emph{single} scalar
generates a non-abelian algebra. Its preconditions are: (i)~target
compactification $S\sim S+2\pi R$, an insertion modifying the primitive of
\cref{def:scalar-primitive}; (ii)~tuning to the self-dual radius, a
codimension-one enhanced-symmetry point; (iii)~two-dimensional conformal
structure, with no $(3{+}1)$-dimensional analogue.\catmark{A}\ (within its
class)
\end{proposition}

\begin{lemma}[Rank Budget]
\label{lem:rank-budget}
Level-one simply-laced current algebras $\widehat{\mathfrak{g}}_1$ require
$\mathrm{rank}(\mathfrak{g})$ free bosons compactified on the root lattice
of $\mathfrak{g}$. A single scalar therefore yields at most rank one
($A_1$), whereas
$\mathrm{rank}(\mathfrak{g}_{\mathrm{SM}})=2+1+1=4$. Even granting the
two-dimensional enhancement mechanism in a four-dimensional problem, the
scalar primitive is short by three ranks.\catmark{A}\ (within the
Frenkel--Kac class)
\end{lemma}

\paragraph{Matrix backgrounds.}
The complementary candidate replaces the substrate itself: coordinates are
taken as non-commuting Hermitian matrices, as on the fuzzy
sphere~\cite{madore1992} or in matrix models of IKKT
type~\cite{ikkt1997}.

\begin{lemma}[Primitive Enlargement]
\label{lem:primitive-enlargement}
A ``real scalar field'' on the matrix algebra $M_N(\mathbb{C})$ is a
Hermitian matrix $\Phi$ with $N^2$ real components. The passage
$S\mapsto\Phi$ enlarges the on-site representation from dimension one to
dimension $N^2$ before any dynamics, and non-vanishing field strength
$F\sim[X_i,X_j]+\cdots\neq 0$ is purchased by the inserted
noncommutativity, not generated by $S$. Condition (C1) is satisfied by the
insertion I4; the obstruction set of \cref{app:rigidity} is evaded exactly
by abandoning the single-component hypothesis.\catmark{A}
\end{lemma}

\noindent Gauge structure on such backgrounds arises from background
splitting: for $X_i=L_i\otimes\mathbf{1}_n+A_i$, the fluctuation theory is
a $\mathrm{U}(n)$ gauge theory, where $n$ is the multiplicity of the
background block, an integer modulus; block backgrounds form dynamically
in matrix models with the corresponding
terms~\cite{grosse2005,azuma2004,steinacker2010}.\catmark{B}

\begin{proposition}[Automorphism Shadow]
\label{prop:automorphism-shadow}
All automorphisms of $M_N(\mathbb{C})$ are inner (Skolem--Noether), so
$\mathrm{Aut}(M_N(\mathbb{C}))\cong PU(N)$. Any gauge group obtained from
matrix-substrate fluctuations is the (projective) unitary group of the
inserted algebra or of a background-selected subalgebra: the gauge
structure is the shadow of the insertion, never of $S$.\catmark{A}
\end{proposition}

\begin{proposition}[Wedderburn Moduli]
\label{prop:wedderburn-moduli}
(i)~The matrix dimension $N$ is unconstrained by any theorem; commutative
and large-$N$ limits exist for all $N$. (ii)~Spontaneous block formation
breaks $\mathrm{U}(N)\to\prod_i\mathrm{U}(n_i)$, but the partition $(n_i)$
depends on potential parameters and entropy, with no theorem selecting
$(3,2,1)$~\cite{azuma2004}. (iii)~All factors so produced are complex
unitary algebras $M_{n_i}(\mathbb{C})$; the quaternionic summand
$\mathbb{H}\subset M_2(\mathbb{C})$ is the fixed-point algebra of an
antilinear structure $J$ with $J^2=-1$, which the complex matrix substrate
does not supply: choosing $J$ and its KO-dimension is exactly the
real-structure axiom of \cref{thm:cc-classification}, external to the
substrate. (iv)~Finite real spectral triples form a moduli space
(Krajewski diagrams)~\cite{krajewski1998,paschkesitarz1998}: the
Wedderburn--Artin invariants $\{(n_i,\mathbb{K}_i)\}$ with
$\mathbb{K}_i\in\{\mathbb{R},\mathbb{C},\mathbb{H}\}$ remain free.
(v)~The scalar $\Phi$ transforms in the adjoint and is selection-inert
with respect to (i)--(iv). Consequently
$\mathbb{C}\oplus\mathbb{H}\oplus M_3(\mathbb{C})$ does not emerge
inherently from a matrix background; it is fixed only by imposing the
axiom system of \cref{thm:cc-classification}.\catmark{A}
\end{proposition}

\subsection{Synthesis}
\label{app:synthesis}

\begin{theorem}[Insertion Necessity; statement about the current theorem corpus]
\label{thm:insertion-necessity}
Let the primitive data be exactly $(\Gamma,S:V\to\mathbb{R})$ with
symmetries implemented by canonical linear transformations. Then:
(1)~every connected compact group acts trivially on $S$
(\cref{lem:singlet-rigidity}, \cref{cor:singlet-charge});
(2)~every holonomy constructible from $S$ alone is exactly flat
(\cref{prop:discrete-exactness});
(3)~the maximal compact connected canonical symmetry of the local mode is
$\mathrm{U}(1)$ (\cref{prop:single-mode-rigidity});
(4)~the string-net route requires inserted fusion data and terminates in
finite gauge structures without chirality (I1--I3);
(5)~the noncommutative-geometric route requires inserted noncommutativity
and an inserted axiom system, with $S$ selection-inert (I4--I7).
Therefore no gapless ab initio derivation $S(x)\to G_{\mathrm{SM}}$ exists
within established theorems. The earliest unavoidable manual insertion is
the replacement of the abelian single-mode local algebra by a noncommutative
local algebra, equivalently the enlargement of the one-dimensional real
on-site representation to representations of dimension at least two. Every
known route performs this replacement exogenously; no theorem performs it
endogenously from single-boson dynamics.\catmark{A}
\end{theorem}

\begin{corollary}[Consistency with the EFT posture]
\label{cor:eft-posture}
Under the current theorem corpus, the only admissible role for $S$ relative
to $G_{\mathrm{SM}}$ is that of a real singlet coupled to an externally
supplied gauge sector, which is precisely the position the main text adopts
(gauge sector coupled, not generated; open form
$S(x)\to\mathcal{A}\to G_{\mathrm{SM}}$ at \catmark{D}, ONT-08). Both
tracks bound the unknown algebra $\mathcal{A}$ from below: Track~1 by ``at
least a unitary fusion category with a chirality-capable extension'' (none
known), Track~2 by ``at least a noncommutative real algebra satisfying a
Chamseddine--Connes-type axiom system''.\catmark{A}
\end{corollary}

\begin{theorem}[Gap Localisation; statement about the current theorem corpus]
\label{thm:gap-localisation}
Neither enhancement mechanism of \cref{app:endogenous} realises
$S(x)\to\mathcal{A}\to G_{\mathrm{SM}}$ ab initio. The missing
mathematical object is one and the same in both tracks: a theorem of
\emph{dynamical noncommutativity generation with fixed invariants}, that
is, an endogenous map from abelian single-mode dynamics on $\Gamma$ to a
noncommutative finite real $C^{*}$-algebra $\mathcal{A}$ satisfying
simultaneously
\textbf{(G1)}~fibre-multiplicity generation $1\to n\ge 2$ without inserted
variables, blocked homologically by \cref{thm:abelianisation-barrier} and
algebraically by \cref{prop:commutative-trivial};
\textbf{(G2)}~theorem-determined Wedderburn invariants
$\{(1,\mathbb{C}),(1,\mathbb{H}),(3,\mathbb{C})\}$, currently free moduli
by \cref{prop:wedderburn-moduli};
\textbf{(G3)}~an endogenous real structure $J$ of KO-dimension~6,
currently a choice by \cref{prop:cc-axiom-status}; and
\textbf{(G4)}~chirality-capable matter, blocked by
\cref{prop:chirality-obstruction}. The unique established endogenous
enhancement (\cref{prop:frenkel-kac}) satisfies a weakened (G1) only in
$1{+}1$ dimensions, under compactification and radius tuning, and is
rank-bounded at one against a requirement of four
(\cref{lem:rank-budget}). No element of the established corpus satisfies
(G1)--(G4) jointly.\catmark{A}
\end{theorem}

\begin{corollary}[Quantitative lower bounds on $\mathcal{A}$]
\label{cor:algebra-lower-bounds}
The obstruction set converts the open form
$S(x)\to\mathcal{A}\to G_{\mathrm{SM}}$ (ONT-08, \catmark{D}) into
quantitative lower bounds on any admissible candidate: $\mathcal{A}$ must
carry (i)~local internal multiplicity at least two (non-abelian fibre);
(ii)~rank resources at least four; (iii)~a real structure with the
$J^2$-signs of KO-dimension~6; and (iv)~a chirality mechanism evading the
obstruction set of \cref{prop:chirality-obstruction}. Any candidate
failing one bound is excluded before dynamical analysis; the ledger fields
\texttt{SCALAR\_ONLY\_COMPATIBLE} and \texttt{GROUP\_SELECTION\_MECHANISM}
of \cref{subsec:arm1-mechanisms} remain mandatory for every
candidate.\catmark{A/D}
\end{corollary}

\begin{governancebox}[Forbidden Inferences]
Mechanism relevance is not framework validation: neither the string-net
construction~\cite{levinwen2005} nor the spectral
classification~\cite{cc2008} supports any claim that \UIDT\ derives
$G_{\mathrm{SM}}$. Conversely, this appendix is a gap localisation, not an
impossibility theorem for all conceivable enriched routes; no such theorem
exists. \Cref{prop:frenkel-kac} does not license the claim that a scalar
generates $\mathrm{SU}(2)$ in $3{+}1$ dimensions, and
\cref{prop:wedderburn-moduli} does not license the claim that matrix
models derive the Standard Model. $G_{\mathrm{SM}}$ is, on all examined
routes, selected, not forced.
\end{governancebox}

\subsection{Appendix Claims Register}
\label{app:claims-register}

\begin{longtable}{p{0.08\textwidth} p{0.42\textwidth} p{0.12\textwidth} p{0.28\textwidth}}
\caption{Appendix claims register. Classes follow the schema of \cref{subsec:governance-boundary}; falsification or downgrade conditions are stated per claim.}\\
\toprule
\textbf{ID} & \textbf{Claim} & \textbf{Class} & \textbf{Falsification / downgrade} \\
\midrule
\endfirsthead
\toprule
\textbf{ID} & \textbf{Claim} & \textbf{Class} & \textbf{Falsification / downgrade} \\
\midrule
\endhead
APP-01 & Singlet rigidity (\cref{lem:singlet-rigidity}, \cref{cor:singlet-charge}). & \catmark{A} & Nontrivial continuous one-dimensional real representation of a connected compact group (impossible). \\
APP-02 & Discrete exactness (\cref{prop:discrete-exactness}). & \catmark{A} & Cycle with non-telescoping holonomy within the construction (impossible). \\
APP-03 & Single-mode rigidity, maximal compact $\mathrm{U}(1)$ (\cref{prop:single-mode-rigidity}). & \catmark{A} & Compact connected $K\not\subseteq\mathrm{U}(1)$ inside $SL(2,\mathbb{R})$ (impossible). \\
APP-04 & UFC finiteness excludes exact emergent compact connected non-abelian Lie gauge groups (\cref{prop:ufc-finiteness}). & \catmark{A} & A UFC with infinitely many simples (contradiction in terms). \\
APP-05 & Gapped $(3{+}1)$D bosonic orders carry finite gauge content (\cref{thm:3p1-classification}). & \catmark{A}\ under hypotheses & An extension of the classification admitting continuous gauge content would narrow the scope. \\
APP-06 & No established non-abelian chiral lattice gauge construction with SM fermion content. & \catmark{B} & A rigorous construction (e.g.\ a completed symmetric-mass-generation programme) resolves I3. \\
APP-07 & Commutative $\mathcal{A}$ yields a trivial gauge sector (\cref{prop:commutative-trivial}). & \catmark{A} & None within the stated noncommutative-geometry axioms. \\
APP-08 & $\mathbb{C}\oplus\mathbb{H}\oplus M_3(\mathbb{C})$ is conditional on the Chamseddine--Connes axioms; $S$ is selection-inert (\cref{thm:cc-classification}, \cref{prop:cc-axiom-status}). & \catmark{A}\ conditional & Derivation of the axioms from an independent principle would strengthen; known relaxations (Pati--Salam) confirm non-uniqueness. \\
APP-09 & Insertion necessity (\cref{thm:insertion-necessity}). & \catmark{A}\ re current corpus & A theorem generating a UFC or a noncommutative $\mathcal{A}_F$ endogenously from single-boson data. \\
APP-10 & $S(x)\to\mathcal{A}\to G_{\mathrm{SM}}$ as research programme (ONT-08). & \catmark{D} & Remains \catmark{D}\ until $\mathcal{A}$ is explicitly defined and externally stress-tested. \\
APP-11 & Abelianisation Barrier: loop-generated transformation algebras have identically vanishing bracket (\cref{thm:abelianisation-barrier}, \cref{prop:loop-phase-structure}). & \catmark{A} & A functorial non-abelian bracket on graph homology (impossible: $H_1=\pi_1^{\mathrm{ab}}$). \\
APP-12 & The $K_4$ dimension coincidence with $\mathfrak{su}(2)$ is algebraically empty (\cref{prop:k4-fallacy}). & \catmark{A} & None within the construction (bracket obstruction). \\
APP-13 & Free fundamental groups are universal donors and select no target group (\cref{rem:universal-donor}). & \catmark{A} & A theorem restricting $\mathrm{Hom}(F_3,G)$ to a unique compact $G$. \\
APP-14 & Frenkel--Kac enhancement is the unique endogenous datum; two-dimensional, tuned, rank-bounded at one (\cref{prop:frenkel-kac}, \cref{lem:rank-budget}). & \catmark{A}\ within its class & A $(3{+}1)$-dimensional analogue or a rank-amplification mechanism would revise the budget. \\
APP-15 & Matrix backgrounds: curvature and multiplicity are purchased by insertion; the gauge group is the automorphism shadow $PU(N)$ (\cref{lem:primitive-enlargement}, \cref{prop:automorphism-shadow}). & \catmark{A} & None within the construction (Skolem--Noether). \\
APP-16 & Wedderburn invariants are free moduli; $\mathbb{C}\oplus\mathbb{H}\oplus M_3(\mathbb{C})$ is not selected by a matrix substrate (\cref{prop:wedderburn-moduli}). & \catmark{A}/\catmark{B} & A dynamical theorem fixing the Wedderburn invariants and the real structure $J$ endogenously. \\
APP-17 & Gap localisation (G1)--(G4) and quantitative lower bounds on $\mathcal{A}$ (\cref{thm:gap-localisation}, \cref{cor:algebra-lower-bounds}). & \catmark{A}\ re current corpus / \catmark{D} & Discovery of a mechanism satisfying (G1)--(G4) jointly. \\
\bottomrule
\end{longtable}
\newpage
\paragraph{Appendix reproduction note.}
Following the template of \cref{subsec:evidence-grammar}, the
machine-checkable identities of this appendix are accompanied by
executable commands. The telescoping identity of \cref{prop:discrete-exactness} was
verified in exact rational arithmetic over one thousand randomly generated
cycles of lengths 3--40, with residual exactly zero in every case:
\begin{verbatim}
python3 -c "from fractions import Fraction; import random; random.seed(7);
ok=all(sum(S[i]-S[(i+1)%n] for i in range(n))==0
 for n,S in ((m,[Fraction(random.randint(-10**6,10**6),
 random.randint(1,10**3)) for _ in range(m)])
 for m in (random.randint(3,40) for _ in range(1000)))); print(ok)"
# -> True (1000/1000 cycles; residual exactly 0; exact arithmetic,
#    so this certifies the identity itself, not a numerical approximation)
\end{verbatim}
The homology rank entering \cref{prop:loop-phase-structure} and
\cref{prop:k4-fallacy} is likewise machine-checkable in exact integer
arithmetic: the cycle rank of $K_4$ is
$b_1=|E|-\mathrm{rank}\,\partial=6-3=3$,
\begin{verbatim}
python3 -c "import sympy as sp; B=sp.Matrix([[1,-1,0,0],[1,0,-1,0],
[1,0,0,-1],[0,1,-1,0],[0,1,0,-1],[0,0,1,-1]]); print(B.shape[0]-B.rank())"
# -> 3 (exact rank over the rationals; the vanishing bracket of
#    u(1)^3 is definitional and requires no computation)
\end{verbatim}
The remaining \catmark{A}-level statements of this appendix are
group-theoretic and operator-algebraic facts (maximal compact subgroups of
$\mathbb{R}^{\times}$ and $SL(2,\mathbb{R})$; Stone--von Neumann; finiteness
of simples in a unitary fusion category; the Hurewicz factorisation through
$H_1$ underlying \cref{thm:abelianisation-barrier}; the Frenkel--Kac
construction; Skolem--Noether): standard results, cited rather
than machine-verified. \cref{thm:3p1-classification},
\cref{thm:cc-classification}, and \cref{prop:frenkel-kac} hold under their
stated hypotheses; their scope, not their truth, is the open
variable.\catmark{A}

%=============================================================================
% APPENDIX B — CONNES-VACUUM ATTRACTOR PROGRAMME (PATCH 005; advisory only)
%=============================================================================
\newpage
\section{Dynamical Attractors toward the Connes Vacuum: Established Inputs, Programme Requirements, and the Residual Freezing Gap}
\label{app:connes-attractor}

\begin{governancebox}[Status of This Appendix]
This appendix is an advisory audit analysis and a \emph{programme charter},
not a result. It does not establish, and must not be cited as establishing,
a thermodynamic emergence of the algebra
$\mathbb{C}\oplus\mathbb{H}\oplus M_3(\mathbb{C})$, a dynamical necessity of
KO-dimension~6, or an anomaly-theoretic uniqueness of the $(1,2,3)$
partition. Established external inputs are carried at \catmark{B};
conditional selections at \catmark{B}\ under stated hypotheses; the
attractor content itself as programme requirements at \catmark{D}\ or
conjectures at \catmark{E}. The appendix operates under the gap
localisation of \cref{thm:gap-localisation}, the lower bounds of
\cref{cor:algebra-lower-bounds}, and the target-leakage theorem of
\cref{prop:target-leakage}. An AI-generated research memo in the
developmental record proposes the mechanisms surveyed here; per the corpus
discipline of \cref{subsec:corpus-topology} (native extraction, verified
identifiers), that memo is contextual \catmark{E}\ material and is not
cited as evidence. The research form
$S(x)\to\mathcal{A}\to G_{\mathrm{SM}}$ remains \catmark{D}\ (ONT-08).
\end{governancebox}

\subsection{From External Postulates to Dynamical Selection: Programme Statement}
\label{app:b-programme}

The classification of \cref{thm:cc-classification} reaches
$\mathbb{C}\oplus\mathbb{H}\oplus M_3(\mathbb{C})$ only through the
externally imposed axiom set localised as insertion point I5
(\cref{prop:cc-axiom-status}). The attractor programme asks whether each
postulate can be replaced by a dynamical or thermodynamic selection
principle acting on a larger configuration space, so that the Connes
vacuum would arise as a stable attractor rather than as an axiom. The
programme is stated here in falsification-disciplined form; none of its
attractor claims is established, and the statement of the programme is a
\catmark{A}-class fact about the corpus, not about nature.\catmark{A}

\begin{definition}[Connes-vacuum attractor programme]
\label{def:attractor-programme}
The programme assigns to each external postulate of
\cref{thm:cc-classification} a candidate dynamical replacement:
(a)~irreducibility, replaced by condensation of a dominant connected
sector; (b),(c)~the real structure, grading, and KO-signs, replaced by a
stability property of charge-conjugation-symmetric configurations
(\cref{app:b-ko6}); (d)~the Wedderburn data, replaced by free-energy
selection on the partition moduli of \cref{prop:wedderburn-moduli}
(\cref{app:b-condensation}); (e)~the order-one condition, replaced by the
unresolved freezing mechanism of \cref{app:b-freezing}. Success conditions
are (G1)--(G4) of \cref{thm:gap-localisation}; every candidate mechanism
must additionally satisfy the lower bounds of
\cref{cor:algebra-lower-bounds}.\catmark{D}
\end{definition}

\subsection{Matrix Condensation and the Wedderburn Moduli Problem}
\label{app:b-condensation}

\paragraph{What is established.}
Zero-dimensional matrix ensembles of IKKT type~\cite{ikkt1997} and their
bosonic deformations exhibit condensation transitions in which
block-diagonal backgrounds form dynamically, breaking
$\mathrm{U}(N)\to\prod_i\mathrm{U}(n_i)$; fuzzy-sphere backgrounds
nucleate and dissolve in first-order transitions governed by the
competition between the classical potential and the entropy of
fluctuations, and off-diagonal (bifundamental) modes acquire masses
growing with the separation of the condensed
blocks~\cite{azuma2004,grosse2005,steinacker2010}.\catmark{B} In
free-energy language, $F[(n_i)]=U[(n_i)]-T\,S[(n_i)]$ on the partition
moduli: entropy-dominated regimes favour either a single dominant block or
near-symmetric partitions, while potential-dominated regimes track the
discrete minima of the engineered potential. In no established regime is
the asymmetric partition $(1,2,3)$ singled out by a structural
principle.\catmark{B}

\begin{remark}[The moduli problem restated]
\label{rem:b-moduli-restated}
\cref{prop:wedderburn-moduli} remains binding at every temperature: the
condensed factors are complex unitary algebras $M_{n_i}(\mathbb{C})$; the
partition $(n_i)$ is a function of couplings, chemical potentials, and
entropy, not of a selection theorem; and the quaternionic summand
$\mathbb{H}$ requires the antilinear structure of \cref{app:b-ko6}, which
the complex ensemble does not supply. Entropy maximisation is therefore
structurally blind to the $(1,2,3)$ hierarchy, and any claim that the
Connes algebra is the unique global thermodynamic minimum of an
unconstrained ensemble is rejected.\catmark{A}
\end{remark}

\begin{governancebox}[Programme Requirement PR-B1: Partition Selection \catmark{D}]
Exhibit a concrete matrix action and a coupling/temperature window in
which the global free-energy minimiser over Wedderburn partitions is
$(1,2,3)$, with the minimiser stable as $N\to\infty$, under the
anti-target-leakage protocol of \cref{prop:target-leakage}: all parameters
fixed by observables independent of the partition; the partition read out
post hoc by a sealed evaluation; null controls run on scrambled reference
partitions (for example $(1,2,4)$ and $(2,2,2)$), which must \emph{not} be
selected by the same procedure. Downgrade condition: if the minimiser
tracks engineered parameters rather than a structural principle, the
mechanism is selection by hand and is recorded at \catmark{E}.
\end{governancebox}

\subsection{Anomaly Freedom, Unimodularity, and the Quiver Flow Conjecture}
\label{app:b-anomaly}

\paragraph{What is established.}
Within the noncommutative-geometric Standard Model the physical gauge
group is the unimodular unitary group of $\mathcal{A}_F$, and the
unimodularity condition is tied to the cancellation of the chiral gauge
anomalies of the spectral fermion content~\cite{agbm1995,cc2008}: anomaly
freedom prunes the unitary group of an \emph{already chosen} algebra to
$G_{\mathrm{SM}}$ up to a finite kernel; it does not select the
algebra.\catmark{B} The configuration space on which any selection
dynamics would act is likewise established: finite real spectral triples
are classified combinatorially by Krajewski
diagrams~\cite{krajewski1998,paschkesitarz1998}, and the order-one
condition acts there as a combinatorial admissibility constraint on the
diagram, restricting which vertex pairs the components of $D_F$ may
connect.\catmark{A}

\paragraph{What is not established.}
No renormalisation-group flow on the space of Krajewski diagrams is
currently defined. By the scaffold discipline of
\cref{subsec:gferg-scaffold}, a claimed flow must specify, in advance, the
flow equation, the regulator analogue on the discrete moduli, the
truncation, the boundary conditions, and the failure criteria; absent
these, renormalisation-group vocabulary on quiver space is not a
calculation. Moreover, uniqueness manifestly fails at the level of the
known classification: the moduli space of anomaly-consistent finite
triples contains families beyond the Standard-Model diagram, and the
relaxation of the order-one condition enlarges the solution set to
Pati--Salam~\cite{ccvs2013,farnsworth2015}. Accordingly the following is
recorded as a conjecture, not a theorem.\catmark{A}

\begin{remark}[Conjecture C-B1: anomaly-constrained quiver attractor, status \catmark{E}]
\label{rem:quiver-conjecture}
Within the anomaly-free, order-one-admissible stratum of Krajewski
diagrams there exists a well-defined coarse-graining flow whose unique
stable fixed point at three summands is the $(1,2,3)$ diagram. No such
flow is defined; no uniqueness argument exists; the conjecture inherits
the full burden of \cref{prop:target-leakage}, since a flow whose fixed
point is verified only against the known Standard-Model diagram predicts
nothing.\catmark{E}
\end{remark}

\subsection{Real Structure and KO-Dimension 6: A Conditional Selection}
\label{app:b-ko6}

\paragraph{What is established, conditionally.}
The KO-dimension of a real spectral triple is fixed modulo~8 by the sign
triple $(\varepsilon,\varepsilon',\varepsilon'')$ in $J^2=\varepsilon$,
$JD=\varepsilon'DJ$, $J\gamma=\varepsilon''\gamma J$; KO-dimension~6
carries $(\varepsilon,\varepsilon',\varepsilon'')=(1,1,-1)$. Within the
almost-commutative spectral Standard Model, imposing simultaneously
(i)~the elimination of the fermion quadrupling of the Euclidean
formulation and (ii)~the admissibility of Majorana mass terms for
right-handed neutrinos selects KO-dimension~6 for the finite
triple~\cite{barrett2007,connes2006,ccm2007}.\catmark{B} With
$\varepsilon''=-1$ the finite Dirac operator admits the block form
\[
D_F=\begin{pmatrix} S & T^{*}\\[2pt] T & \bar S\end{pmatrix},
\qquad T=T^{T},
\]
where $S$ contains the Dirac masses and the symmetric block $T$ couples
right-handed neutrinos to their conjugates, enabling the seesaw
mechanism~\cite{connes2006,ccm2007}.\catmark{B} The selection is
\emph{conditional}: requirements (i) and (ii) are physical inputs, not
consequences of any dynamics. Necessity of KO-dimension~6 is therefore
necessity relative to (i) and (ii), and nothing stronger may be asserted;
formulations such as ``absolute mathematical necessity'' or ``rigorously
proven not to be an insertion'' are rejected wherever they appear in the
developmental record.\catmark{A}

\paragraph{What is open.}
A dynamical origin of the real structure $J$ itself, for instance through
a charge-conjugation-stability or orientifold-type projection of a matrix
ensemble, is not established. The heuristic that unoriented configurations
are entropically or energetically disfavoured on a stochastic substrate is
exactly that, a heuristic: no defined ensemble, measure, or flow currently
demonstrates it, and \cref{prop:wedderburn-moduli}(iii) records that the
complex substrate does not supply $J$.\catmark{A}

\begin{governancebox}[Programme Requirement PR-B2: Endogenous Real Structure \catmark{D}]
Exhibit a substrate dynamics whose stable configurations are invariant
under an antilinear involution $J$ with the KO-dimension-6 sign triple,
where $J$-invariance emerges as a stability or extremality property of the
ensemble rather than as a restriction imposed on the configuration space
at the outset, and show that configurations without such $J$ are
dynamically disfavoured. Downgrade condition: if $J$-invariance is imposed
on the measure by hand (an orientifold projection inserted at step zero),
the mechanism is an insertion and PR-B2 is not met.
\end{governancebox}

\subsection{The Residual Freezing Gap and a Spectral Falsification Protocol}
\label{app:b-freezing}

\paragraph{Layer-C declaration.}
The deepest unresolved step of the programme is the dynamical origin of
the bipartite order-one condition \emph{prior to} any flow: the condition
is doing genuine selection work, since its relaxation enlarges the algebra
to Pati--Salam~\cite{ccvs2013,farnsworth2015}, yet no mechanism explains
why the commutant structure of an emergent geometry should freeze into the
first-order form at all. This is recorded as an open question, not as a
difficulty in the course of being resolved.\catmark{A}

\begin{openquestion}[The freezing gap]
\label{oq:freezing-gap}
What dynamical or thermodynamic principle, acting on the configuration
space of finite real spectral triples, enforces the order-one condition
$[[D_F,a],Jb^{*}J^{-1}]=0$ as a stability or extremality property, prior
to and independently of any imposed axiom? No candidate principle is
currently established; the programme requirements PR-B1 and PR-B2 are
necessary inputs to any answer but are not sufficient for it.\catmark{D}
\end{openquestion}

\paragraph{Falsification handle.}
If a concrete attractor realisation is ever exhibited, it acquires a
discriminating observable through the spectral statistics of the condensed
finite Dirac operator. Random Dirac ensembles over finite spectral triples
are an established, simulable arena~\cite{barrettglaser2016}, and with
$Z(\beta+it)=\operatorname{Tr}e^{-(\beta+it)D_F^{2}}$ the spectral form
factor $g(t;\beta)=\lvert Z(\beta+it)\rvert^{2}/Z(\beta)^{2}$
distinguishes random-matrix spectral rigidity (the dip-ramp-plateau
profile, with the symmetry class fixed by the sign triple) from Poissonian
level statistics~\cite{cotler2017}. The protocol below is a \catmark{D}\
proposal; no simulation or measurement of it exists for any \UIDT-relevant
ensemble.

\begin{killswitch}
If a concrete attractor realisation (PR-B1 together with PR-B2) is
exhibited and the spectral form factor of its condensed phase is
incompatible with the rigidity class fixed by that realisation's own sign
triple $(\varepsilon,\varepsilon',\varepsilon'')$, the realisation is
falsified. If, instead, stable condensed phases are exhibited whose
algebras lack $G_{\mathrm{SM}}$ entirely, any uniqueness reading of the
Connes vacuum as attractor is invalidated and Conjecture C-B1 is to be
withdrawn. Absent any realisation, ONT-08 remains \catmark{D}\ and no
evidence class of the main text is altered by this appendix.\catmark{D}
\end{killswitch}

\subsection{Appendix B Claims Register}
\label{app:b-claims-register}

\begin{longtable}{p{0.10\textwidth} p{0.40\textwidth} p{0.12\textwidth} p{0.28\textwidth}}
\caption{Appendix B claims register. Classes follow the schema of \cref{subsec:governance-boundary}; falsification or downgrade conditions are stated per claim.}\\
\toprule
\textbf{ID} & \textbf{Claim} & \textbf{Class} & \textbf{Falsification / downgrade} \\
\midrule
\endfirsthead
\toprule
\textbf{ID} & \textbf{Claim} & \textbf{Class} & \textbf{Falsification / downgrade} \\
\midrule
\endhead
APP-B-01 & Block condensation with entropy/potential competition and massive bifundamentals is established in matrix models (\cref{app:b-condensation}). & \catmark{B} & Source audit of the cited lattice/matrix literature; downgrade on failed audit. \\
APP-B-02 & Entropy maximisation is blind to $(1,2,3)$; no functional selects the partition (\cref{rem:b-moduli-restated}). & \catmark{A}\ re current corpus & A selection theorem satisfying PR-B1 under \cref{prop:target-leakage}. \\
APP-B-03 & Anomaly freedom is tied to unimodularity and prunes $\mathcal{U}(\mathcal{A}_F)$ to $G_{\mathrm{SM}}$; it does not select $\mathcal{A}_F$ (\cref{app:b-anomaly}). & \catmark{B} & Source audit; a theorem deriving the algebra from anomaly data alone would supersede. \\
APP-B-04 & The order-one condition is a combinatorial admissibility constraint on Krajewski diagrams; the anomaly-consistent stratum is non-unique (\cref{app:b-anomaly}). & \catmark{A} & Exhibition of uniqueness at three summands would revise. \\
APP-B-05 & Quiver-attractor uniqueness is Conjecture C-B1 (\cref{rem:quiver-conjecture}). & \catmark{E} & A defined flow with regulator, truncation, boundary conditions, and a leak-free uniqueness proof would upgrade. \\
APP-B-06 & KO-dimension 6 is selected conditionally by fermion-doubling resolution plus Majorana admissibility; necessity is relative to these inputs (\cref{app:b-ko6}). & \catmark{B}\ conditional & A dynamical derivation of inputs (i)/(ii) or of $J$ would strengthen; otherwise permanent. \\
APP-B-07 & The dynamical origin of the order-one condition is open; the freezing gap stands (\cref{oq:freezing-gap}). & \catmark{A}/\catmark{D} & A principle answering \cref{oq:freezing-gap} under PR-B1/PR-B2 discipline. \\
APP-B-08 & Spectral-form-factor protocol as discriminator for any future realisation (\cref{app:b-freezing}). & \catmark{D} & Executable only upon a concrete realisation; the killswitch criterion then applies. \\
\bottomrule
\end{longtable}

\paragraph{Appendix B reproduction note.}
The established input of \cref{app:b-anomaly} has a machine-checkable
core: the four anomaly sums of one Standard-Model generation vanish
identically in exact rational arithmetic for the unimodular hypercharge
assignment
$(Y_Q,Y_{u^c},Y_{d^c},Y_L,Y_{e^c})=(\tfrac{1}{6},-\tfrac{2}{3},\tfrac{1}{3},-\tfrac{1}{2},1)$:
\begin{verbatim}
python3 -c "from fractions import Fraction as F;
Y={'Q':(F(1,6),6),'uc':(F(-2,3),3),'dc':(F(1,3),3),'L':(F(-1,2),2),
'ec':(F(1),1)};
print(sum(m*y**3 for y,m in Y.values()),
 sum(m*y for y,m in Y.values()),
 2*Y['Q'][0]+Y['uc'][0]+Y['dc'][0],
 3*Y['Q'][0]+Y['L'][0])"
# -> 0 0 0 0  ([U(1)]^3, grav^2-U(1), [SU(3)]^2-U(1), [SU(2)]^2-U(1);
#    exact arithmetic; certifies the assignment, not any UIDT claim)
\end{verbatim}
This check certifies the anomaly arithmetic of the established input only;
no attractor statement of this appendix is machine-checkable, because no
attractor realisation exists to check.\catmark{A}

\begin{governancebox}[Forbidden Inferences]
Nothing in this appendix licenses the statements that the Connes vacuum
emerges thermodynamically, that KO-dimension~6 is dynamically necessary,
that anomaly cancellation selects the $(1,2,3)$ algebra, or that \UIDT\
derives $G_{\mathrm{SM}}$. The established inputs are external literature
results on matrix condensation, unimodularity, and conditional
KO-dimension selection; the attractor content consists of two programme
requirements (PR-B1, PR-B2), one conjecture (C-B1), and one open question
(\cref{oq:freezing-gap}). \Cref{thm:gap-localisation} and
\cref{cor:algebra-lower-bounds} remain the binding constraints, and
emergence wording is rejected wherever it appears in the developmental
record.
\end{governancebox}

%=============================================================================
% BIBLIOGRAPHY
%=============================================================================
\newpage
\begin{thebibliography}{9}
\footnotesize

\bibitem{uidt-ontology}
P.~Rietz, \textit{Vacuum Information Density as the Fundamental Geometric Scalar: UIDT Ontology v3.9} (this work). Zenodo. DOI: \href{https://doi.org/10.5281/zenodo.20319634}{10.5281/zenodo.20319634}.

\bibitem{uidt-framework}
P.~Rietz, \textit{UIDT Framework v3.9 (parent record)}. Zenodo. DOI: \href{https://doi.org/10.5281/zenodo.17835200}{10.5281/zenodo.17835200}.

\bibitem{uidt-clay}
P.~Rietz, \textit{UIDT Clay-Submission Record}. Zenodo. DOI: \href{https://doi.org/10.5281/zenodo.18003018}{10.5281/zenodo.18003018}. Historical provenance record; cited as documentation only, not as an active claim.

\bibitem{bib-audit-note}
\textbf{Bibliography audit note.} The external references below were resolved to verifiable primary sources (DOI or arXiv) in a dedicated literature audit. In keeping with the framework's evidence discipline, citation of these sources changes no \UIDT\ evidence class: external literature supports the quantitative pure-gauge benchmarks, is contested on the interpretive physics, explicitly identifies open chiral-fermion obstructions, is interpretive for the consciousness frameworks, and refutes any MeV-to-neural coupling. A small number of secondary venue/page details remain flagged for final primary confirmation before submission.

% Cluster 1: pure SU(3) lattice benchmarks
\bibitem{athenodorou2020}
A.~Athenodorou, M.~Teper, ``The glueball spectrum of SU(3) gauge theory in 3+1 dimensions,'' JHEP \textbf{11} (2020) 172. arXiv:2007.06422.
\bibitem{athenodorou2021}
A.~Athenodorou, M.~Teper, ``SU(N) gauge theories in 3+1 dimensions: glueball spectrum, string tensions and topology,'' JHEP \textbf{12} (2021) 082. arXiv:2106.00364.
\bibitem{morningstar1999}
C.~J.~Morningstar, M.~Peardon, ``The glueball spectrum from an anisotropic lattice study,'' Phys.\ Rev.\ D \textbf{60} (1999) 034509. arXiv:hep-lat/9901004.
\bibitem{morningstar2024}
C.~Morningstar, ``Update on Glueballs,'' PoS \textbf{LATTICE2024} (2024) 004. DOI 10.22323/1.466.0004. arXiv:2502.02547.
\bibitem{michael1989}
C.~Michael, M.~Teper, ``The glueball spectrum in SU(3),'' Nucl.\ Phys.\ B \textbf{314} (1989) 347.
\bibitem{luscher2010}
M.~L\"uscher, ``Properties and uses of the Wilson flow in lattice QCD,'' JHEP \textbf{08} (2010) 071. arXiv:1006.4518.
\bibitem{borsanyi2012}
S.~Bors\'anyi et al.\ (BMW), ``High-precision scale setting in lattice QCD,'' JHEP \textbf{09} (2012) 010. arXiv:1203.4469.
\bibitem{chitop2025}
S.~D\"urr, G.~Fuwa, ``Topological susceptibility and excess kurtosis in SU(3) Yang--Mills theory,'' Phys.\ Rev.\ D \textbf{113} (2026) 054508. arXiv:2501.08217.
\bibitem{durr2007}
S.~D\"urr, Z.~Fodor, C.~Hoelbling, T.~Kurth, ``Precision study of the SU(3) topological susceptibility in the continuum,'' arXiv:hep-lat/0612021.
\bibitem{dallabrida2019}
M.~Dalla Brida, A.~Ramos, ``The gradient flow coupling at high-energy and the scale of SU(3) Yang--Mills theory,'' Eur.\ Phys.\ J.\ C \textbf{79} (2019) 720. arXiv:1905.05147.

% Cluster 1b: non-Abelian Casimir effect on the lattice
\bibitem{chernodub2023}
M.~N.~Chernodub, V.~A.~Goy, A.~V.~Molochkov, A.~S.~Tanashkin, ``Boundary states and non-Abelian Casimir effect in lattice Yang--Mills theory,'' Phys.\ Rev.\ D \textbf{108} (2023) 014515. arXiv:2302.00376.
\bibitem{ngwenya2025}
B.~A.~Ngwenya, A.~K.~Rothkopf, W.~A.~Horowitz, ``The Non-Abelian Casimir Effect for Plates, Symmetrical Tube and Box on the Lattice,'' arXiv:2507.21333 (2025); B.~A.~Ngwenya, ``The Casimir effect in non-Abelian gauge theories on the lattice,'' Ph.D.\ thesis, University of Cape Town (2025), arXiv:2505.02875.

% Cluster 2: noncommutative geometry SM
\bibitem{chamseddine1997}
A.~H.~Chamseddine, A.~Connes, ``The Spectral Action Principle,'' Commun.\ Math.\ Phys.\ \textbf{186} (1997) 731. arXiv:hep-th/9606001.
\bibitem{chamseddine1996}
A.~H.~Chamseddine, A.~Connes, ``Universal Formula for Noncommutative Geometry Actions,'' Phys.\ Rev.\ Lett.\ \textbf{77} (1996) 4868.
\bibitem{ccm2007}
A.~H.~Chamseddine, A.~Connes, M.~Marcolli, ``Gravity and the standard model with neutrino mixing,'' Adv.\ Theor.\ Math.\ Phys.\ \textbf{11} (2007) 991. arXiv:hep-th/0610241.
\bibitem{cc2012}
A.~H.~Chamseddine, A.~Connes, ``Resilience of the Spectral Standard Model,'' JHEP \textbf{09} (2012) 104.

\bibitem{cc2006}
A.~H.~Chamseddine, A.~Connes, ``Scale invariance in the spectral action,'' J.\ Math.\ Phys.\ \textbf{47} (2006) 063504. arXiv:hep-th/0512169.

\bibitem{ccvs2013}
A.~H.~Chamseddine, A.~Connes, W.~D.~van Suijlekom, ``Beyond the Spectral Standard Model: Emergence of Pati--Salam Unification,'' JHEP \textbf{11} (2013) 132. arXiv:1304.8050.

\bibitem{connes2013}
A.~Connes, ``On the spectral characterization of manifolds,'' J.\ Noncommut.\ Geom.\ \textbf{7} (2013) 1. arXiv:0810.2088.
\bibitem{fuzzy1012}
M.~Ihl, C.~Sachse, C.~S\"amann, ``Fuzzy Scalar Field Theory as Matrix Quantum Mechanics,'' JHEP \textbf{03} (2011) 091. arXiv:1012.3568. DOI 10.1007/JHEP03(2011)091. (Candidate source for the induced-matrix-structure mechanism; \catmark{D}.)

% Cluster 3: string-net / topological order
\bibitem{levinwen2005}
M.~A.~Levin, X.-G.~Wen, ``String-net condensation: A physical mechanism for topological phases,'' Phys.\ Rev.\ B \textbf{71} (2005) 045110. arXiv:cond-mat/0404617.

\bibitem{kitaev2003}
A.~Yu.~Kitaev, ``Fault-tolerant quantum computation by anyons,'' Ann.\ Phys.\ \textbf{303} (2003) 2. arXiv:quant-ph/9707021.

\bibitem{mueger2003}
M.~M\"uger, ``From subfactors to categories and topology II: The quantum double of tensor categories and subfactors,'' J.\ Pure Appl.\ Algebra \textbf{180} (2003) 159. arXiv:math/0111205.

\bibitem{eno2005}
P.~Etingof, D.~Nikshych, V.~Ostrik, ``On fusion categories,'' Ann.\ Math.\ \textbf{162} (2005) 581. arXiv:math/0203060.

\bibitem{lankongwen2018}
T.~Lan, L.~Kong, X.-G.~Wen, ``Classification of (3+1)D bosonic topological orders: The case when pointlike excitations are all bosons,'' Phys.\ Rev.\ X \textbf{8} (2018) 021074. arXiv:1704.04221.

\bibitem{lanwen2019}
T.~Lan, X.-G.~Wen, ``Classification of (3+1)D bosonic topological orders (II): The case when some pointlike excitations are fermions,'' Phys.\ Rev.\ X \textbf{9} (2019) 021005. arXiv:1801.08530.
\bibitem{wen2003}
X.-G.~Wen, ``Quantum order from string-net condensations and the origin of light and massless fermions,'' Phys.\ Rev.\ D \textbf{68} (2003) 065003. arXiv:hep-th/0302201. DOI 10.1103/PhysRevD.68.065003. [Article number recertified to 065003 against APS primary metadata; the previously-circulating 024501 is a cross-contamination artifact and is incorrect.]
\bibitem{walkerwang2012}
K.~Walker, Z.~Wang, ``(3+1)-TQFTs and topological insulators,'' Front.\ Phys.\ \textbf{7} (2012) 150. arXiv:1104.2632.
\bibitem{vonkeyserlingk2013}
C.~W.~von Keyserlingk, F.~J.~Burnell, S.~H.~Simon, ``Three-dimensional topological lattice models with surface anyons,'' Phys.\ Rev.\ B \textbf{87} (2013) 045107. arXiv:1208.5128.

% Cluster 4: quantum link / D-theory
\bibitem{brower1999}
R.~Brower, S.~Chandrasekharan, U.-J.~Wiese, ``QCD as a quantum link model,'' Phys.\ Rev.\ D \textbf{60} (1999) 094502. arXiv:hep-lat/9704106.
\bibitem{brower2004}
R.~Brower, S.~Chandrasekharan, S.~Riederer, U.-J.~Wiese, ``D-theory: field quantization by dimensional reduction of discrete variables,'' Nucl.\ Phys.\ B \textbf{693} (2004) 149. arXiv:hep-lat/0309182. DOI 10.1016/j.nuclphysb.2004.06.007.
\bibitem{beard1998}
B.~B.~Beard, R.~Brower, S.~Chandrasekharan, D.~Chen, A.~Tsapalis, U.-J.~Wiese, ``D-theory: field theory via dimensional reduction of discrete variables'' (proceedings), Nucl.\ Phys.\ B Proc.\ Suppl.\ \textbf{63} (1998) 775. arXiv:hep-lat/9709120.
\bibitem{wiese2022}
U.-J.~Wiese, ``From quantum link models to D-theory: a resource efficient framework,'' Phil.\ Trans.\ R.\ Soc.\ A \textbf{380} (2022) 20210068. arXiv:2107.09335.

% Cluster 5: gauged tensor networks / PEPS
\bibitem{tagliacozzo2014}
L.~Tagliacozzo, A.~Celi, M.~Lewenstein, ``Tensor Networks for Lattice Gauge Theories with Continuous Groups,'' Phys.\ Rev.\ X \textbf{4} (2014) 041024.
\bibitem{haegeman2015}
J.~Haegeman, K.~Van Acoleyen, N.~Schuch, J.~I.~Cirac, F.~Verstraete, ``Gauging quantum states,'' Phys.\ Rev.\ X \textbf{5} (2015) 011024.
\bibitem{cirac2021}
J.~I.~Cirac, D.~P\'erez-Garc\'ia, N.~Schuch, F.~Verstraete, ``Matrix product states and projected entangled pair states,'' Rev.\ Mod.\ Phys.\ \textbf{93} (2021) 045003.
\bibitem{nielsen1981}
H.~B.~Nielsen, M.~Ninomiya, ``Absence of neutrinos on a lattice,'' Nucl.\ Phys.\ B \textbf{185} (1981) 20; \textbf{193} (1981) 173.

\bibitem{weinbergwitten1980}
S.~Weinberg, E.~Witten, ``Limits on Massless Particles,'' Phys.\ Lett.\ B \textbf{96} (1980) 59. DOI 10.1016/0370-2693(80)90212-9.

% Cluster 6: thermodynamic gravity
\bibitem{jacobson1995}
T.~Jacobson, ``Thermodynamics of Spacetime: The Einstein Equation of State,'' Phys.\ Rev.\ Lett.\ \textbf{75} (1995) 1260. arXiv:gr-qc/9504004.
\bibitem{padmanabhan2010}
T.~Padmanabhan, ``Thermodynamical Aspects of Gravity: New insights,'' Rep.\ Prog.\ Phys.\ \textbf{73} (2010) 046901. arXiv:0911.5004.
\bibitem{verlinde2011}
E.~Verlinde, ``On the Origin of Gravity and the Laws of Newton,'' JHEP \textbf{04} (2011) 029. arXiv:1001.0785.
\bibitem{chirco2014}
G.~Chirco, H.~M.~Haggard, A.~Riello, C.~Rovelli, ``Spacetime thermodynamics without hidden degrees of freedom,'' Phys.\ Rev.\ D \textbf{90} (2014) 044044. arXiv:1401.5262.

% Cluster 6b: information-theoretic / structural-realist prior art and emergent-gravity constraints
\bibitem{wheeler1990}
J.~A.~Wheeler, ``Information, Physics, Quantum: The Search for Links,'' in W.~H.~Zurek (ed.), \emph{Complexity, Entropy, and the Physics of Information}, Addison-Wesley (1990), pp.~3--28.
\bibitem{carrollsingh2018}
S.~M.~Carroll, A.~Singh, ``Mad-Dog Everettianism: Quantum Mechanics at Its Most Minimal,'' arXiv:1801.08132 (2018); in \emph{What is Fundamental?}, The Frontiers Collection, Springer (2019), DOI 10.1007/978-3-030-11301-8\_10.
\bibitem{jenkins2009}
A.~Jenkins, ``Constraints on emergent gravity,'' Int.\ J.\ Mod.\ Phys.\ D \textbf{18} (2009) 2249. arXiv:0904.0453.
\bibitem{carlip2014}
S.~Carlip, ``Challenges for Emergent Gravity,'' Stud.\ Hist.\ Phil.\ Sci.\ B \textbf{46} (2014) 200. arXiv:1207.2504.

% Cluster 7: observer / neuroscience comparison (interpretive)
\bibitem{friston2010}
K.~Friston, ``The free-energy principle: a unified brain theory?,'' Nat.\ Rev.\ Neurosci.\ \textbf{11} (2010) 127.
\bibitem{metzinger2003}
T.~Metzinger, \textit{Being No One: The Self-Model Theory of Subjectivity}, MIT Press (2003).
\bibitem{hoffman2014}
D.~Hoffman, C.~Prakash, ``Objects of Consciousness,'' Front.\ Psychol.\ \textbf{5} (2014) 577.
\bibitem{hameroff2014}
S.~Hameroff, R.~Penrose, ``Consciousness in the universe: A review of the `Orch OR' theory,'' Phys.\ Life Rev.\ \textbf{11} (2014) 39.
\bibitem{tononi2004}
G.~Tononi, ``An information integration theory of consciousness,'' BMC Neurosci.\ \textbf{5} (2004) 42.

% Cluster 8: scale-bridge / decoupling
\bibitem{attwell2001}
D.~Attwell, S.~B.~Laughlin, ``An energy budget for signaling in the grey matter of the brain,'' J.\ Cereb.\ Blood Flow Metab.\ \textbf{21} (2001) 1133.
\bibitem{harris2012}
J.~J.~Harris, R.~Jolivet, D.~Attwell, ``Synaptic energy use and supply,'' Neuron \textbf{75} (2012) 762.
\bibitem{levy2021}
W.~B.~Levy, V.~G.~Calvert, ``Communication consumes 35 times more energy than computation in the human cortex,'' PNAS \textbf{118} (2021) e2008173118.
\bibitem{appelquist1975}
T.~Appelquist, J.~Carazzone, ``Infrared singularities and massive fields,'' Phys.\ Rev.\ D \textbf{11} (1975) 2856.

% v4.x G_SM-gap research frontiers (literature-status, [D]/[E])
\bibitem{cc2008}
A.~H.~Chamseddine, A.~Connes, ``Why the Standard Model,'' J.\ Geom.\ Phys.\ \textbf{58} (2008) 38. arXiv:0706.3688. DOI 10.1016/j.geomphys.2007.09.011.
\bibitem{ccm2015}
A.~H.~Chamseddine, A.~Connes, V.~Mukhanov, ``Quanta of Geometry: Noncommutative Aspects,'' Phys.\ Rev.\ Lett.\ \textbf{114} (2015) 091302. arXiv:1409.2471.
\bibitem{grosse2005}
H.~Grosse, H.~Steinacker, ``Finite Gauge Theory on Fuzzy $CP^2$,'' Nucl.\ Phys.\ B \textbf{707} (2005) 145. arXiv:hep-th/0407089. DOI 10.1016/j.nuclphysb.2004.11.058.
\bibitem{azuma2004}
T.~Azuma, S.~Bal, K.~Nagao, J.~Nishimura, ``Nonperturbative studies of fuzzy spheres in a matrix model with the Chern--Simons term,'' JHEP \textbf{05} (2004) 005. arXiv:hep-th/0401038. DOI 10.1088/1126-6708/2004/05/005.
\bibitem{farnsworth2015}
S.~Farnsworth, L.~Boyle, ``Rethinking Connes' approach to the standard model of particle physics via non-commutative geometry,'' New J.\ Phys.\ \textbf{17} (2015) 023021. arXiv:1408.5367. DOI 10.1088/1367-2630/17/2/023021.
\bibitem{aastrup2025}
J.~Aastrup, J.~M.~Grimstrup, ``On the emergence of an almost-commutative spectral triple from a geometric construction on a configuration space,'' (2025). arXiv:2504.03391.
\bibitem{luscher1999}
M.~L\"uscher, ``Abelian chiral gauge theories on the lattice with exact gauge invariance,'' Nucl.\ Phys.\ B \textbf{549} (1999) 295. arXiv:hep-lat/9811032. DOI 10.1016/S0550-3213(99)00115-7.
\bibitem{golterman2023}
M.~Golterman, Y.~Shamir, ``Propagator zeros and lattice chiral gauge theories,'' (2023). arXiv:2311.12790; and ``Constraints on the symmetric mass generation paradigm,'' (2025). arXiv:2505.20436.
\bibitem{dupuis2021}
N.~Dupuis et al., ``The nonperturbative functional renormalization group and its applications,'' Phys.\ Rept.\ \textbf{910} (2021) 1. arXiv:2006.04853. DOI 10.1016/j.physrep.2021.01.001.
\bibitem{schnoerr2013}
D.~Schnoerr, I.~Boettcher, J.~M.~Pawlowski, C.~Wetterich, ``Error estimates and specification parameters for functional renormalization,'' (2013). arXiv:1301.4169.

% Appendix A endogenous-enhancement cluster (added in 004-appA-gsm; entries
% resolved to primary sources; flagged for the regular DOI sweep)
\bibitem{wilson1974}
K.~G.~Wilson, ``Confinement of quarks,'' Phys.\ Rev.\ D \textbf{10} (1974) 2445. DOI 10.1103/PhysRevD.10.2445.
\bibitem{frenkelkac1980}
I.~B.~Frenkel, V.~G.~Kac, ``Basic representations of affine Lie algebras and dual resonance models,'' Invent.\ Math.\ \textbf{62} (1980) 23.
\bibitem{goddardolive1986}
P.~Goddard, D.~Olive, ``Kac--Moody and Virasoro algebras in relation to quantum physics,'' Int.\ J.\ Mod.\ Phys.\ A \textbf{1} (1986) 303.
\bibitem{kuranishi1951}
M.~Kuranishi, ``On everywhere dense imbedding of free groups in Lie groups,'' Nagoya Math.\ J.\ \textbf{2} (1951) 63.
\bibitem{madore1992}
J.~Madore, ``The fuzzy sphere,'' Class.\ Quantum Grav.\ \textbf{9} (1992) 69.
\bibitem{ikkt1997}
N.~Ishibashi, H.~Kawai, Y.~Kitazawa, A.~Tsuchiya, ``A large-$N$ reduced model as superstring,'' Nucl.\ Phys.\ B \textbf{498} (1997) 467. arXiv:hep-th/9612115.
\bibitem{steinacker2010}
H.~Steinacker, ``Emergent geometry and gravity from matrix models: an introduction,'' Class.\ Quantum Grav.\ \textbf{27} (2010) 133001. arXiv:1003.4134.
\bibitem{krajewski1998}
T.~Krajewski, ``Classification of finite spectral triples,'' J.\ Geom.\ Phys.\ \textbf{28} (1998) 1. arXiv:hep-th/9701081.
\bibitem{paschkesitarz1998}
M.~Paschke, A.~Sitarz, ``Discrete spectral triples and their symmetries,'' J.\ Math.\ Phys.\ \textbf{39} (1998) 6191. arXiv:q-alg/9612029.

% Appendix B Connes-vacuum attractor cluster (added in 005-appB-connes;
% entries resolved to primary sources; flagged for the regular DOI sweep)
\bibitem{agbm1995}
E.~Alvarez, J.~M.~Gracia-Bond\'ia, C.~P.~Mart\'in, ``Anomaly cancellation and the gauge group of the standard model in non-commutative geometry,'' Phys.\ Lett.\ B \textbf{364} (1995) 33. arXiv:hep-th/9506115.
\bibitem{connes2006}
A.~Connes, ``Noncommutative geometry and the standard model with neutrino mixing,'' JHEP \textbf{11} (2006) 081. arXiv:hep-th/0608226.
\bibitem{barrett2007}
J.~W.~Barrett, ``A Lorentzian version of the non-commutative geometry of the standard model of particle physics,'' J.\ Math.\ Phys.\ \textbf{48} (2007) 012303. arXiv:hep-th/0608221.
\bibitem{barrettglaser2016}
J.~W.~Barrett, L.~Glaser, ``Monte Carlo simulations of random non-commutative geometries,'' J.\ Phys.\ A \textbf{49} (2016) 245001. arXiv:1510.01377.
\bibitem{cotler2017}
J.~S.~Cotler et al., ``Black holes and random matrices,'' JHEP \textbf{05} (2017) 118. arXiv:1611.04650.

% Cosmology inputs (Stratum-I, calibration only)
\bibitem{planck2018}
Planck Collaboration, ``Planck 2018 results. VI. Cosmological parameters,'' Astron.\ Astrophys.\ \textbf{641} (2020) A6. arXiv:1807.06209.
\bibitem{riess2022}
A.~G.~Riess et al., ``A Comprehensive Measurement of the Local Value of the Hubble Constant\ldots (SH0ES),'' Astrophys.\ J.\ Lett.\ \textbf{934} (2022) L7. arXiv:2112.04510.
\bibitem{desidr2}
DESI Collaboration, ``DESI DR2 Results II: Measurements of Baryon Acoustic Oscillations and Cosmological Constraints,'' (2025). arXiv:2503.14738.
\bibitem{kids2025}
KiDS Collaboration (A.~H.~Wright et al.), ``KiDS-Legacy: Cosmological constraints from cosmic shear with the complete Kilo-Degree Survey,'' arXiv:2503.19441 (2025).
\bibitem{des2026}
DES Collaboration (T.~M.~C.~Abbott et al.), ``Dark Energy Survey Year 6 Results: Cosmological Constraints from Galaxy Clustering and Weak Lensing,'' arXiv:2601.14559 (2026).

\end{thebibliography}

\end{document}
